% இந்தியத் தமிழ்ப் பதிப்பின் முழு திருத்தத்தக்க வாசிப்பு மூலம்.
% கட்டமைப்பு மற்றும் பக்கச் சரிபார்ப்பின் விவரங்கள் முழு மூல ZIP இன் சான்றுகளில் உள்ளன.
% கணித உள்ளடக்கமும் மொழிபெயர்ப்பும் இக்கோப்பில் நேரடியாக உள்ளன.
% எழுத்துருக்கள், upstream styles/assets/bibliography மற்றும் translation styles முழு மூல ZIP இல் தேவை.
\documentclass[11pt,a4paper,openany]{memoir}
\newcommand{\olpath}{../upstream}
\usepackage[no-math]{fontspec}
\ifLuaTeX
\setmainfont{Nirmala UI}[Script=Tamil,Renderer=HarfBuzz,WordSpace=1.3,ItalicFont={Nirmala UI},ItalicFeatures={FakeSlant=0.15}]
\setsansfont{Nirmala UI}[Script=Tamil,Renderer=HarfBuzz,WordSpace=1.3]
\else
\setmainfont{Nirmala UI}[Script=Tamil,Renderer=OpenType,WordSpace=1.3,ItalicFont={Nirmala UI},ItalicFeatures={FakeSlant=0.15}]
\setsansfont{Nirmala UI}[Script=Tamil,Renderer=OpenType,WordSpace=1.3]
\fi
\setmonofont{Consolas}
\usepackage{xr-hyper}
\input{\olpath/sty/open-logic.sty}
\ifXeTeX
  \XeTeXlinebreaklocale "ta"
  \XeTeXlinebreakskip=0pt plus 1pt
  \XeTeXgenerateactualtext=1
\fi
\setlrmarginsandblock{25mm}{25mm}{*}
\setulmarginsandblock{22mm}{25mm}{*}
\checkandfixthelayout
\linespread{1.17}
\emergencystretch=2em
\input{../translation/tamil-config.sty}
\input{../translation/tamil-pdf-math.sty}
% Latin diacritics retain their spelling; Tamil remains in Nirmala UI.
% fontspec must be loaded with no-math so legacy mathematical accent fonts
% are not replaced by a text font lacking the mathematical bar glyph.
\newfontfamily\TALatinFont{Cambria}
\RenewDocumentCommand\L{}{{\TALatinFont\char"0141}}
\RenewDocumentCommand\l{}{{\TALatinFont\char"0142}}
\NewCommandCopy\TAOriginalAcute\'
\RenewDocumentCommand\'{m}{{\TALatinFont\TAOriginalAcute{#1}}}
\usepackage{newunicodechar}
\newunicodechar{Ł}{{\TALatinFont\char"0141}}
\newunicodechar{ł}{{\TALatinFont\char"0142}}
\newunicodechar{ś}{{\TALatinFont\char"015B}}
\AtBeginEnvironment{thebibliography}{\TALatinFont}
% cleveref's list joins and enumerate labels are generated prose; localize them
% after babel has selected the document language.
\AtBeginDocument{%
  \crefname{enumi}{படி}{படிகள்}%
  \Crefname{enumi}{படி}{படிகள்}%
  \crefname{enumii}{படி}{படிகள்}%
  \Crefname{enumii}{படி}{படிகள்}%
  \renewcommand{\crefpairconjunction}{ மற்றும் }%
  \renewcommand{\crefmiddleconjunction}{, }%
  \renewcommand{\creflastconjunction}{ மற்றும் }%
  \renewcommand{\crefpairgroupconjunction}{ மற்றும் }%
  \renewcommand{\crefmiddlegroupconjunction}{, }%
  \renewcommand{\creflastgroupconjunction}{, மற்றும் }%
  \renewcommand{\crefrangeconjunction}{--}%
}
% The upstream starred induction case and exercise marker contain hard-coded
% English even when their source arguments are translated.
\RenewDocumentCommand\indcase { s t{!} m m +m }{%
  \DeclareDocumentMacro \indfrm {#3}%
  \DeclareDocumentMacro \indfrmp {#3}%
  \DeclareDocumentMacro \indcomplex {#4}%
  \IfBooleanTF{#1}{$#3$ அணு வாய்பாடு: }{$#3 \ident #4$: }%
  \IfBooleanTF{#2}{பயிற்சி.}{#5}}

\tagtrue{novice,math,compsci,FOL,prvTrue,prvFalse,prvEx,prvAll}
\addto\captionsenglish{%
  \renewcommand{\chaptername}{அத்தியாயம்}%
  \renewcommand{\partname}{பகுதி}%
  \renewcommand{\appendixname}{இணைப்பு}%
  \renewcommand{\tablename}{அட்டவணை}%
}
\setlocalecaption{english}{photocredits}{படங்களுக்கான நன்றியுரை}
\setlocalecaption{english}{history}{வரலாற்றுக் குறிப்புகள்}
\setlocalecaption{english}{reading}{மேலும் படிக்க}
\setlocalecaption{english}{appendix}{இணைப்பு}
\setlocalecaption{english}{appendices}{இணைப்புகள்}
\setlocalecaption{english}{definition}{வரையறை}
\setlocalecaption{english}{example}{எடுத்துக்காட்டு}
\setlocalecaption{english}{lemma}{துணைத்தேற்றம்}
\setlocalecaption{english}{proposition}{கூற்று}
\setlocalecaption{english}{corollary}{தொடர்விளைவு}
\setlocalecaption{english}{problem}{பயிற்சி}
\setlocalecaption{english}{problems}{பயிற்சிகள்}
\setlocalecaption{english}{theorem}{தேற்றம்}
\setlocalecaption{english}{remark}{குறிப்பு}
\setlocalecaption{english}{axiom}{அடிகோள்}
\setlocalecaption{english}{note}{குறிப்பு}
\setlocalecaption{english}{case}{நிலை}
\setlocalecaption{english}{convention}{மரபு}
\addto\captionsenglish{\renewcommand{\proofname}{நிறுவல்}}
\addto\captionsenglish{\renewcommand{\figurename}{படம்}}
\crefname{page}{பக்கம்}{பக்கங்கள்}
\Crefname{page}{பக்கம்}{பக்கங்கள்}
\crefname{part}{பகுதி}{பகுதிகள்}
\Crefname{part}{பகுதி}{பகுதிகள்}
\crefname{chapter}{அத்தியாயம்}{அத்தியாயங்கள்}
\Crefname{chapter}{அத்தியாயம்}{அத்தியாயங்கள்}
\crefname{section}{பிரிவு}{பிரிவுகள்}
\Crefname{section}{பிரிவு}{பிரிவுகள்}
\crefname{subsection}{பிரிவு}{பிரிவுகள்}
\Crefname{subsection}{பிரிவு}{பிரிவுகள்}
\crefname{figure}{படம்}{படங்கள்}
\Crefname{figure}{படம்}{படங்கள்}
\crefname{table}{அட்டவணை}{அட்டவணைகள்}
\Crefname{table}{அட்டவணை}{அட்டவணைகள்}
\crefname{equation}{சமன்பாடு}{சமன்பாடுகள்}
\Crefname{equation}{சமன்பாடு}{சமன்பாடுகள்}
\crefname{defn}{வரையறை}{வரையறைகள்}
\Crefname{defn}{வரையறை}{வரையறைகள்}
\crefname{thm}{தேற்றம்}{தேற்றங்கள்}
\Crefname{thm}{தேற்றம்}{தேற்றங்கள்}
\crefname{ex}{எடுத்துக்காட்டு}{எடுத்துக்காட்டுகள்}
\Crefname{ex}{எடுத்துக்காட்டு}{எடுத்துக்காட்டுகள்}
\crefname{lem}{துணைத்தேற்றம்}{துணைத்தேற்றங்கள்}
\Crefname{lem}{துணைத்தேற்றம்}{துணைத்தேற்றங்கள்}
\crefname{prop}{கூற்று}{கூற்றுகள்}
\Crefname{prop}{கூற்று}{கூற்றுகள்}
\crefname{prob}{பயிற்சி}{பயிற்சிகள்}
\Crefname{prob}{பயிற்சி}{பயிற்சிகள்}
\crefname{rem}{குறிப்பு}{குறிப்புகள்}
\Crefname{rem}{குறிப்பு}{குறிப்புகள்}
\crefname{cor}{தொடர்விளைவு}{தொடர்விளைவுகள்}
\Crefname{cor}{தொடர்விளைவு}{தொடர்விளைவுகள்}
\crefname{axiom}{அடிகோள்}{அடிகோள்கள்}
\Crefname{axiom}{அடிகோள்}{அடிகோள்கள்}
\crefname{note}{குறிப்பு}{குறிப்புகள்}
\Crefname{note}{குறிப்பு}{குறிப்புகள்}
\crefname{case}{நிலை}{நிலைகள்}
\Crefname{case}{நிலை}{நிலைகள்}
\crefname{conv}{மரபு}{மரபுகள்}
\Crefname{conv}{மரபு}{மரபுகள்}
\AtBeginEnvironment{align*}{\fontsize{9.5}{12}\selectfont}
\hypersetup{pdftitle={திறந்த தருக்கவியல் — மூல மாற்று வடிவங்கள்},pdfauthor={Open Logic Project; தமிழ் மொழிபெயர்ப்புத் திட்டம்},pdfsubject={உறைந்த OpenLogic மூலத்தின் மாற்று வடிவங்கள்; இயந்திரத் தமிழ் மொழிபெயர்ப்பு}}
\addto\captionsenglish{\renewcommand{\contentsname}{பொருளடக்கம்}\renewcommand{\bibname}{மேற்கோள் நூல்}}
\tagtrue{prfSC,prfND,prfTab,prfAX}
\includeenv{editorial}
\externaldocument{tamil-complete}[tamil-complete.pdf]
\makeatletter
\newcommand\TADriverHeadingMode{\def\ol@chaptercs{section}\def\ol@partcs{section}}
\makeatother
% Alternate drivers retain their headings, editorial prose and import
% inventories. Their shared chapters are read in the principal volume.
\NewDocumentCommand\TADriverImport{s o m o}{%
  \par\noindent மூல உள்ளீடு:
  \IfBooleanT{#1}{\texttt{*}}%
  \IfNoValueF{#2}{\texttt{\detokenize{#2}/}}%
  \texttt{\detokenize{#3}.tex}%
  \IfNoValueF{#4}{\quad பிரிவு வடிவம்: \texttt{\detokenize{#4}}}%
  \par
}
\NewDocumentCommand\TACompanionUnit{m m}{%
  \clearpage
  \noindent மூல அலகு: \texttt{#1}\par
  \subfile{../translation/#2}%
}
\NewDocumentCommand\TACompanionDriver{m m m m}{%
  \clearpage
  \noindent மூல அலகு: \texttt{#1}\par
  \begingroup
  \TADriverHeadingMode
  \olchapterid{#3}{#4}%
  \RenewDocumentCommand\olimport{s o m o}{%
    \IfBooleanTF{##1}{%
      \IfNoValueTF{##2}{%
        \IfNoValueTF{##4}{\TADriverImport*{##3}}{\TADriverImport*{##3}[##4]}%
      }{%
        \IfNoValueTF{##4}{\TADriverImport*[##2]{##3}}{\TADriverImport*[##2]{##3}[##4]}%
      }%
    }{%
      \IfNoValueTF{##2}{%
        \IfNoValueTF{##4}{\TADriverImport{##3}}{\TADriverImport{##3}[##4]}%
      }{%
        \IfNoValueTF{##4}{\TADriverImport[##2]{##3}}{\TADriverImport[##2]{##3}[##4]}%
      }%
    }%
  }%
  \subfile{../translation/#2}%
  \endgroup
}
\begin{document}
\raggedright
\pagestyle{plain}
\chapter*{மூல மாற்று வடிவங்களும் கூடுதல் பகுதிகளும்}
\noindent இந்தியத் தமிழ்ப் பதிப்பு — ta-Taml-IN

\noindent மூல நூல்: Open Logic Project.
\href{https://github.com/OpenLogicProject/OpenLogic}{மூலத் திட்டம்},
\href{https://creativecommons.org/licenses/by/4.0/}{CC BY 4.0}.
மூலத் திருத்தம்: \texttt{9620cc73f9c8e0ad003c514a5d3748f29611c4c0}.

\noindent முதல் 570 அலகுகளின் இயந்திரத் தமிழாக்கமும் திருத்தங்களும்
OpenAI Codex — GPT-5.6 Sol, Ultra சிந்தனை நிலையில் செய்யப்பட்டன.
மீதமுள்ள 152 அலகுகளின் தமிழாக்கமும் பின்னைய பதிப்புத் திருத்தங்களும்
OpenAI Codex — GPT-6 Sol, Ultra சிந்தனை நிலையில் செய்யப்பட்டன.
சுயாதீனமான தாய்மொழி மதிப்புரை செய்யப்பட்டதாகக் கூறப்படவில்லை.
இந்தத் துணைநூலும் முதன்மை நூலும்
சேர்ந்து உறைந்த மூலத்தின் 722 அலகுகளையும் கணக்கிடுகின்றன.

\noindent முதன்மை வாசிப்பு நூலின் 695 மூல அலகுகளுடன்,
இங்கு 27 தனித்த மூல அலகுகள் சேர்க்கப்படுகின்றன.
அவற்றில் நான்கு மாற்று அத்தியாயம் அல்லது பகுதி அமைப்புக் கோப்புகள்.
அவற்றின் தமிழ்த் தலைப்புகள், ஆசிரியர் குறிப்புகள் மற்றும் மூல
உள்ளீட்டுப் பட்டியல்கள் இங்கே உள்ளன; ஏற்கெனவே முதன்மை நூலில்
இடம்பெறும் அத்தியாயங்கள் மீண்டும் அச்சிடப்படவில்லை.
மற்ற 23 அலகுகளின் முழுப் பாடப்பொருள் இங்கே இடம்பெறுகிறது.
அனைத்து 722 அலகுகளின் திருத்தத்தக்க மூலக் கோப்புகளும்
முழு மூலத் தொகுப்பில் உள்ளன.

\noindent \href{tamil-complete.pdf}{முதன்மை வாசிப்பு நூலைத் திறக்க}.
மூலத்தின் மாற்று வடிவங்கள் சில இடங்களில் அதே அடையாளங்களைப்
பயன்படுத்துகின்றன. இந்நூலில் உள்ள உள்ளூர் முடிவுகளின் இணைப்புகள்
இந்நூலுக்குள்ளும், பிற மேற்கோள்கள் முதன்மை நூலுக்கும் செல்கின்றன.
உறைந்த மூலத்தில் இல்லாத நிறுவல் அல்லது விளக்கம் சேர்க்கப்பட்டதாகக்
கூறப்படவில்லை.

\clearpage
\begingroup
\sloppy
\renewcommand{\cftsectionfont}{\small}
\tableofcontents*
\endgroup
\chapter{முதல்தரத் தருக்கத்தின் மாற்று வடிவங்கள்}
\clearpage
\noindent மூல அலகு: \texttt{OLP-0643}\par

% BEGIN SOURCE UNIT OLP-0643 content/first-order-logic/axiomatic-deduction/provability.tex
% Exact modular target bytes: 20160; SHA256: 3856cfa15db843b807b92faff65d6230a05253c2e4dc4aecd4e726b6b47a5eff
% SEGMENT OLP-0643-S01


\iftag{FOL}
      {\olfileid{fol}{axd}{prv}}
      {\olfileid{pl}{axd}{prv}}

\olsection{\usetoken{S}{derivability} இன் பண்புகள்}

\begin{prop}[ஒருபோக்குத்தன்மை] $\Gamma \subseteq \Delta$ மற்றும்
$\Gamma \Proves !A$ எனில், $\Delta \Proves !A$ உம் ஆகும்.
\end{prop}

\begin{proof} $\Gamma$ இன் எந்த முடிவுறு உட்கணம்~$\Gamma_0$ உம்
$\Delta$ இன் முடிவுறு உட்கணமாகும். ஆகவே $\Gamma_0 \Proves !A$ என்பதற்கான
!!a{derivation}, $\Delta \Proves !A$ என்பதையும் காட்டுகிறது.
\end{proof}

% SEGMENT OLP-0643-S02
\begin{prop}\ollabel{prop:provability} $\Gamma \Proves !A$ என்பது,
$\Gamma\cup \{\lnot!A \}$ முரணுடையதாக இருந்தால், அப்போதும் அப்போது
மட்டுமே ஆகும். \end{prop}

\begin{prop} \begin{enumerate}
\item \ollabel{prop:provability-contr} $\Gamma \Proves !A$ மற்றும்
$\Gamma \cup \{ !A\} \Proves \lfalse$ எனில், $\Gamma$ முரணுடையது.

\item \ollabel{prop:provability-lnot} $\Gamma \cup \{!A\}
\Proves \lfalse$ எனில், $\Gamma \Proves \lnot !A$.

\item \ollabel{prop:provability-exhaustive} $\Gamma \cup \{!A\}
\Proves \lfalse$ மற்றும் $\Gamma \cup \{\lnot !A\} \Proves
\lfalse$ எனில், $\Gamma \Proves \lfalse$.

\tagitem{prvOr}{\ollabel{prop:provability-lor-left} $\Gamma \cup \{!A\}
\Proves \lfalse$ மற்றும் $\Gamma \cup \{!B\} \Proves
\lfalse$ எனில், $\Gamma \cup \{!A \lor !B\} \Proves \lfalse$.}{}

\tagitem{prvOr}{\ollabel{prop:provability-lor-right} $\Gamma
\Proves !A$ அல்லது $\Gamma \Proves !B$ எனில், $\Gamma
\Proves !A \lor !B$.}{}

\tagitem{prvAnd}{\ollabel{prop:provability-land-left} $\Gamma
\Proves !A \land !B$ எனில், $\Gamma \Proves !A$ மற்றும்
$\Gamma \Proves !B$ ஆகிய இரண்டும் கிடைக்கும்.}{}

\tagitem{prvAnd}{\ollabel{prop:provability-land-right} $\Gamma
\Proves !A$ மற்றும் $\Gamma \Proves !B$ எனில், $\Gamma
\Proves !A \land !B$.}{}

\tagitem{prvIf}{\ollabel{prop:provability-mp} $\Gamma \Proves
!A$ மற்றும் $\Gamma \Proves !A \lif !B$ எனில், $\Gamma
\Proves !B$.}{}

\tagitem{prvIf}{\ollabel{prop:provability-lif} $\Gamma \Proves
\lnot !A$ அல்லது $\Gamma \Proves !B$ எனில், $\Gamma \Proves
!A \lif !B$.}{} \end{enumerate} \end{prop}

% SEGMENT OLP-0643-S03
\begin{proof}
\begin{enumerate}
\item $\Gamma_0 \cup \{!A\} \Proves \lfalse$ எனக் கொள்வோம்.

\begin{derivation}
1. & $\Gamma_0 \cup \{!A\} \Proves \lfalse$ & கருதுகோள் \\
2. & $\Gamma_1 \Proves !A$ & கருதுகோள் \\
3. & $\Gamma_0 \cup \Gamma_1 \Proves \lfalse$ & வெட்டு \\
\end{derivation}

$\Gamma_0 \subseteq \Gamma$ மற்றும் $\Gamma_1 \subseteq \Gamma$ என்பதால்,
$\Gamma_0 \cup \Gamma_1 \subseteq \Gamma$. எனவே $\Gamma \Proves \lfalse$.

% SOURCE CORRECTION TA-FOL-PRV-001: the frozen derivation proves the
% opposite-direction proposition, rather than this negation conclusion.
\item $\Gamma_0 \cup\{!A\} \Proves \lfalse$ எனக் கொள்வோம்.

\begin{derivation}
1. & $\Gamma_0 \cup\{!A\} \Proves \lfalse$ & கருதுகோள் \\
2. & $\Gamma_0 \Proves !A \lif \lfalse$ & கழித்தல் தேற்றம் \\
3. & $\Gamma_0 \Proves (!A \lif \lfalse) \lif \lnot !A$ & Ax0 \\
4. & $\Gamma_0 \Proves \lnot !A$ & MP 2, 3 \\
\end{derivation}

\item $\Gamma \cup \{ !A \} \Proves \lfalse$ மற்றும்
$\Gamma_1 \cup \{\lnot !A \} \Proves \lfalse$ எனக் கொள்வோம்.

\begin{derivation}
1. & $\Gamma_0 \cup\{!A\} \Proves \lfalse$ & கருதுகோள் \\
2. & $\Gamma_0 \Proves !A \lif \lfalse$ & கழித்தல் தேற்றம் \\
3. & $\Gamma_1 \cup \{\lnot !A \} \Proves \lfalse$ & கருதுகோள் \\
4. & $\Gamma_0 \Proves (!A \lif \lfalse) \lif \lnot !A$ & Ax0 \\
5. & $\Gamma_0 \Proves \lnot !A$ & MP 2, 4 \\
6. & $\Gamma_0 \cup \Gamma_1 \Proves \lfalse$ & வெட்டு 3, 5 \\
\end{derivation}

$\Gamma_0 \subseteq \Gamma$ மற்றும் $\Gamma_1 \subseteq \Gamma$ என்பதால்,
$\Gamma_0 \cup \Gamma_1 \subseteq \Gamma$. எனவே $\Gamma \Proves \lfalse$.

% SEGMENT OLP-0643-S04
% prop:provability-lor-left
\tagitem{defOr}{}{\iftag{probOr}{பயிற்சி.}{பயிற்சி.}}

% prop:provability-lor-right
\tagitem{defOr}{}{\iftag{probOr}{பயிற்சி.}{
$\Gamma_0 \Proves !A$ எனக் கொள்வோம்.

\begin{derivation}
1. & $\Gamma_0 \Proves !A$ & கருதுகோள் \\
2. & $\Gamma_0 \Proves !A \lif (!A \lor !B)$ & Ax0 \\
3. & $\Gamma_0 \Proves !A \lor !B$ & MP 1, 2 \\
\end{derivation}

எனவே $\Gamma \Proves !A \lor !B$. $\Gamma \Proves !B$ என்ற
நிலைக்கான நிறுவலும் இதேபோன்றது.}}

% prop:provability-land-left
\tagitem{defAnd}{}{\iftag{probAnd}{பயிற்சி.}{
$\Gamma_0 \Proves !A \land !B$ எனக் கொள்வோம்.

\begin{derivation}
1. & $\Gamma_0 \Proves !A \land !B$ & கருதுகோள் \\
2. & $\Gamma_0 \Proves (!A\land !B) \lif !A $ & Ax0 \\
% SOURCE CORRECTION TA-FOL-PRV-002: MP yields A, not A-or-B.
3. & $\Gamma_0 \Proves !A$ & MP 1, 2 \\
\end{derivation}

ஆகவே $\Gamma \Proves !A$. இதேபோன்ற மற்றொரு !!{derivation} மூலம்
$\Gamma \Proves !B$ என்பதையும் காட்டலாம்.}}

% prop:provability-land-right
\tagitem{defAnd}{}{\iftag{probAnd}{பயிற்சி.}{பயிற்சி.}}

% prop:provability-mp
\tagitem{defIf}{}{\iftag{probIf}{பயிற்சி.}{
இரு வருவித்தல்களிலும் பயன்படும் முடிவுறு கருதுகோள்களை ஒன்றிணைத்து,
அவற்றுக்கு ஒரு பொதுவான முடிவுறு துணைக்கணத்தை எடுத்துக்கொள்ளலாம்.
$\Gamma_0 \Proves !A$ மற்றும் $\Gamma_0 \Proves !A \lif !B $
ஆகிய இரண்டையும் கருதுவோம்.

\begin{derivation}
1. & $\Gamma_0 \Proves !A$ & கருதுகோள் \\
2. & $\Gamma_0 \Proves !A \lif !B $ & கருதுகோள் \\
3. & $\Gamma_0 \Proves !B$ & MP 1, 2 \\
\end{derivation}

% SOURCE CORRECTION TA-FOL-PRV-003: Gamma_1 is absent from this proof.
$\Gamma_0 \subseteq \Gamma$ என்பதால், $\Gamma \Proves !B$.}}

% prop:provability-lif
\tagitem{defIf}{}{\iftag{probIf}{பயிற்சி.}{
முதலில் $\Gamma \Proves \lnot !A$ எனக் கொள்வோம்.

\begin{derivation}
1. & $\Gamma_0 \Proves \lnot !A$ & கருதுகோள் \\
2. & $\Gamma_0 \Proves \lnot !A \lif (!A \lif !B) $ & Ax0 \\
3. & $\Gamma_0 \Proves !A \lif !B$ & MP 1, 2 \\
\end{derivation}

$\Gamma \Proves !B$ என்பதற்கான நிறுவலும் இதேபோன்றது:

\begin{derivation}
1. & $\Gamma_0 \Proves !B$ & கருதுகோள் \\
2. & $\Gamma_0 \Proves !B \lif (!A \lif !B) $ & Ax0 \\
3. & $\Gamma_0 \Proves !A \lif !B$ & MP 1, 2 \\
\end{derivation}}}

\end{enumerate}
\end{proof}

% SEGMENT OLP-0643-S05
\begin{probtag}{probOr,probAnd,probIf} \olref[fol][axd][prv]{prop:provability}
இன் நிறுவலை நிறைவு செய்க.
\end{probtag}

% SEGMENT OLP-0643-S06
\begin{thm}[பொதுமைப்படுத்தல்] \ollabel{thm:Generalization}
$\Gamma \Proves !A$ என்றும், $\Gamma$ இல் உள்ள எந்த !!{formula} விலும்
$x$ கட்டற்றதாக இல்லை என்றும் வைத்துக்கொண்டால், $\Gamma
\Proves \lforall[x][!A]$.
\end{thm}

\begin{proof} $\mathsf{Thm}_\Gamma$ மீது தொகுத்தறிதல் செய்வோம். $!A$
ஓர் அடிகோள் எனில், $\lforall[x][!A]$ உம் ஓர் அடிகோள். $!A\in \Gamma$
எனில், $!A$ இல் $x$ கட்டற்றதாக இல்லை. ஆகவே $!A \lif
\lforall[x][!A]$ ஓர் அடிகோள் (\textbf{Ax3}); மோடஸ் போனென்ஸால்
$\Gamma \Proves \lforall[x][!A]$. இப்போது $\Gamma \Proves !B$ மற்றும்
$\Gamma \Proves !B \lif !A$ என்பவற்றிலிருந்து மோடஸ் போனென்ஸால் $!A$
கிடைக்கிறது என்க. தொகுத்தறிதல் கருதுகோளின்படி $\Gamma \Proves
\lforall[x][!B]$ மற்றும் $\Gamma \Proves \lforall[x][( !B \lif
!A)]$. மேலும் \textbf{Ax2} இன்படி
\[ \Gamma \Proves
\lforall[x][( !B \lif !A)] \lif (\lforall[x][ !B] \lif \lforall[x][!A]), \]
எனவே மோடஸ் போனென்ஸை இருமுறை பயன்படுத்தினால், தேவையானபடி
$\Gamma \Proves \lforall[x][!A]$ கிடைக்கும்.
\end{proof}

% SEGMENT OLP-0643-S07
\begin{thm}\ollabel{thm:WeakGen} (\emph{மாறிலிக் குறிகள் குறித்த பலவீனமான
பொதுமைப்படுத்தல்})
$\Gamma \Proves !A$ என்றும், $c$ என்பது $\Gamma$ இல் தோன்றாத
!!a{constant} என்றும் வைத்துக்கொள்வோம். அப்போது $!A$ இல் தோன்றாத
ஒரு மாறி~$x$ உண்டு; அதற்கு $\Gamma \Proves
\lforall[x][\Subst{!A}{x}{c}]$. மேலும் இந்த நிறுவலில் $c$ இடம்பெறாது.
\end{thm}

\begin{proof} $!A_1,\ldots,!A_n$ என்பது $\Gamma$ இலிருந்து $!A$ க்கான
நிறுவல் என்றும் $!A=!A_n$ என்றும் கொள்வோம். $!A_1,\ldots,!A_n$ இல்
தோன்றாத மாறி~$x$ ஐத் தேர்ந்தெடுத்து, புதிய வரிசையைக் கருதுவோம்:
$\Subst{!A_1}{x}{c},\ldots,\Subst{!A_n}{x}{c}.$
இந்த வரிசை $\Gamma$ இலிருந்து $\Subst{!A}{x}{c}$ க்கான நிறுவல்.
உண்மையில், ஒவ்வொரு $i=1,\ldots,n$ க்கும்:
\begin{itemize}
% SOURCE CORRECTION TA-FOL-PRV-004: the axiom case concerns A_i.
\item $!A_i$ ஓர் அடிகோள் எனில், $\Subst{!A_i}{x}{c}$ உம் அடிகோளே;
\item $!A_i \in \Gamma$ எனில், $c$ ஆனது $\Gamma$ இல் இல்லாததால்,
$\Subst{!A_i}{x}{c} = !A_i \in \Gamma$;
\item $!A_j$ மற்றும் $!A_j \lif !A_i$ ஆகியவற்றிலிருந்து $!A_i$
பெறப்பட்டிருந்தால், $\Subst{!A_j}{x}{c}$ மற்றும்
$\Subst{(!A_j \lif !A_i)}{x}{c}
= \Subst{!A_j}{x}{c} \lif \Subst{!A_i}{x}{c}$ ஆகியவற்றிலிருந்து
மோடஸ் போனென்ஸால் $\Subst{!A_i}{x}{c}$ கிடைக்கும்.
\end{itemize}
புதிய வரிசையில் !!{constant}~$c$ இனி தோன்றாது. இப்போது
$\Subst{!A}{x}{c}$ இன் !!{derivation} இல் இடம்பெறும் $\Gamma$ வின்
!!{formula}s அனைத்தையும் கொண்டது $\Gamma'$ என்க. அப்போது $\Gamma'$ இன்
எந்த !!{formula} விலும் $x$ கட்டற்றதாக இல்லை; மேலும் $\Gamma'
\Proves \Subst{!A}{x}{c}$. பொதுமைப்படுத்தலின்படி
(\olref{thm:Generalization}), $\Gamma'
\Proves \lforall[x][\Subst{!A}{x}{c}]$; ஒருபோக்குத்தன்மையால்
தேவையானபடி $\Gamma \Proves \lforall[x][\Subst{!A}{x}{c}]$.
\end{proof}

% SEGMENT OLP-0643-S08
\begin{explain}
\olref{thm:WeakGen} இல் $x$ ஆனது $!A$ இல் எங்குமே தோன்றக் கூடாது
என்ற நிபந்தனை உள்ளது. அதற்குப் பதிலாக, $x$ ஆனது $!A$ இல் கட்டற்றதாக
இருக்கக் கூடாது என்றும், $!A$ இல் $c$ க்காக !!{free for} ஆக இருக்க
வேண்டும் என்றும் கூறும் மெலிந்த நிபந்தனைகளைப் பெறுவதே நமது நோக்கம்.
கட்டுண்ட !!{variable} யை மாற்றுவதன் மூலம் இதைச் செய்யலாம்; அதை முதலில்
நிறுவுவோம்.
\end{explain}

% SEGMENT OLP-0643-S09
\begin{lem}
\ollabel{lem:free-for}
$!A$ இல் $c$ க்காக $x$ இடைநிறுத்தத்தக்கதாகவும், $!A$ இல் $y$
கட்டற்றதாக இல்லாமலும் இருந்தால், $\Subst{!A}{y}{c}$ இல் $y$ க்காக $x$
!!{free for} ஆகும்.
\end{lem}

\begin{proof}
% SOURCE CORRECTION TA-FOL-PRV-005: test the same substituted formula
% as in the lemma statement and the following sentence.
$\Subst{!A}{y}{c}$ இல் $y$ க்காக $x$ !!{free for} ஆகவில்லை என்க.
அப்போது $\Subst{!A}{y}{c}$ இல் கட்டற்றதாகத் தோன்றும் $y$ வின் ஓர்
இடம், $\lforall[x]$ என்ற அளவையடையின் ஆளுகைக்குள் இருக்கும். ஆனால்
கருதுகோளின்படி $!A$ இல் $y$ கட்டற்றதாக இல்லை. எனவே அவ்விடங்கள் அனைத்தும்
$\Subst{}{y}{c}$ என்ற இடைநிறுத்தலிலிருந்து வந்தவை. ஆகவே $c$ இன் ஓர்
இடம் $\lforall[x]$ இன் ஆளுகைக்குள் உள்ளது; அதனால் $!A$ இல் $c$
க்காக $x$ !!{free for} ஆகாது.
\end{proof}

% SEGMENT OLP-0643-S10
\begin{lem}[கட்டுண்ட மாறியை மாற்றுதல்]
\ollabel{lem:ChangeBdVar}
$!A$ இல் $x$, $y$ இரண்டும் கட்டற்றதாக இல்லாமல், இரண்டும் $!A$ இல் $c$
க்காக !!{free for} ஆக இருந்தால்,
$\Proves \lforall[x][\Subst{!A}{x}{c}] \equiv
\lforall[y][\Subst{!A}{y}{c}]$; இந்த நிறுவலில் $c$ இடம்பெறாது.
\end{lem}

\begin{proof}
கருதுகோள்களின்படி $!A$ இல் $c$ க்காக $y$ !!{free for} ஆகும்;
$!A$ இல் $x$ கட்டற்றதாக இல்லை. எனவே முந்தைய
\olref{lem:free-for} படி, $\Subst{!A}{x}{c}$ இல் $x$ க்காக $y$
!!{free for} ஆகும். ஆகவே
$\lforall[x][\Subst{!A}{x}{c} \lif \Subst{!A}{x}{c} \Subst{}{y}{x}]$
ஓர் அடிகோள் (\textbf{Ax1}). $!A$ இல் $x$ கட்டற்றதாக இல்லாததால்,
$\Subst{!A}{x}{c} \Subst{}{y}{x} = \Subst{!A}{y}{c}$. எனவே
\[
\Proves \lforall[x][\Subst{!A}{x}{c}] \lif \Subst{!A}{y}{c},
\]
கழித்தல் தேற்றத்தால் $\lforall[x][\Subst{!A}{x}{c}] \Proves
\Subst{!A}{y}{c}$. மேலும் $!A$ இல் $y$ கட்டற்றதாக இல்லாததால்,
$\lforall[x][\Subst{!A}{x}{c}]$ இலும் அது கட்டற்றதாக இல்லை. எனவே
பொதுமைப்படுத்தலால்
$\lforall[x][\Subst{!A}{x}{c}] \Proves
\lforall[y][\Subst{!A}{y}{c}]$; கழித்தல் தேற்றத்தை மீண்டும்
பயன்படுத்தினால் $\Proves
\lforall[x][\Subst{!A}{x}{c}] \lif \lforall[y][\Subst{!A}{y}{c}]$.
$\Proves \lforall[y][\Subst{!A}{y}{c}] \lif \lforall[x]
[\Subst{!A}{x}{c}]$ என்பதற்கான நிறுவல் இதன் சமச்சீரான வடிவமே.
ஆகவே தர்க்கமெய்மையால் முடிவு கிடைக்கிறது.
\end{proof}

% SEGMENT OLP-0643-S11
\begin{thm}[மாறிலிக் குறிகளின் வலிமையான பொதுமைப்படுத்தல்]
\ollabel{thm:StrongGen}
$\Gamma \Proves !A$ என்றும், !!{constant}~$c$ ஆனது $\Gamma$ இல்
தோன்றாது என்றும், $x$ ஆனது $!A$ இல் கட்டற்றதாக இல்லாமல் $!A$ இல் $c$
க்காக !!{free for} ஆகும் என்றும் வைத்துக்கொள்வோம். அப்போது
$\Gamma \Proves \lforall[x][\Subst{!A}{x}{c}]$.
\end{thm}

\begin{proof}
$\Gamma \Proves !A$ என்றும் $c$ ஆனது $\Gamma$ இல் இல்லை என்றும்
கொண்டால், பலவீனமான பொதுமைப்படுத்தலால் $!A$ இல் தோன்றாத
!!a{variable}~$y$ ஒன்று உண்டு; அதற்கு $\Gamma \Proves
\lforall[y][\Subst{!A}{y}{c}]$. $!A$ இல் $y$ எங்குமே தோன்றாததால்,
அது $!A$ இல் கட்டற்றதாகவும் இல்லை; $!A$ இல் $c$ க்காக
!!{free for} ஆகவும் இருக்கிறது. மேலும், கருதுகோளின்படி $x$ உம்
$!A$ இல் கட்டற்றதாக இல்லாமல் $!A$ இல் $c$ க்காக !!{free for} ஆகிறது.
எனவே கட்டுண்ட !!{variable} யை மாற்றுவதற்கான நிபந்தனைகள் நிறைவேறுகின்றன:
$\Proves \lforall[x][\Subst{!A}{x}{c}] \equiv
\lforall[y][\Subst{!A}{y}{c}]$. இதிலிருந்து $\Gamma \Proves
\lforall[x][\Subst{!A}{x}{c}]$ கிடைக்கும்.
\end{proof}

% SEGMENT OLP-0643-S12
\begin{thm}
\ollabel{thm:provability-quantifiers}
\begin{tagenumerate}{prvEx,prvAll}
\tagitem{prvEx}{$\Gamma \Proves !A(t)$ எனில், $\Gamma \Proves
  \lexists[x][!A(x)]$.}{}
\tagitem{prvAll}{$\Gamma \Proves \lforall[x][!A(x)]$ எனில், $\Gamma
  \Proves !A(t)$.}{}
\end{tagenumerate}
\end{thm}

\begin{proof}
\begin{tagenumerate}{prvEx,prvAll}
%% இருப்பளவைக்கான அடிகோள்கள் தேவை
\tagitem{prvEx}{பயிற்சி.}{}
\tagitem{prvAll}{ $\Gamma_0 \Proves \lforall[x][!A(x)]$ எனக் கொள்வோம்.

\begin{derivation}
1. & $\Gamma_0 \Proves \lforall[x][!A(x)]$ & கருதுகோள் \\
2. & $\Gamma_0 \Proves \lforall[x][!A(x)] \lif \Subst{!A}{t}{x}$ & Ax1 \\
3. & $\Gamma_0 \Proves \Subst{!A}{t}{x}$ & MP 1, 2 \\
\end{derivation}}{}

\end{tagenumerate}
\end{proof}


% END SOURCE UNIT OLP-0643

\clearpage
\noindent மூல அலகு: \texttt{OLP-0644}\par

% BEGIN SOURCE UNIT OLP-0644 content/first-order-logic/completeness/maximally-consistent-sets.tex
% Exact modular target bytes: 12622; SHA256: e1042903b8d8d8a517a7e31e0903af09a8d0bc28810fdae6b6902cac093a7437
% SEGMENT OLP-0644-S01


\olfileid{fol}{com}{mcs}
\olsection{அதிகபட்ச முரணற்ற \usetoken{P}{sentence}-கணங்கள்}

\begin{defn}[அதிகபட்ச முரணற்ற கணம்]
\ollabel{mcs}
!!{sentence}s கொண்ட கணம்~$\Gamma$, பின்வரும் இரண்டு நிபந்தனைகளும்
நிறைவேறினால், அப்போதும் அப்போது மட்டுமே \emph{அதிகபட்ச முரணற்றது}:
\begin{enumerate}
\item $\Gamma$ முரணற்றது; மேலும்
\item $\Gamma \subsetneq \Gamma'$ எனில், $\Gamma'$ முரணுடையது.
\end{enumerate}
\end{defn}

% SEGMENT OLP-0644-S02
\begin{explain}
மேலுள்ளதற்குச் சமானமான மற்றொரு வரையறை இதுதான்: வாக்கியக்
கணம்~$\Gamma$, பின்வரும் இரண்டு நிபந்தனைகளும் நிறைவேறினால்,
அப்போதும் அப்போது மட்டுமே \emph{அதிகபட்ச முரணற்றது}:
\begin{enumerate}
\item $\Gamma$ முரணற்றது; மேலும்
\item $\Gamma \cup \{ !A \}$ முரணற்றது எனில், $!A \in \Gamma$.
\end{enumerate}
வேறு சொற்களில், ஏற்கெனவே~$\Gamma$ இல் இல்லாத !!{sentence}s ஐ
அதிகபட்ச முரணற்ற கணம்~$\Gamma$ உடன் சேர்த்தால், அதன் விளைவான பெரிய கணம்
முரணுடையதாகிவிடும்.
\end{explain}

% SEGMENT OLP-0644-S03
\begin{explain}
நிறைவு நிறுவலில் அதிகபட்ச முரணற்ற கணங்கள் முக்கியமானவை. ஏனெனில்,
!!{sentence}s கொண்ட எந்த முரணற்ற கணம்~$\Gamma$ உம் ஓர் அதிகபட்ச
முரணற்ற கணம்~$\Gamma^*$ இல் அடங்கும் என்பதை உறுதிசெய்யலாம். மேலும்,
அதிகபட்ச முரணற்ற கணம், ஒவ்வொரு !!{sentence}~$!A$ க்கும் $!A$
அல்லது அதன் மறுப்பு~$\lnot !A$ இவற்றில் ஒன்றைக் கொண்டிருக்கும்.
இது குறிப்பாக அணு !!{sentence}s க்கும் பொருந்துகிறது. எனவே
!!{constant}s ஐ ஏற்றவாறு சேர்த்து விரிவாக்கிய மொழியில் உள்ள
அதிகபட்ச முரணற்ற கணத்திலிருந்து, எந்த அணு !!{sentence}s ஆனவை
$\Gamma^*$ இல் உள்ளன என்பதன்படி !!{predicate}s இன் பொருள்விளக்கம்
வரையறுக்கப்படும் !!a{structure} ஒன்றைக் கட்டமைக்கலாம். இந்த
!!{structure} ஆனது~$\Gamma^*$ இல் உள்ள அனைத்து !!{sentence}s ஐயும்
(ஆகவே~$\Gamma$ இல் உள்ள அனைத்தையும்) மெய்யாக்கும் என்று பின்னர்
காட்டலாம். இந்தப் பின்னைய கூற்றை நிறுவ, $\lnot !A \in
\Gamma^*$ அப்போதும் அப்போது மட்டுமே $!A \notin \Gamma^*$;
$(!A \lor !B) \in \Gamma^*$ அப்போதும் அப்போது மட்டுமே $!A
\in \Gamma^*$ அல்லது $!B \in \Gamma^*$; இதுபோன்ற பண்புகள் தேவை.
\end{explain}

% SEGMENT OLP-0644-S04
\begin{prop}
\ollabel{prop:mcs}
$\Gamma$ அதிகபட்ச முரணற்றது என்க. அப்போது:
\begin{enumerate}
\item \ollabel{prop:mcs-prov-in} $\Gamma \Proves !A$ எனில், $!A \in
  \Gamma$.

\item \ollabel{prop:mcs-either-or} எந்த $!A$ க்கும், $!A \in
  \Gamma$ அல்லது $\lnot !A \in \Gamma$.

\tagitem{prvAnd}{\ollabel{prop:mcs-and} $(!A \land !B) \in \Gamma$
  அப்போதும் அப்போது மட்டுமே $!A \in \Gamma$, $!B \in \Gamma$ இரண்டும்.}{}

\tagitem{prvOr}{\ollabel{prop:mcs-or} $(!A \lor !B) \in \Gamma$
  அப்போதும் அப்போது மட்டுமே $!A \in \Gamma$ அல்லது $!B \in \Gamma$.}{}

\tagitem{prvIf}{\ollabel{prop:mcs-if} $(!A \lif !B) \in \Gamma$
  அப்போதும் அப்போது மட்டுமே $!A \notin \Gamma$ அல்லது $!B \in \Gamma$.}{}
\end{enumerate}
\end{prop}

% SEGMENT OLP-0644-S05
\begin{proof}
பின்வரும் அனைத்திலும் $\Gamma$ அதிகபட்ச முரணற்றது எனக் கொள்வோம்.
\begin{enumerate}
\item $\Gamma \Proves !A$ எனில், $!A \in \Gamma$.

$\Gamma \Proves !A$ என்க. இதற்கு மாறாக $!A
\notin \Gamma$ எனக் கொள்வோம். $\Gamma$ அதிகபட்ச முரணற்றது என்பதால்,
$\Gamma \cup \{!A\}$ முரணுடையது; எனவே $\Gamma \cup \{!A\} \Proves
\lfalse$. பின்வரும் மேற்கோளின்படி,
\tagrefs{prfSC/{fol:seq:prv:prop:provability-contr},prfND/{fol:ntd:prv:prop:provability-contr}},
$\Gamma$ முரணுடையது. இது $\Gamma$ முரணற்றது என்ற கருதுகோளுக்கு
முரணாகும். ஆகவே $!A \notin
\Gamma$ என்பது இயலாது; எனவே $!A \in \Gamma$.

\item எந்த $!A$ க்கும், $!A \in \Gamma$ அல்லது $\lnot !A \in \Gamma$.

இதற்கு மாறாக, ஏதோ~$!A$ க்கு $!A \notin \Gamma$ மற்றும்
$\lnot !A \notin \Gamma$ இரண்டும் எனக் கொள்வோம். $\Gamma$ அதிகபட்ச
முரணற்றது என்பதால், $\Gamma \cup \{!A\}$ மற்றும்
$\Gamma \cup \{\lnot !A\}$ இரண்டும் முரணுடையவை. எனவே
$\Gamma \cup \{!A\} \Proves \lfalse$ மற்றும் $\Gamma \cup
\{\lnot !A\} \Proves \lfalse$. பின்வரும் மேற்கோளின்படி,
\tagrefs{prfSC/{fol:seq:prv:prop:provability-exhaustive},prfND/{fol:ntd:prv:prop:provability-exhaustive}},
$\Gamma$ முரணுடையது; இது முரண்பாடு. ஆகவே அத்தகைய
!!a{sentence}~$!A$ இருக்க முடியாது. இதனால் ஒவ்வொரு $!A$ க்கும்,
$!A \in \Gamma$ அல்லது $\lnot !A
\in \Gamma$.

% SEGMENT OLP-0644-S06
\tagitem{defAnd}{}{%
\iftag{probAnd}{பயிற்சி.}{%
$(!A \land !B) \in \Gamma$ அப்போதும் அப்போது மட்டுமே
$!A \in \Gamma$, $!B \in \Gamma$ இரண்டும்:

முன்னோக்கிய திசைக்கு, $(!A \land !B) \in \Gamma$ என்க. அப்போது
$\Gamma \Proves !A \land !B$. பின்வரும் மேற்கோளின்படி,
\tagrefs{prfSC/{fol:seq:ppr:prop:provability-land-left},prfND/{fol:ntd:ppr:prop:provability-land-left}},
$\Gamma \Proves !A$ மற்றும் $\Gamma \Proves !B$. தேவையானவாறு,
\olref{prop:mcs-prov-in} படி $!A \in \Gamma$ மற்றும் $!B \in \Gamma$.

மறுதிசைக்கு, $!A \in \Gamma$ மற்றும் $!B \in
\Gamma$ என்க. அப்போது $\Gamma \Proves !A$ மற்றும்
$\Gamma \Proves !B$. பின்வரும் மேற்கோளின்படி,
\tagrefs{prfSC/{fol:seq:ppr:prop:provability-land-right},prfND/{fol:ntd:ppr:prop:provability-land-right}},
$\Gamma \Proves !A \land !B$. \olref{prop:mcs-prov-in} படி,
$(!A \land !B) \in \Gamma$.}}

% SEGMENT OLP-0644-S07
\tagitem{defOr}{}{%
\iftag{probOr}{பயிற்சி.}{%
$(!A \lor !B) \in \Gamma$ அப்போதும் அப்போது மட்டுமே
$!A \in \Gamma$ அல்லது $!B \in \Gamma$.

முன்னோக்கிய திசையின் எதிர்மறைநேர்மாற்றை நிறுவ, $!A
\notin \Gamma$ மற்றும் $!B \notin \Gamma$ என்க. நாம் $(!A \lor
!B) \notin \Gamma$ என்பதைக் காட்ட வேண்டும். $\Gamma$ அதிகபட்ச
முரணற்றது என்பதால், $\Gamma \cup \{!A\} \Proves \lfalse$ மற்றும்
$\Gamma \cup \{!B\} \Proves \lfalse$. பின்வரும் மேற்கோளின்படி,
\tagrefs{prfSC/{fol:seq:ppr:prop:provability-lor},prfND/{fol:ntd:ppr:prop:provability-lor}},
$\Gamma \cup \{(!A \lor !B)\}$ முரணுடையது. ஆகவே தேவையானபடி,
$(!A \lor !B)
\notin \Gamma$.

மறுதிசைக்கு, $!A \in \Gamma$ அல்லது $!B \in
\Gamma$ என்க. அப்போது $\Gamma \Proves !A$ அல்லது
$\Gamma \Proves !B$. பின்வரும் மேற்கோளின்படி,
\tagrefs{prfSC/{fol:seq:ppr:prop:provability-lor},prfND/{fol:ntd:ppr:prop:provability-lor}},
$\Gamma \Proves !A \lor !B$. \olref{prop:mcs-prov-in} படி,
தேவையானபடி $(!A \lor
!B) \in \Gamma$.}}

% SEGMENT OLP-0644-S08
\tagitem{defIf}{}{%
\iftag{probIf}{பயிற்சி.}{%
$(!A \lif !B) \in \Gamma$ அப்போதும் அப்போது மட்டுமே
$!A \notin \Gamma$ அல்லது $!B \in \Gamma$:

முன்னோக்கிய திசைக்கு, $(!A \lif !B) \in \Gamma$ என்க. இதற்கு
மாறாக $!A \in \Gamma$ மற்றும் $!B \notin \Gamma$ எனக் கொள்வோம்.
இந்தக் கருதுகோள்களின்படி $\Gamma \Proves !A \lif !B$ மற்றும்
$\Gamma \Proves !A$. பின்வரும் மேற்கோளின்படி,
\tagrefs{prfSC/{fol:seq:ppr:prop:provability-lif-left},prfND/{fol:ntd:ppr:prop:provability-lif-left}},
$\Gamma \Proves !B$. ஆனால் \olref{prop:mcs-prov-in} படி $!B \in
\Gamma$; இது $!B \notin \Gamma$ என்ற கருதுகோளுக்கு முரணாகும்.

மறுதிசையில், முதலில் $!A \notin
\Gamma$ என்ற நிலையை எடுத்துக்கொள்வோம். \olref{prop:mcs-either-or}
படி $\lnot !A \in \Gamma$; ஆகவே $\Gamma \Proves \lnot !A$.
பின்வரும் மேற்கோளின்படி,
\tagrefs{prfSC/{fol:seq:ppr:prop:provability-lif-right},prfND/{fol:ntd:ppr:prop:provability-lif-right}},
$\Gamma \Proves !A \lif !B$. மீண்டும் \olref{prop:mcs-prov-in}
படி, தேவையானபடி $(!A \lif !B) \in \Gamma$.

இப்போது $!B \in \Gamma$ என்ற நிலையை எடுத்துக்கொள்வோம். அப்போது
$\Gamma \Proves !B$. பின்வரும் மேற்கோளின்படி,
\tagrefs{prfSC/{fol:seq:ppr:prop:provability-lif-right},prfND/{fol:ntd:ppr:prop:provability-lif-right}},
$\Gamma \Proves !A \lif !B$. \olref{prop:mcs-prov-in} படி,
$(!A \lif
!B) \in \Gamma$.}}
\end{enumerate}
\end{proof}

% SEGMENT OLP-0644-S09
\begin{prob}
\olref[fol][com][mcs]{prop:mcs} இன் நிறுவலை நிறைவு செய்க.
\end{prob}


% END SOURCE UNIT OLP-0644

\clearpage
\noindent மூல அலகு: \texttt{OLP-0645}\par

% BEGIN SOURCE UNIT OLP-0645 content/first-order-logic/syntax-and-semantics/introduction.tex
% Exact modular target bytes: 9518; SHA256: a378bf53bd76cc2760151e8a96f14ce8d19cb70d00689ae781150bcc98ae6b4f



% SEGMENT OLP-0645-S01
\olfileid{fol}{syn}{int}

\olsection{அறிமுகம்}

% SEGMENT OLP-0645-S02
முதல்தர தருக்கத்தின் கோட்பாட்டையும் மேல்கோட்பாட்டையும் வளர்க்க, அதன்
வெளிப்பாடுகளின் தொடரமைப்பையும் பொருண்மையியலையும் முதலில் வரையறுக்க வேண்டும்.
முதல்தர தருக்கத்தின் வெளிப்பாடுகள் சொற்களும் !!{formula}s உம் ஆகும். சொற்கள்
!!{variable}s, !!{constant}s மற்றும் !!{function}s ஆகியவற்றிலிருந்து
உருவாக்கப்படுகின்றன. !!^{formula}s, மறுபுறம், !!{predicate}s மற்றும் சொற்களுடன்
சேர்ந்து உருவாகின்றன (இவை மிகச் சிறிய, ``அணு'' !!{formula}s); பின்னர் அணு
!!{formula}s இலிருந்து தருக்கத் தொடுப்புக்குறிகளையும் அளவையடைகளையும் கொண்டு
மிகச் சிக்கலானவற்றை உருவாக்கலாம். உருவாக்க விதிகளை எழுதுவதற்குப் பல வழிகள்
உள்ளன; அவற்றில் ஒரு சாத்தியமான வழியை மட்டுமே தருகிறோம். பிற அமைப்புகள்
வேறு குறிகளைத் தேர்ந்தெடுக்கலாம், முதன்மையான தொடுப்புக்குறிகளின் வேறு
தொகுப்பைத் தேர்ந்தெடுக்கலாம், அடைப்புக்குறிகளை வேறுவிதமாகப் பயன்படுத்தலாம்
(அல்லது போலிஷ் குறியீட்டில் போல அவற்றைப் பயன்படுத்தாமலும் இருக்கலாம்).
எல்லா அணுகுமுறைகளுக்கும் பொதுவானது என்னவென்றால், உருவாக்க விதிகள் சொற்களின்
மற்றும் !!{formula}s இன் தொகுப்பை \emph{தொகுத்தறிதல் வழியாக} வரையறுக்கின்றன.
அது முறையாகச் செய்யப்பட்டால், ஒவ்வொரு வெளிப்பாடும் உருவாக்க விதிகளின்படி
அடிப்படையில் ஒரே ஒரு வழியில் மட்டுமே பெறப்படும். வெளிப்பாடுகள்
\emph{தனித்த வாசிப்புடையவை} ஆகும் என்ற இந்தத் தொகுத்தறிதல் வரையறை,
அதே முறையை---தொகுத்தறிதல் வரையறையைப் பயன்படுத்தி---அவற்றுக்கு
பொருள்களையும் வழங்க வழி செய்கிறது.

% SEGMENT OLP-0645-S03
வெளிப்பாடுகளுக்குப் பொருள் வழங்குவது பொருண்மையியலின் களம். பொருண்மையியலின்
மையக் கருத்து !!a{structure} இல் நிறைவுறுத்தல். !!^a{structure} மொழியின்
கட்டுமானக் கூறுகளுக்குப் பொருள் வழங்குகிறது: !!a{domain} என்பது வெறுமையற்ற
பொருட்களின் தொகுப்பு. அளவையடைகள் இப்பொருட்களத்தில் இயங்குவதாகப்
பொருள்படுத்தப்படுகின்றன; !!{constant}s பொருட்களத்தின் உறுப்புகளுக்கு
ஒதுக்கப்படுகின்றன; !!{function}s !!{domain} இலிருந்து அதற்கே செல்லும்
சார்புகளுக்கு ஒதுக்கப்படுகின்றன; !!{predicate}s !!{domain} மீது உள்ள
தொடர்புகளுக்கு ஒதுக்கப்படுகின்றன. !!{domain} மற்றும் அடிப்படைச்
சொற்களஞ்சியத்திற்கான ஒதுக்கீடுகள் சேர்ந்து !!a{structure} ஐ உருவாக்குகின்றன.
!!^{variable}s !!{formula}s இல் தோன்றலாம்; அவற்றுக்கு பொருண்மையியல்
வழங்க, !!{domain} இன் !!{element}s ஐ அவற்றுக்கும் ஒதுக்க வேண்டும்---இதுவே
மாறி ஒதுக்கீடு. இறுதியாக, நிறைவுறுத்தல் தொடர்பு இக்கூறுகளை ஒன்றிணைக்கிறது.
!!^a{formula} ஒன்று !!a{structure}~$\Struct{M}$ இல்,
மாறி ஒதுக்கீடு~$s$ சார்பாக நிறைவுறுத்தப்படலாம்; இதை
$\Sat{M}{!A}[s]$ என்று எழுதுகிறோம். இத்தொடர்பும்~$!A$ இன் கட்டமைப்பின்
மீது தொகுத்தறிதலால் வரையறுக்கப்படுகிறது; உதாரணமாக, தருக்கத் தொடுப்புக்குறிகளின்
மெய்மை அட்டவணைகளைப் பயன்படுத்தி $!A \land !B$ இன் நிறைவுறுத்தலை
$!A$ மற்றும் ~$!B$ ஆகியவற்றின் நிறைவுறுத்தல் (அல்லது நிறைவுறுத்தாமை) அடிப்படையில்
வரையறுக்கிறோம். மாறி ஒதுக்கீடு, !!{formula}~$!A$
!!a{sentence} ஆக, அதாவது கட்டற்ற மாறிகள் இல்லாததாக, இருந்தால் பொருத்தமற்றது
என்பது பின்னர் தெரியவரும்; எனவே !!{sentence}s !!{structure}s இல்
வெறுமனே நிறைவுறுத்தப்படுகின்றனவா இல்லையா என்று பேசலாம்.

% SEGMENT OLP-0645-S04
வாக்கியங்களுக்கான நிறைவுறுத்தல் தொடர்பு~$\Sat{M}{!A}$ அடிப்படையில்,
செல்லுபடியாகும் தன்மை, உட்கிடை, நிறைவுசெய்யத்தக்க தன்மை ஆகிய அடிப்படைப்
பொருண்மையியல் கருத்துகளை வரையறுக்கலாம். ஒரு வாக்கியம் செல்லுபடியாகும்,
$\Entails !A$, எனில் ஒவ்வொரு கட்டமைப்பும் அதை நிறைவுறுத்தும்.
!!{sentence}s கொண்ட ஒரு கணத்தால் அது உட்கிடையாகிறது,
$\Gamma \Entails !A$, எனில்~$\Gamma$ இல் உள்ள எல்லா !!{sentence}s ஐயும்
நிறைவுறுத்தும் ஒவ்வொரு !!{structure} உம்~$!A$ ஐயும் நிறைவுறுத்தும்.
மேலும், வாக்கியங்கள் கொண்ட ஒரு கணம் நிறைவுசெய்யத்தக்கது எனில், அதிலுள்ள
எல்லா !!{sentence}s ஐயும் ஒரே நேரத்தில் நிறைவுறுத்தும் !!{structure} ஒன்று
உள்ளது. !!{formula}s தொகுத்தறிதலால் வரையறுக்கப்பட்டுள்ளன; நிறைவுறுத்தலும்
!!{formula}s இன் கட்டமைப்பின் மீது தொகுத்தறிதலால் வரையறுக்கப்படுகிறது.
எனவே, நமது பொருண்மையியலின் பண்புகளை நிறுவவும் வரையறுக்கப்பட்ட
பொருண்மையியல் கருத்துகளைத் தொடர்புபடுத்தவும் தொகுத்தறிதலைப் பயன்படுத்தலாம்.


% END SOURCE UNIT OLP-0645

\clearpage
\noindent மூல அலகு: \texttt{OLP-0646}\par
\begingroup
\TADriverHeadingMode
\olchapterid{fol}{syn}

% BEGIN SOURCE UNIT OLP-0646 content/first-order-logic/syntax-and-semantics/syntax-and-semantics.tex
% Exact modular target bytes: 718; SHA256: 8d0f26f61639efd51a2de832519cf54920e61a5b7646501c72eebdbd29347bf3



% SEGMENT OLP-0646-S01
\olchapter{fol}{syn}{தொடரமைப்பும் பொருண்மையியலும்}

\TADriverImport{introduction}

\TADriverImport{first-order-languages}

\TADriverImport{terms-formulas}

\TADriverImport{unique-readability}

\TADriverImport{main-operator}

\TADriverImport{subformulas}

\TADriverImport{formation-sequences}

\TADriverImport{free-vars-sentences}

\TADriverImport{substitution}

\TADriverImport{structures}

\TADriverImport{covered-structures}

\TADriverImport{satisfaction}

\TADriverImport{assignments}

\TADriverImport{extensionality}

\TADriverImport{semantic-notions}

\OLEndChapterHook


% END SOURCE UNIT OLP-0646

\endgroup


\chapter{பிற தருக்கப் பகுதிகள்}
% The definition of C precedes the theorem about representability in Q.
\clearpage
\noindent மூல அலகு: \texttt{OLP-0648}\par

% BEGIN SOURCE UNIT OLP-0648 content/incompleteness/representability-in-q/c.tex
% Exact modular target bytes: 2711; SHA256: 4e9682d78ec1d51f5e174542b12cb12d7b966265b9aca73a2d99c44d2aac6b7c



% SEGMENT OLP-0648-S01
\olfileid{inc}{req}{cee}
\olsection{சார்புக் கணம் $C$}

$C$ என்பது பின்வருவனவற்றைக் கொண்ட மிகச் சிறிய சார்புக் கணமாக இருக்கட்டும்:
\begin{enumerate}
\item $0$,
\item அடுத்தெண்,
\item கூட்டல்,
\item பெருக்கல்,
\item வீழ்த்தல்கள், மற்றும்
\item சமத்துவத்தின் சிறப்பியல்புச் சார்பு, $\Char{=}$;
\end{enumerate}
மேலும் பின்வருவனவற்றின் கீழ் அது மூடப்பட்டிருக்கட்டும்:
\begin{enumerate}
\item சேர்ப்பு, மற்றும்
\item ஒழுங்கான சார்புகளுக்குப் பயன்படுத்தப்படும் வரம்பற்ற தேடல்.
\end{enumerate}

% SEGMENT OLP-0648-S02
கடைசிக் கட்டுப்பாட்டின் பொருள், விளைவு எங்கும் வரையறுக்கப்பட்டதாக
இருக்கும்போது மட்டுமே $\mu$ செயல்பாட்டைப் பயன்படுத்தலாம் என்பதே.
இதைப் \emph{பொது மீள்வரையறைச்} சார்புகளின் வரையறையுடன் ஒப்பிடுக:
இங்கு கூட்டல், பெருக்கல், $\Char{=}$ ஆகியவற்றைச் சேர்த்துள்ளோம்;
ஆனால் அடிப்படை மீள்வரையறையை நீக்கியுள்ளோம்.

கூட்டல், பெருக்கல், $\Char{=}$ ஆகியவை மீள்வரையறைச் சார்புகள் என்பதால்,
$C$ இலுள்ள அனைத்தும் மீள்வரையறைச் சார்புகளே என்பது தெளிவு.
இதன் மறுதிசையும் உண்மை என்று காட்டுவோம்; அதாவது $C$ இல் உள்ள மற்ற
செயல்முறைகளைக் கொண்டு அடிப்படை மீள்வரையறையைச் செய்ய முடியும்.


% END SOURCE UNIT OLP-0648

\clearpage
\noindent மூல அலகு: \texttt{OLP-0647}\par

% BEGIN SOURCE UNIT OLP-0647 content/incompleteness/representability-in-q/c-representable.tex
% Exact modular target bytes: 18390; SHA256: 683a5b0e79cebf557b1117c44cf1465fb45b66c9fbf2afe503721fca666d444d



% SEGMENT OLP-0647-S01
\olfileid{inc}{req}{cre}
\olsection{$C$ இலுள்ள சார்புகள் $\Th{Q}$ இல் பிரதிநிதித்துவப்படுத்தத்தக்கவை}

$C$ இலுள்ள ஒவ்வொரு சார்பும் $\Th{Q}$ இல் பிரதிநிதித்துவப்படுத்தத்தக்கது
என்று நாம் காட்ட வேண்டும். இறுதியில், $C$ இலுள்ள ஒவ்வொரு $k$-இடச்
சார்பு $f(x_0,\dots,x_{k-1})$ க்கும் அதைப் பிரதிநிதித்துவப்படுத்தும்
வாய்பாடு $!A_f(x_0,\dots,x_{k-1},y)$ ஒன்றை எவ்வாறு ஒதுக்குவது என்று
காட்ட வேண்டும்.

% SEGMENT OLP-0647-S02
\begin{lem}
இயல் எண்கள் $n$, $m$ கொடுக்கப்பட்டிருக்க, $n \neq m$ எனில்
$Q \Proves \num n \neq \num m$.
\end{lem}

% SEGMENT OLP-0647-S03
\begin{proof}
ஒவ்வொரு $m$ க்கும், $n \neq m$ எனில்
$Q \Proves \eq/[\num n][\num m]$ என்று காட்ட $n$ மீது
தொகுத்தறிதலைப் பயன்படுத்துவோம்.

அடிப்படை நிலையில் $n = 0$. $m$ ஆனது $0$ க்குச் சமமல்ல எனில், ஏதோ
இயல் எண் $k$ க்கு $m = k + 1$. நம்மிடம்
$\lforall[x][\eq/[0][x']]$ என்று கூறும் அடிகோள் உள்ளது.
அளவையடை அடிகோளைப் பயன்படுத்தி $x$ க்குப் பதிலாக $\num k$ ஐ இட்டால்,
$\eq/[0][\num k']$ என்று முடிக்கலாம். ஆனால் $\num k'$ என்பது
$\num m$ தான்.

தொகுத்தறிதல் படியில், கூற்று $n$ க்கு உண்மை என்று எடுகொண்டு,
$n+1$ ஐக் கருதுவோம். $m$ தன்னிச்சையான இயல் எண்ணாக இருக்கட்டும்.
இரண்டு வாய்ப்புகள் உள்ளன: $m = 0$, அல்லது ஏதோ $k$ க்கு
$m = k+1$. முதல் நேர்வை மேலே கண்டவாறே கையாளலாம். இரண்டாம் நேர்வில்
$n+1 \neq k+1$ என்று கொள்க. அப்போது $n \neq k$.
$n$ க்கான தொகுத்தறிதல் எடுகோளால்
$\Th{Q} \Proves \eq/[\num n][\num k]$. நம்மிடம்
$\lforall[x][\lforall[y][\eq[x'][y'] \lif \eq[x][y]]]$
என்று கூறும் அடிகோள் உள்ளது. அளவையடை அடிகோளைப் பயன்படுத்தி,
$\eq[\num n'][\num k'] \lif \eq[\num n][\num
  k]$ ஐப் பெறுகிறோம். கூற்றுத் தருக்கத்தைப் பயன்படுத்தி,
$\Th{Q}$ இல்
$\eq/[\num n][\num k] \lif \eq/[\num n'][\num k']$
என்று முடிக்கலாம். விதிப்படி உட்படுத்தலால்
$\eq/[\num n'][\num k']$ ஐப் பெறுகிறோம். $\num k'$ என்பது
$\num m$ என்பதால், இதுவே வேண்டியது.
\end{proof}

% SEGMENT OLP-0647-S04
இந்தத் துணைத்தேற்றம் கூறுவது மிக அதிகமல்ல: வெவ்வேறு எண்ணுருக்கள்
வெவ்வேறு பொருட்களைக் குறிக்கின்றன என்பதை $\Th{Q}$ நிறுவ முடியும்
என்பதே அதன் சாரம். எடுத்துக்காட்டாக, $0'' \neq 0'''$ ஐ
$\Th{Q}$ நிறுவுகிறது. ஆனால் இது பொதுவாகப் பொருந்துகிறது என்று
காட்டுவதற்குச் சற்றுக் கவனம் தேவை. நாம் தொகுத்தறிதலைப் பயன்படுத்தினாலும்,
அது $\Th{Q}$ க்கு \emph{வெளியே} நடைபெறும் தொகுத்தறிதல் என்பதையும்
கவனிக்க.

சுழியம், அடுத்தெண், கூட்டல், பெருக்கல், சமத்துவத்தின் சிறப்பியல்புச்
சார்பு, வீழ்த்தல்கள் ஆகியவற்றை நம்மால் பிரதிநிதித்துவப்படுத்த முடியும்.
ஒவ்வொரு நேர்விலும் பொருத்தமான பிரதிநிதித்துவ வாய்பாடு முற்றிலும்
நேரடியானது; எடுத்துக்காட்டாக, சுழியம் பின்வரும் வாய்பாட்டால்
பிரதிநிதித்துவப்படுத்தப்படுகிறது:
\[
y = 0,
\]
அடுத்தெண் பின்வரும் வாய்பாட்டால் பிரதிநிதித்துவப்படுத்தப்படுகிறது:
\[
x_0' = y,
\]
கூட்டல் பின்வரும் வாய்பாட்டால் பிரதிநிதித்துவப்படுத்தப்படுகிறது:
\[
x_0 + x_1 = y.
\]
தொடர்புடைய கூற்றுகளை $\Th{Q}$ நிறுவ முடியும் என்று காட்டுவதில்தான்
பணி உள்ளது; எடுத்துக்காட்டாக, மேலுள்ள வாய்பாடு கூட்டலைப்
பிரதிநிதித்துவப்படுத்துகிறது என்பதற்கு, ஒவ்வோர் இயல் எண் சோடி
$m$,~$n$ க்கும் $\Th{Q}$ பின்வருவதை நிறுவுகிறது என்று காட்ட வேண்டும்:
\[
\eq[\num n + \num m][\num {n+m}]
\]
மேலும் பின்வருவதையும் நிறுவுகிறது:
\[
\lforall[y][(\eq[(\num n + \num m)][y] \lif \eq[y][\num {n+m}])].
\]

% SEGMENT OLP-0647-S05
சேர்ப்பு பற்றி என்ன? $h$ பின்வருமாறு வரையறுக்கப்படுகிறது என்று கொள்க:
\[
h(x_0,\dots,x_{l-1}) = f(g_0(x_0,\dots,x_{l-1}), \dots,
g_{k-1}(x_0,\dots,x_{l-1})).
\]
இங்கு $f,g_0,\dots,g_{k-1}$ ஆகிய சார்புகளை முறையே
பிரதிநிதித்துவப்படுத்தும் $!A_f, !A_{g_0}, \dots,
!A_{g_{k-1}}$ என்ற வாய்பாடுகளை ஏற்கெனவே கண்டுள்ளோம். அப்போது
$h$ ஐப் பிரதிநிதித்துவப்படுத்தும் $!A_h$ என்ற வாய்பாட்டை,
$!A_h(x_0,\dots,x_{l-1},y)$ பின்வருவதாக வரையறுப்பதன் மூலம்
அமைக்கலாம்:
% SOURCE CORRECTION TA-INC-041: keep all representing conjuncts in the quantifier scope.
\begin{multline*}
  \lexists[z_0\dots][\lexists[z_{k-1}][(!A_{g_0}(x_0,\dots,x_{l-1},z_0) \land
  \dots \land !A_{g_{k-1}}(x_0,\dots,x_{l-1},z_{k-1}) \land\\
  !A_f(z_0,\dots,z_{k-1},y))]]
\end{multline*}

% SEGMENT OLP-0647-S06
இறுதியாக வரம்பற்ற தேடலைக் கருதுவோம். $g(x,\vec z)$ ஒழுங்கானதாகவும்,
$!A_g(x,\vec z,y)$ என்ற வாய்பாட்டால் $\Th{Q}$ இல்
பிரதிநிதித்துவப்படுத்தத்தக்கதாகவும் இருக்கட்டும். $f$ ஆனது
$f(\vec z) = \mu x \; g(x,\vec z)$ என்பதால் வரையறுக்கப்படட்டும்.
$f$ ஐப் பிரதிநிதித்துவப்படுத்தும் $!A_f(\vec z,y)$ என்ற
வாய்பாட்டைக் கண்டறிய விரும்புகிறோம். இயல்பான தேர்வு இதுதான்:
\[
!A_f(\vec z,y) \ident !A_g(y,\vec z,0) \land \lforall[w][(w < y
\lif \lnot !A_g(w,\vec z,0))].
\]

இது செயல்படுகிறது என்பதைச் சில கூடுதல் துணைத்தேற்றங்களின் உதவியால்
காட்டலாம்; எடுத்துக்காட்டாக:

\begin{lem}
  ஒவ்வொரு மாறி $x$ க்கும் ஒவ்வோர் இயல் எண் $n$ க்கும்,
  $\Th{Q}$ ஆனது $\eq[(x'
  + \num n)][(x + \num n)']$ ஐ நிறுவுகிறது.
\end{lem}

இது, $\Th{Q}$ ஆனது
$\lforall[x][\lforall[y][(x' + y) = (x +y)']]$
ஐ நிறுவுகிறது என்று கூறுவதைவிட வலுக்குறைந்தது என்பதை மீண்டும்
குறிப்பிட வேண்டும்; அதை நிறுவ முடியும் என்ற வலுவான கூற்று பொய்.

% SEGMENT OLP-0647-S07
\begin{proof}
வழக்கம்போல $n$ மீது தொகுத்தறிதலால் நிறுவுவோம். அடிப்படை நிலையில்
$n = 0$; $\Th{Q}$ ஆனது
$\eq[(x'+0)][(x+0)']$ ஐ நிறுவுகிறது என்று காட்ட வேண்டும்.
ஆனால் நம்மிடம் பின்வருவன உள்ளன:
\begin{eqnarray*}
& & \eq[(x' + 0)][x'] \quad \text{அடிகோள் 4 இன்படி} \\
& & \eq[(x + 0)][x] \quad \text{அடிகோள் 4 இன்படி} \\
& & \eq[(x+0)'][x'] \quad \text{இரண்டாம் வரியின்படி} \\
& & \eq[(x' + 0)][(x + 0)'] \quad \text{முதல், மூன்றாம் வரிகளின்படி}
\end{eqnarray*}
தொகுத்தறிதல் படியில், $\Th{Q}$ இல்
$\eq[(x' + \num n)][(x + \num n)']$ ஐ
வருவித்துள்ளோம் என்று எடுகொள்கிறோம். $\num{n+1}$ என்பது
$\num n'$ என்பதால், $\Th{Q}$ ஆனது
$\eq[(x' + \num n')][(x + \num n')']$
ஐ நிறுவுகிறது என்று காட்ட வேண்டும். பின்வரும் சமத்துவத் தொடரை
$\Th{Q}$ இல் வருவிக்கலாம்:
\begin{eqnarray*}
(x' + \num n') & \eq & (x' + \num n)' \quad \text{அடிகோள் 5 இன்படி} \\
& \eq & (x + \num n')' \quad \text{தொகுத்தறிதல் எடுகோளின்படி}
\end{eqnarray*}
\end{proof}

% SEGMENT OLP-0647-S08
\begin{lem}
\begin{enumerate}
\item $\Th{Q}$ ஆனது $\lnot (x < \num 0)$ ஐ நிறுவுகிறது.
\item ஒவ்வோர் இயல் எண் $n$ க்கும், $\Th{Q}$ பின்வருவதை நிறுவுகிறது:
\[
x < \num {n+1} \lif (\eq[x][\num 0] \lor \dots \lor \eq[x][\num n]).
\]
\end{enumerate}
\end{lem}

% SEGMENT OLP-0647-S09
\begin{proof}
1 ஐயும் 2 இன் ஒரு பகுதியையும் முறைசாரா வகையில் நிறுவுவோம்
(அதாவது முறைசார் !!{derivation} ஐ அமைப்பதற்கான குறிப்புகளை மட்டுமே
தருவோம்).

1 ஆம் பகுதிக்கு, $<$ இன் வரையறையின்படி
$\Th{Q}$ இல்
$\lnot \lexists[y][\eq[(y'+x)][0]]$
ஐ நிறுவ வேண்டும்; இது முதல்தர தருக்கத்தின் அடிகோள்களையும் விதிகளையும்
பயன்படுத்தினால்
$\lforall[y][\eq/[(y' +
    x)][0]]$
என்பதற்குச் சமம். எண்ணம் இதுதான்: $\eq[(y'+x)][0]$
என்று கொள்க. $x$ ஆனது 0 எனில், $\eq[(y' + 0)][0]$
கிடைக்கும். ஆனால் $\Th{Q}$ இன் அடிகோள் 4 இன்படி
$\eq[(y' + 0)][y']$; அடிகோள் 2 இன்படி
$\eq/[y'][0]$; இது முரண்.
$x$ ஆனது $0$ அல்ல எனில், அடிகோள் 3 இன்படி
$\eq[x][z']$ ஆகும் $z$ ஒன்று உள்ளது. ஆனால் அப்போது
$\eq[(y' + z')][0]$ கிடைக்கும். அடிகோள்~5 இன்படி
$\eq[(y' + z)'][0]$ கிடைக்கிறது; இதுவும் அடிகோள்~2 உடன்
முரண்படுகிறது. இரு நேர்வுகளிலும்
$\lforall[y][\eq/[(y' +x)][0]]$ பெறப்படுகிறது.

2 ஆம் பகுதிக்கு, $n$ மீது தொகுத்தறிதலைப் பயன்படுத்துக.
$n =0$ என்ற அடிப்படை நிலையைக் கருதுவோம். அப்போது
$x < \num 1 \lif \eq[x][\num
  0]$ என்று காட்ட வேண்டும். $x < \num 1$ என்க.
$<$ க்கான வரையறுக்கும் அடிகோளின்படி
$\lexists[y][\eq[(y'+x)][0']]$ கிடைக்கும்.
அப்பண்புடைய $y$ ஒன்று இருப்பதாகக் கொள்க; எனவே
$\eq[y'+x][0']$.

$\eq[x][0]$ என்று காட்ட வேண்டும். அடிகோள் 3 இன்படி,
$x$ ஆனது 0 அல்ல எனில், அது ஏதோ $z$ க்கு $z'$ ஆகும்.
அப்போது $\eq[(y' + z')][0']$ கிடைக்கும்.
$\Th{Q}$ இன் அடிகோள்~5 இன்படி
$\eq[(y'+z)'][0']$ கிடைக்கும். அடிகோள் 1 இன்படி
$\eq[(y' +
    z)][0]$ கிடைக்கும். ஆனால், வரையறையின்படி இதற்கு
$z < 0$ என்று பொருள்; இது 1 ஆம் பகுதிக்கு முரணானது.
\end{proof}

% SEGMENT OLP-0647-S10
\begin{explain}
கணிக்கத்தக்க சார்புகளின் கணத்தை, $\Th{Q}$ இல்
பிரதிநிதித்துவப்படுத்தத்தக்க சார்புகளின் கணமாக வகைப்படுத்தலாம் என்று
காட்டியுள்ளோம். உண்மையில், நிறுவல் இன்னும் பொதுவானது.
பிரதிநிதித்துவப்படுத்தத்தக்கது என்பதன் வரையறையிலிருந்து,
$\Th{Q}$ ஐ விரிவாக்கும் எந்தக் கோட்பாட்டிலும்
(அல்லது $\Th{Q}$ ஐப் பொருள்படுத்தக்கூடிய கோட்பாட்டிலும்)
கணிக்கத்தக்க சார்புகளைப் பிரதிநிதித்துவப்படுத்தலாம் என்பதை
எளிதில் காணலாம். மறுபுறம், நிறுவல் என்ற கருத்து கணிக்கத்தக்கதாக
இருக்கும் எந்த !!{derivation} அமைப்பிலும் பிரதிநிதித்துவப்படுத்தத்தக்க
ஒவ்வொரு சார்பும் கணிக்கத்தக்கது. ஆகவே, எடுத்துக்காட்டாக,
கணிக்கத்தக்க சார்புகளின் கணம், பியானோ எண்கணிதத்தில் அல்லது
செர்மெலோ--ஃபிராங்கல் கணக் கோட்பாட்டில்கூட
பிரதிநிதித்துவப்படுத்தத்தக்க சார்புகளின் கணமாக வகைப்படுத்தப்படலாம்.
கோடல் குறிப்பிட்டதுபோல இது சற்றே வியப்பூட்டுகிறது.
நிறுவத்தன்மையைப் பொறுத்தவரை, எந்தக் கோட்பாட்டைக் கருதுகிறோம்
என்பதைப் பொறுத்துக் கேள்விகள் மிகவும் மாறுகின்றன என்பதைக் காண்போம்:
தோராயமாகச் சொன்னால், அடிகோள்கள் வலுப்பெறும்போது கூடுதலாக நிறுவலாம்.
ஆனால் பரந்த வகை அடிகோள் கோட்பாடுகளில், பிரதிநிதித்துவப்படுத்தத்தக்க
சார்புகள் துல்லியமாகக் கணிக்கத்தக்க சார்புகளே.
\end{explain}


% END SOURCE UNIT OLP-0647

\clearpage
\noindent மூல அலகு: \texttt{OLP-0649}\par

% BEGIN SOURCE UNIT OLP-0649 content/intuitionistic-logic/semantics/propositions.tex
% Exact modular target bytes: 2175; SHA256: a2f7d73f97439eeb98ec3019d748d4da08036f4be6a745761001e78dcd2bdef3



% SEGMENT OLP-0649-S01
\olfileid{int}{sem}{pro}

\olsection{கூற்றுகள்}

!!a{relational model}~$\mModel{M}$ இல் !!a{formula}~$!A$ மெய்யாகும்
உலகங்களின் கணத்திற்குப் பெயர் தருவது சில வேளைகளில் பயனுள்ளதாக இருக்கும்.
அதை~$!A$ வரையறுக்கும் \emph{கூற்று} என்று அழைத்து,
$\Prop{M}{!A}$ என எழுதுகிறோம்.

% SEGMENT OLP-0649-S02
\begin{defn}
  $V$ என்ற சார்பை, $\Prop{M}{!A} \colon \Frm
  \to \Pow{W}$ என்ற சார்பாகப் பின்வருமாறு தொகுத்தறிதல் முறையில்
  விரிவுபடுத்துகிறோம்:
  \begin{enumerate}
  \item $\Prop{M}{\lfalse} = \emptyset$
  \item $\Prop{M}{p} = V(p)$
  \item $\Prop{M}{\lnot !B} = {}$
    \[
      \Setabs{w \in W}{\text{ஒவ்வொரு $w'$ க்கும், $Rww'$ எனில், $w'
          \notin \Prop{M}{!B}$}}.
    \]
  \item $\Prop{M}{!B \land !C} = \Prop{M}{!B} \cap \Prop{M}{!C}$.
  \item $\Prop{M}{!B \lor !C} = \Prop{M}{!B} \cup \Prop{M}{!C}$.
  \item $\Prop{M}{!B \lif !C} = {}$
    \[
      \Setabs{w \in W}{\text{ஒவ்வொரு $w'$ க்கும், $Rww'$ எனில், $w'
          \in \Prop{M}{!B}$ என்றால் $w' \in \Prop{M}{!C}$}}.
    \]
  \end{enumerate}
\end{defn}

% SEGMENT OLP-0649-S03
\begin{prop}
  \ollabel{prop:propositions-true}
  % SOURCE CORRECTION TA-IL-025: truth at a world is \mSat, not the model constructor.
  $\mSat{M}{!A}[w]$ எனில், எனில் மட்டுமே, $w \in \Prop{M}{!A}$
\end{prop}

% SEGMENT OLP-0649-S04
\begin{proof}
  பயிற்சி.
\end{proof}

% SEGMENT OLP-0649-S05
\begin{prob}
  \olref[int][sem][pro]{prop:propositions-true} ஐ நிறுவுக.
\end{prob}


% END SOURCE UNIT OLP-0649

\clearpage
\noindent மூல அலகு: \texttt{OLP-0650}\par

% BEGIN SOURCE UNIT OLP-0650 content/lambda-calculus/lambda-definability/lists.tex
% Exact modular target bytes: 3561; SHA256: 64d4722f7569c9dd595d22f3f461937307310e939bf523d5f3a69e8f14c983ca



% SEGMENT OLP-0650-S01
\olfileid{lam}{ldf}{lst}
\olsection{பட்டியல்கள்}

$[M_1, M_2, \ldots, M_n]$ என்ற பட்டியலைப் பின்வருமாறு வரையறுக்கலாம்:
\[
  [M_1, M_2, \ldots, M_n] \ident \lambd[sz][s M_1 (s M_2 (\ldots (s M_n z)))]
\]
முறைசாரா வகையில் சொன்னால், பட்டியலைக் குறிக்கும் சொல் அதன்
‘திரட்டுச் சார்பு’ ஆகும்: அது இரண்டு சொற்களை ஏற்கும் சார்பு.
முதலாவது, புதிய உறுப்பையும் இதுவரை திரட்டிய மதிப்பையும் பெற்று
புதிய திரண்ட மதிப்பைத் தரும் சார்பு; இரண்டாவது, தொடக்கத் திரண்ட
மதிப்பு. இவ்வாறு உறுப்புகளை ஒவ்வொன்றாகத் திரட்டி இறுதியில்
விளைவைத் தருகிறது.

% SEGMENT OLP-0650-S02
\begin{digress}
  சார்புசார் நிரலாக்கத்தை அறிந்தவர்களுக்கு, இந்த வரையறை
  \emph{வலப்புற மடிப்பு} என்பதற்கு ஒத்ததாகத் தோன்றும்:
  ஒரு பட்டியல், அதிலுள்ள உறுப்புகள்மீது வலப்புற மடிப்பைச் செய்யும்
  சார்பாக வரையறுக்கப்படுகிறது.
\end{digress}

% SEGMENT OLP-0650-S03
இந்தக் குறியாக்கத்தின் அடிப்படையில் பயனுள்ள சில சார்புகளை எளிதில்
வரையறுக்கலாம்:
\begin{align*}
  \fn{Sum} & \ident
  \lambd[l][l \, (\lambd[xa][\fn{Add} \, x \, a])\,  \num{0}]\\
  \fn{Len} & \ident
  \lambd[l][l \, (\lambd[xa][\fn{Add} \, a \, \num{1}]) \num{0}]
\end{align*}

$\fn{Sum}$ ஆனது சர்ச் எண்ணுருக்களின் பட்டியலில் உள்ளவற்றின்
கூட்டுத்தொகையைக் கணக்கிடுகிறது. தொடக்க மதிப்பு $\num{0}$
ஆக இருக்கும் திரட்டலைப் பட்டியல்மீது செய்கிறது; ஒவ்வோர் உறுப்புக்கும்,
புதிய மதிப்பாக $\fn{Add} \, x\, a$ ஐத் தருகிறது
(இங்கு $x$ தற்போதைய உறுப்பு, $a$ இதுவரை திரண்ட மதிப்பு).
விளைவு எல்லா உறுப்புகளின் கூட்டுத்தொகை.

% SEGMENT OLP-0650-S04
\begin{prob}
  $\fn{Len}$ என்ன செய்கிறது? விளக்குக.
\end{prob}

% SEGMENT OLP-0650-S05
\begin{prob}
  ஒரு பட்டியலைத் தலைகீழாக்கும் சார்பை வரையறுக்க.
\end{prob}

% END SOURCE UNIT OLP-0650

\clearpage
\noindent மூல அலகு: \texttt{OLP-0651}\par

% BEGIN SOURCE UNIT OLP-0651 content/lambda-calculus/syntax/conversion.tex
% Exact modular target bytes: 1533; SHA256: bcbca0e0cdda210137c3fcbd177ecd0d05d7eb1d1a2b2ae16742cede6bebbd3b



% SEGMENT OLP-0651-S01
\olfileid{lam}{int}{con}
\olsection{மாற்றமும் குறைத்தலும்}

% SEGMENT OLP-0651-S02
லாம்டா கலனத்தில் மூன்று வகையான மாற்றங்கள் அல்லது குறைத்தல்கள் உள்ளன.
நாம் காணப்போவதுபோல், இங்கு உருமாற்றம் ஒருதிசையில் நிகழ்கிறதா,
இருதிசைகளிலும் நிகழ்கிறதா என்பதே மாற்றத்திற்கும் குறைத்தலுக்கும்
இடையிலான வேறுபாடு.

% SEGMENT OLP-0651-S03
$\alpha$-மாற்றம் என்பது மாறிகளை மறுபெயரிடுவது பற்றியது.
$\lambda x. x$ உம் $\lambda y.y$ உம் ஒரே சார்பை
(சமனிச் சார்பை) குறிக்கின்றன என்று கருதுவது இயல்பானது.
இந்தக் கருத்தைப் பின்வருமாறு முறைப்படுத்துகிறோம்:

% Frozen source unit ends after this lead-in; the ordinary alpha chapter carries the definition.

% END SOURCE UNIT OLP-0651

\clearpage
\noindent மூல அலகு: \texttt{OLP-0652}\par

% BEGIN SOURCE UNIT OLP-0652 content/many-valued-logic/sequent-calculus/proof-theoretic-notions.tex
% Exact modular target bytes: 10614; SHA256: dd94f9178d9545b8dd84f429d9d40d3dad5adf2d6ab86a976cbeb03cac21fed9
% SEGMENT OLP-0652-S01


\olfileid{mvl}{seq}{prf}

\olsection{நிறுவல்-கோட்பாட்டுக் கருத்துகள்}

\begin{explain}
பல முக்கியமான பொருண்மைக் கருத்துகளை (செல்லுபடித்தன்மை, பின்விளைவு,
நிறைவுறத்தக்கதன்மை) நாம் வரையறுத்ததைப் போலவே, அவற்றுடன் தொடர்புடைய
\emph{நிறுவல்-கோட்பாட்டுக் கருத்துகளை} இப்போது வரையறுக்கிறோம். இவை
!!{structure}s[loc] !!{sentence}s நிறைவுறுவதைச் சார்ந்து வரையறுக்கப்படாமல்,
குறிப்பிட்ட தொடரணிகளின் !!{derivability} அல்லது !!{nonderivability} ஐச்
சார்ந்து வரையறுக்கப்படுகின்றன. இந்தக் கருத்துகள் ஒன்றுபடுகின்றன என்பது ஒரு
முக்கியமான கண்டுபிடிப்பு. அவை ஒன்றுபடுவதே \emph{சரித்தன்மைத் தேற்றம்} மற்றும்
\emph{நிறைவுத் தேற்றம்} ஆகியவற்றின் உள்ளடக்கம்.
\end{explain}

% SEGMENT OLP-0652-S02
\begin{defn}[தேற்றங்கள்]
!!^a{sentence}~$!A$ க்கு $\quad \Sequent !A$ தொடரணியின்
!!a{derivation} ஒன்று~$\Log{LK}$ இல் இருந்தால், அந்த வாக்கியம் ஒரு \emph{தேற்றம்}.
$!A$ தேற்றமாக இருந்தால் $\Proves !A$ என்றும், இல்லையெனில் $\Proves/ !A$
என்றும் எழுதுகிறோம்.
\end{defn}

% SEGMENT OLP-0652-S03
\begin{defn}[!!^{derivability}]
!!^a{sentence}~$!A$ என்பது !!{sentence}s கொண்ட~$\Gamma$ கணத்திலிருந்து
\emph{!!{derivable}} என்பதை $\Gamma \Proves !A$ என்று எழுதுகிறோம். இதன்
பொருள், $\Gamma_0 \subseteq \Gamma$ என்ற முடிவுறு உட்கணமும்,
$\Gamma_0$ விலுள்ள !!{sentence}s[gen] $\Gamma_0'$ என்ற தொடரும் இருந்து,
$\Log{LK}$ ஆனது $\Gamma_0' \Sequent !A$ ஐ !!{derive}s என்பதாகும். $!A$
ஆனது $\Gamma$ இலிருந்து !!{derivable} அல்ல என்றால் $\Gamma \Proves/ !A$
என்று எழுதுகிறோம்.
\end{defn}

% SEGMENT OLP-0652-S04
சுருக்கல், தளர்த்தல், இடமாற்றம் ஆகிய விதிகளால்~$\Gamma_0'$ இல் உள்ள
!!{sentence}s[gen] வரிசையும் எண்ணிக்கையும் பொருட்டல்ல: $\Gamma_0' \Sequent
!A$ என்ற தொடரணி !!{derivable} எனில்,~$\Gamma_0'$ இலுள்ள அதே !!{sentence}s[acc]
கொண்ட எந்த $\Gamma_0''$ க்கும் $\Gamma_0'' \Sequent !A$ தொடரணியும்
!!{derivable}. எடுத்துக்காட்டாக, $\Gamma_0 = \{!B, !C\}$ எனில்,
$\Gamma_0' = \tuple{!B, !B, !C}$, $\Gamma_0'' = \tuple{!C, !C, !B}$
ஆகிய இரண்டும்~$\Gamma_0$ விலுள்ள !!{sentence}s[acc] மட்டுமே கொண்ட தொடர்கள்.
ஒன்றைக் கொண்ட ஒரு தொடரணி வருவிக்கத்தக்கது எனில் மற்றொன்றைக் கொண்ட தொடரணியும்
அவ்வாறே; எடுத்துக்காட்டாக:
\begin{prooftree}
  \AxiomC{}
  \Deduce$!B, !B, !C \fCenter !A$
  \RightLabel{\LeftR{\Contraction}}
  \UnaryInf$!B, !C \fCenter !A$
  \RightLabel{\LeftR{\Exchange}}
  \UnaryInf$!C, !B \fCenter !A$
  \RightLabel{\LeftR{\Weakening}}
  \UnaryInf$!C, !C, !B \fCenter !A$
\end{prooftree}
இனிமேல் $\Gamma_0$ என்பது !!{sentence}s கொண்ட ஒரு முடிவுறு கணமாக இருந்தால்,
$\Gamma_0 \Sequent !A$ என்பது முன்தொடரில்~$\Gamma_0$ விலுள்ள !!{sentence}s[gen]
ஒரு தொடரைக் கொண்ட ஏதேனும் ஒரு தொடரணி என்று கூறுவோம்; தேவையான சுருக்கல்கள்,
இடமாற்றங்கள், தளர்த்தல்கள் ஆகியவற்றை மறைமுகமாகச் சேர்த்துக்கொள்வோம்.

% SEGMENT OLP-0652-S05
\begin{defn}[முரணின்மை]
$\Gamma$ என்ற வாக்கியக் கணம் \emph{முரணுடையது} என்பது, $\Gamma_0 \subseteq
\Gamma$ என்ற ஏதோவொரு முடிவுறு உட்கணம் இருந்து, $\Log{LK}$ ஆனது $\Gamma_0
\Sequent \quad$ ஐ !!{derive}s என்பதாகும். $\Gamma$ முரணுடையதல்ல எனில்,
அதாவது ஒவ்வொரு முடிவுறு $\Gamma_0 \subseteq \Gamma$ க்கும் $\Log{LK}$ ஆனது
$\Gamma_0 \Sequent \quad$ ஐ !!{derive} இயலாது எனில், அது \emph{முரணற்றது}
என்கிறோம்.
\end{defn}

% SEGMENT OLP-0652-S06
\begin{prop}[தற்சுட்டுப் பண்பு]
\ollabel{prop:reflexivity}
$!A \in \Gamma$ எனில், $\Gamma \Proves !A$.
\end{prop}

\begin{proof}
தொடக்கத் தொடரணி $!A \Sequent !A$ !!{derivable}; மேலும் $\{!A\}
\subseteq \Gamma$.
\end{proof}
  
% SEGMENT OLP-0652-S07
\begin{prop}[ஒருபோக்குத்தன்மை]
\ollabel{prop:monotonicity}
$\Gamma \subseteq \Delta$ மற்றும் $\Gamma \Proves !A$ எனில், $\Delta
\Proves !A$.
\end{prop}

\begin{proof}
$\Gamma \Proves !A$ என்க. அதாவது, $\Gamma_0 \subseteq \Gamma$ என்ற ஒரு
முடிவுறு உட்கணம் இருந்து, $\Gamma_0 \Sequent !A$ !!{derivable}. $\Gamma
\subseteq \Delta$ என்பதால், $\Gamma_0$ ஆனது~$\Delta$ வின் முடிவுறு உட்கணமும்
ஆகும். ஆகவே $\Gamma_0 \Sequent !A$ இன் !!{derivation} ஆனது $\Delta
\Proves !A$ என்பதையும் காட்டுகிறது.
\end{proof}

% SEGMENT OLP-0652-S08
\begin{prop}[கடப்புப் பண்பு]
\ollabel{prop:transitivity}
$\Gamma \Proves !A$ மற்றும் $\{!A\} \cup \Delta \Proves !B$ எனில்,
$\Gamma \cup \Delta \Proves !B$.
\end{prop}

\begin{proof}
$\Gamma \Proves !A$ எனில், ஒரு முடிவுறு $\Gamma_0 \subseteq \Gamma$ உம்,
$\Gamma_0 \Sequent !A$ இன் !!a{derivation}~$\pi_0$ உம் உள்ளன. $\{!A\}
\cup \Delta \Proves !B$ எனில், ஏதோவொரு முடிவுறு $\Delta_0 \subseteq
\Delta$ க்கு $!A, \Delta_0 \Sequent !B$ இன் !!a{derivation}~$\pi_1$
உள்ளது. பின்வரும் !!{derivation}[acc] கருதுக:
\begin{prooftree}
  \AxiomC{}
  \RightLabel{$\pi_0$}
  \Deduce$\Gamma_0 \fCenter !A$
  \AxiomC{}
  \RightLabel{$\pi_1$}
  \Deduce$!A, \Delta_0 \fCenter !B$
  \RightLabel{\Cut}
  \BinaryInf$\Gamma_0, \Delta_0 \fCenter !B$
\end{prooftree}
$\Gamma_0 \cup \Delta_0 \subseteq \Gamma \cup \Delta$ என்பதால், இது
$\Gamma \cup \Delta \Proves !B$ என்பதைக் காட்டுகிறது.
\end{proof}

குறிப்பாக, $\Gamma \Proves !A$ மற்றும் $!A \Proves !B$ எனில் $\Gamma
\Proves !B$ என்பது இதன் பொருள் என்பதைக் கவனிக்கவும். மேலும் $!A_1,
\dots, !A_n \Proves !B$ என்றும், ஒவ்வொரு~$i$ க்கும் $\Gamma \Proves !A_i$
என்றும் இருந்தால், $\Gamma \Proves !B$ என்பதும் பின்வருகிறது.

% SEGMENT OLP-0652-S09
\begin{prop}
\ollabel{prop:incons}
$\Gamma$ முரணுடையது என்பது, அப்போதும் அப்போது மட்டுமே, ஒவ்வொரு
வாக்கியம்~$!A$ க்கும் $\Gamma \Proves {!A}$.
\end{prop}

\begin{proof}
பயிற்சி.
\end{proof}

\begin{prob}
\olref[fol][seq][ptn]{prop:incons} ஐ நிறுவுக.
\end{prob}

% SEGMENT OLP-0652-S10
\begin{prop}[இறுக்கத்தன்மை]
  \ollabel{prop:proves-compact}
  \begin{enumerate}
  \item $\Gamma \Proves !A$ எனில், $\Gamma_0 \subseteq \Gamma$ என்ற ஒரு
    முடிவுறு உட்கணம் இருந்து, $\Gamma_0 \Proves !A$.
  \item $\Gamma$ வின் ஒவ்வொரு முடிவுறு உட்கணமும் முரணற்றது எனில்,
    $\Gamma$ முரணற்றது.
  \end{enumerate}
\end{prop}

\begin{proof}
  \begin{enumerate}
    \item $\Gamma \Proves !A$ எனில், $\Gamma_0 \subseteq \Gamma$ என்ற ஒரு
      முடிவுறு உட்கணம் இருந்து, $\Gamma_0 \Sequent !A$ தொடரணிக்கு
      !!a{derivation} உள்ளது. இதனால் $\Gamma_0 \Proves !A$.
    \item $\Gamma$ முரணுடையது எனில், $\Gamma_0 \subseteq \Gamma$ என்ற ஒரு
      முடிவுறு உட்கணம் இருந்து, $\Log{LK}$ ஆனது $\Gamma_0 \Sequent \quad$ ஐ
      !!{derive}s. ஆனால் அப்போது~$\Gamma_0$, $\Gamma$ வின் முரணுடைய ஒரு
      முடிவுறு உட்கணமாகும்.
  \end{enumerate}
\end{proof}


% END SOURCE UNIT OLP-0652

\clearpage
\noindent மூல அலகு: \texttt{OLP-0653}\par

% BEGIN SOURCE UNIT OLP-0653 content/model-theory/basics/nonstandard-arithmetic.tex
% Exact modular target bytes: 24510; SHA256: deef6dbb3b24df8c026ddd938ca4787e31a10bceefe8e7cd0e18a20a22155461
% SEGMENT OLP-0653-S01


\olfileid{mod}{bas}{nsa}
\section{எண்கணிதத்தின் நிலையற்ற மாதிரிகள்}

\begin{defn}
$\Lang{L_N}$ என்பது எண்கணிதத்தின் !!{language} ஆகட்டும். அதில்
!!{constant} $\Obj{0}$, இரு இட !!{predicate} $<$, ஓரிட
!!{function} $\prime$, இரு இட !!{function}s $+$ மற்றும் $\times$
ஆகியவை உள்ளன.
\begin{enumerate}
\item எண்கணிதத்தின் \emph{நிலையான மாதிரி} என்பது $\Lang{L_N}$
  க்கான !!{structure}~$\Struct{N}$. அதன் ஆட்களம்
  $\Nat = \{ 0, 1, 2, \dots\}$; அதில் $\Obj{0}$ என்பது $0$ ஆகவும்,
  $<$ என்பது~$\Nat$ மீதான சிறியது-என்னும் தொடர்பாகவும்,
  $\prime$, $+$, $\times$ ஆகியவை முறையே~$\Nat$ மீதான
  அடுத்தெண், கூட்டல், பெருக்கல் செயல்களாகவும் பொருள்கொள்ளப்படுகின்றன.
\item \emph{மெய் எண்கணிதம்} என்பது $\Theory{N}$ கோட்பாடு.
\end{enumerate}
\end{defn}

% SEGMENT OLP-0653-S02
$\Lang{L_N}$ இல் செயல்படும்போது, $\Obj{0}$ இற்கு அடுத்தெண்
செயலியை~$n$ முறைப் பயன்படுத்திப் பெறும்
$\Obj{0}^{\prime\dots\prime}$ வடிவிலான ஒவ்வொரு சொல்லையும்
$\num{n}$ எனச் சுருக்கி எழுதுகிறோம்.

\begin{defn}
$\Lang{L_N}$ க்கான !!{structure}~$\Struct{M}$ ஆனது
$\Struct{N} \iso \Struct{M}$ எனில், எனில் மட்டுமே,
\emph{நிலையானது}.
\end{defn}

% SEGMENT OLP-0653-S03
\begin{thm}\ollabel{thm:non-std}
மெய் எண்கணிதத்திற்கு நிலையற்ற !!{enumerable}
மாதிரிகள் உள்ளன.
\end{thm}

\begin{proof}
புதிய !!{constant}~$c$ ஒன்றைச் சேர்த்து $\Lang{L_N}$ ஐ விரிவாக்கி,
பின்வரும் கோட்பாட்டைக் கருதுக:
\[
\Theory{N} \cup \Setabs{\num{n} < c}{n \in \Nat}.
\]
இக்கோட்பாட்டின் ஒவ்வொரு முடிவுறு பகுதியும் நிறைவுறத்தக்கது.
ஆகவே இறுக்கத்தன்மையால் இதற்கு மாதிரி~$\Struct{M}$ ஒன்று உண்டு;
கீழ்நோக்கிய L\"owenheim--Skolem தேற்றத்தால் அதை
!!{enumerable} ஆக எடுத்துக்கொள்ளலாம்.
$\Domain{M}$ என்பது~$\Struct{M}$ இன் ஆட்களம் என்க. அப்போது
% SOURCE CORRECTION TA-MOD-001: use the defined L_N, and call all interpretations symbols.
$\Lang{L_N}$ இன் தருக்கமல்லாத குறியீடுகளுக்கான
$\Struct{M}$ இன் பொருள்கோள்கள் பின்வருமாறு இருக்கட்டும்:
$\mathbf{z} = \Assign{\Obj{0}}{M} \in \Domain{M}$,
% SOURCE CORRECTION TA-MOD-002: relation on the domain, not on the structure object.
${\prec} =
\Assign{<}{M} \subseteq \Domain{M}^2$,
$* = \Assign{\prime}{M} \colon
\Domain{M} \to \Domain{M}$, மற்றும்
$\oplus = \Assign{+}{M}, \otimes =
\Assign{\times}{M} \colon \Domain{M}^2 \to \Domain{M}$.
ஒவ்வொரு $x \in \Domain{M}$ க்கும், $*$ ஐ~$x$ இற்குப்
பயன்படுத்திப் பெறும் $\Domain{M}$ இன் உறுப்பை $x^*$ என எழுதுகிறோம்.

இப்போது $\Struct{N}$ க்கும் $\Struct{M}$ க்கும் இடையில்
ஓரினச்சார்பு~$h$ ஒன்று உண்டென்றால், ஏதோவொரு
$n \in \Nat$ க்கு $h(n) = \Assign{c}{M}$ ஆக வேண்டும்.
எனவே $s(x) = n$ ஆகும்~$\Struct{N}$ இன் ஏதேனும் ஒதுக்கீடு~$s$
ஐ எடுத்துக்கொள்க. அப்போது
$\Sat{N}{\eq[\num{n}][x]}[s]$; \olref[iso]{thm:isom} இன்
நிரூபிப்பின்படி
$\Sat{M}{\eq[\num{n}][x]}[h \circ s]$ உம் பொருந்தும்.
ஆகவே $\Assign{c}{M} =
\mathbf{z}^{*\cdots *}$ ($*$ ஐ~$n$ முறைத் தொடர்ச்சியாகப்
பயன்படுத்தியது). ஆனால்
$\Sat{M}{\num{n} < c}$ என்றும் $\prec$ தற்சுட்டற்றது
என்றும் கருதப்பட்டதால் இது இயலாது. எனவே $\Struct{M}$
நிலையற்றது.
\end{proof}

% SEGMENT OLP-0653-S04
\begin{prob}
கணம்~$X$ மீதான தொடர்பு~$R$ ஆனது \emph{நன்கு நிறுவப்பட்டது}
என்பது,~$R$ இல் முடிவிலா இறங்கும் சங்கிலிகள் இல்லாதபோது,
அப்போது மட்டுமே; அதாவது
$\dots x_2Rx_1Rx_0$ ஆக அமையும்
$x_0$, $x_1$, $x_2$,~\dots ஆகியவை~$X$ இல் இல்லாதபோது.
செர்மெலோ--ஃபிரேங்கல் கணக் கோட்பாடு~$ZF$ முரணற்றது
என்று கொண்டு, $ZF$ இன் நன்கு நிறுவப்படாத மாதிரிகள்
உள்ளன என்று காட்டுக; அதாவது,
$\dots x_2 \in x_1 \in x_0$ ஆக அமையும்
மாதிரிகள்~$\mathfrak{M}$ உள்ளன என்று காட்டுக.
\end{prob}

% SEGMENT OLP-0653-S05
\olref{thm:non-std} இல் கிடைத்த நிலையற்ற மாதிரி~$\Struct{M}$
நிலையான மாதிரியுடன் அடிப்படைச் சமானமானது. எனவே
$\Struct{M}$ இன் பல பண்புகளை வருவிக்கலாம்.
இந்தப் பிரிவின் எஞ்சிய பகுதி அந்தப் பணியை மேற்கொள்கிறது;
அதன் மூலம் $\Theory{N}$ இன் !!{enumerable} நிலையற்ற
மாதிரிகளுக்குத் துல்லியமான தன்மைப்படுத்தலைப் பெறுவோம்.

\begin{enumerate}
\item $\Domain{M}$ இன் எந்த உறுப்பும் தன்னைவிட
  $\prec$-சிறியது அல்ல: $\lforall[x][\lnot x < x]$
  என்ற வாக்கியம் $\Struct{N}$ இல் மெய்; ஆகவே
  $\Struct{M}$ இலும் மெய்.
\item அதே போன்ற காரணத்தால், $\prec$ ஆனது
  $\Domain{M}$ இன் \emph{நேரியல் வரிசை}; அதாவது
  $\Domain{M}$ மீதான முழுமையான, தற்சுட்டற்ற,
  கடப்புப் பண்புள்ள தொடர்பு.
\item $\mathbf{z}$ என்பது $\Domain{M}$ இன்
  $\prec$-மீச்சிறு உறுப்பு.
\item $\Domain{M}$ இன் ஒவ்வோர் உறுப்பும் அதன்
  $*$-அடுத்தெண்ணைவிட $\prec$-சிறியது; மேலும்
  $x$ ஐவிடப் பெரிய $\Domain{M}$ உறுப்புகளில்
  $\prec$-மீச்சிறியது $x^*$.

% SEGMENT OLP-0653-S06
\item $\Struct{M}$ இல்~$\Nat$ க்கு ஓரினமான
  $\prec$-தொடக்கப் பகுதி ஒன்று உள்ளது:
  $\mathbf{z}, \mathbf{z}^*, \mathbf{z}^{**}, \dots$.
  இதனை $\Domain{M}$ இன் \emph{நிலையான பகுதி}
  என்கிறோம். $\Domain{M}$ இன் மற்ற எந்த உறுப்பும்
  \emph{நிலையற்றது}. நிலையற்ற உறுப்புகள்
  $\Domain{M}$ இல் இருக்கவேண்டும்; இல்லையெனில்
  \olref{thm:non-std} இன் நிரூபிப்பிலுள்ள சார்பு~$h$
  ஓரினச்சார்பாகிவிடும். $\Struct{M}$ இன் இந்நிலையான
  பகுதியில் மதிப்பெடுக்கும் !!{variable}s ஆக
  $n, m, \dots$ ஐப் பயன்படுத்துகிறோம்.
\item ஒவ்வொரு நிலையற்ற உறுப்பும் எந்த நிலையான
  உறுப்பையும்விடப் பெரியது. ஏனெனில்
  ஒவ்வொரு $n \in \Nat$ க்கும்,
  \[
  \Sat{N}{\lforall[z][(\lnot(\eq[z][\Obj{0}] \lor
  \dots \lor \eq[z][\num{n}]) \lif \num{n} < z)]},
  \]
  எனவே $z \in \Domain{M}$ நிலையான உறுப்புகள்
  அனைத்திலிருந்தும் வேறுபட்டால், அவை அனைத்தையும்விட
  அது \emph{பெரியது}.
\item $\mathbf{z}$ தவிர $\Domain{M}$ இன் ஒவ்வோர்
  உறுப்பும் $\Domain{M}$ இன் ஒரேயொரு உறுப்பின்
  $*$-அடுத்தெண்; அந்த முன்னெண்ணை $^*x$
  எனக் குறிக்கிறோம். $x = y^*$ எனில், $x$, $y$
  ஆகியவற்றுள் ஒன்று நிலையானதாக இருந்தால் இரண்டும்
  நிலையானவை; ஒன்று நிலையற்றதாக இருந்தால் இரண்டும்
  நிலையற்றவை.

% SEGMENT OLP-0653-S07
\item $\Domain{M}$ மீதான சமானத் தொடர்பு~$\approx$
  ஐப் பின்வருமாறு வரையறுக்கிறோம்: ஏதோவொரு
  \emph{நிலையான}~$n$ க்கு
  $x \oplus n = y$ அல்லது $y \oplus n =x$
  எனில், எனில் மட்டுமே, $x \approx y$.
  வேறுவிதமாகச் சொன்னால், $x$, $y$ இரண்டும்
  முடிவுறு தொலைவில் இருந்தால், எனில் மட்டுமே,
  $x \approx y$. $n$, $m$ நிலையானவை எனில்
  $n \approx m$. $x$ இன் \emph{தொகுதி}
  என்பது அதன் சமான வகுப்பு
  $[x] = \Setabs{y}{x
  \approx y}$.

% SEGMENT OLP-0653-S08
\item $x \prec y$ மற்றும் $x \not\approx y$
  என்க. $\Sat{N}{\lforall[x][\lforall[y][(x < y \lif (x' < y \lor x' =
      y))]]}$ என்பதால், $x^* \prec y$ அல்லது
  $x^* = y$. இரண்டாவது வாய்ப்பு
  $x \approx y$ என்பதைக் கொடுப்பதால் இயலாது;
  % SOURCE CORRECTION TA-MOD-003: retain the strengthened successor inequality.
  ஆகவே $x^* \prec y$. அதேபோல்,
  $x \prec y$ மற்றும் $x \not\approx y$
  எனில், $x \prec {^*y}$. எனவே
  $x \prec y$ மற்றும் $x \not\approx y$
  எனில், $w \approx x$ ஆகும் ஒவ்வோர்
  உறுப்பும், $v \approx y$ ஆகும் ஒவ்வோர்
  உறுப்பைவிடவும் $\prec$-சிறியது.
  % SOURCE CORRECTION TA-MOD-007: only non-standard blocks extend in both directions.
  இதனால் ஒவ்வொரு நிலையற்ற தொகுதி~$[x]$ உம்
  இரு திசைகளிலும் முடிவிலாத சங்கிலியாக அமைகிறது:
  \[
  \cdots \prec  {^{**}x} \prec {^*}x \prec x \prec x^* \prec x^{**}
  \prec \cdots
  \]
  முழுக்களின் வரிசை வகையைக் கொண்டதால் இது
  $Z$-சங்கிலி எனப்படும்.

% SEGMENT OLP-0653-S09
% SOURCE CORRECTION TA-MOD-004: order only distinct blocks.
\item $\prec$ வரிசையை வெவ்வேறு தொகுதிகளுக்கும்
  உயர்த்தலாம்: $x \prec y$ ஆகவும் அவை
  வெவ்வேறு தொகுதிகளைச் சேர்ந்தவையாகவும் இருந்தால்,
  $x$ இன் தொகுதி $y$ இன் தொகுதியைவிடச் சிறியது.
  நிலையற்ற உறுப்பைக் கொண்ட தொகுதி
  \emph{நிலையற்ற} தொகுதி. நிலையான தொகுதியே
  மீச்சிறு தொகுதி.

% SEGMENT OLP-0653-S10
\item மீச்சிறு நிலையற்ற தொகுதி இல்லை:
  $y$ நிலையற்றது எனில், $x \prec y$,
  $x$ நிலையற்றது, $x \not\approx y$
  என அமையும் உறுப்பு ஒன்று உண்டு.
  நிரூபிப்பு: நிலையான மாதிரி~$\Struct{N}$ இல்
  ஒவ்வோர் எண்ணும் இரண்டால் வகுபடும்; மீதி ஒன்று
  இருக்கலாம். அதாவது
  % SOURCE CORRECTION TA-MOD-005: every y has some half, not every x is its half.
  $\Sat{N}{\lforall[y][\lexists[x][(y = x + x \lor y = x + x +
        \Obj{0}')]]}$.
  அடிப்படைச் சமானத்தால், ஒவ்வொரு
  $y \in \Domain{M}$ க்கும்
  $x \oplus x = y$ அல்லது
  $x \oplus x \oplus \mathbf{z}^*= y$
  ஆக அமையும் $x \in \Domain{M}$ ஒன்று உண்டு.
  $x$ நிலையானதாக இருந்தால் $y$ உம்
  நிலையானதாகிவிடும்; ஆகவே $x$ நிலையற்றது.
  நிலையற்ற உறுப்பு பூச்சியைவிடப் பெரியது
  என்பதால், அது அதன் இரட்டிப்பைவிடச் சிறியது.
  மேலும் $x$, $y$ வெவ்வேறு தொகுதிகளைச்
  சேர்ந்தவை; அதாவது $x \not\approx y$.
  அப்படியல்லாமல், ஏதோவொரு நிலையான~$n$
  க்கு $x \oplus n = y$ என்க.
  $y = x \oplus x$ எனில்,
  $x \oplus n = x \oplus x$;
  கூட்டலின் நீக்கல் விதியால்
  ($\Struct{N}$ இலும், ஆகவே $\Struct{M}$
  இலும் அது பொருந்தும்) $x = n$ கிடைக்கும்.
  அப்போது $x$ நிலையானதாகிவிடும்.
  $y = x \oplus x \oplus \mathbf{z}^*$
  எனும்போதும் இதே போன்றது.
% SEGMENT OLP-0653-S11
\item இதே போன்ற வாதத்தால், மீப்பெரு தொகுதியும் இல்லை.
% SEGMENT OLP-0653-S12
\item தொகுதிகளின் வரிசை செறிந்தது:
  $[x]$ ஆனது~$[y]$ ஐவிடச் சிறிதாகவும்
  ($x \not\approx y$), அவை வேறுபட்டவையாகவும்
  இருந்தால், அவற்றுக்கு இடையில் இரண்டிலிருந்தும்
  வேறுபட்ட தொகுதி~$[z]$ ஒன்று உண்டு.
  $x \prec y$ என்க. முன்போல,
  $x \oplus y$ இரண்டால் வகுபடும்
  (மீதி இருக்கலாம்). எனவே
  $x \oplus y = u \oplus u$ அல்லது
  $x \oplus y = u \oplus u \oplus \mathbf{z}^*$
  ஆக அமையும் $u
  \in \Domain{M}$ ஒன்று உண்டு.
  $u$ என்பது $x$, $y$ ஆகியவற்றின்
  சராசரி; ஆகவே அவற்றுக்கிடையில் உள்ளது.
  $x \oplus y = u \oplus u$ எனக் கொள்க
  (மற்ற நேர்வும் இதே போன்றது).
  $u \approx x$ எனில், ஏதோவொரு
  நிலையான~$n$ க்கு
  \[
  x \oplus y = x \oplus n \oplus x \oplus n,
  \]
  ஆகவே $y = x \oplus n \oplus n$;
  இது கருதுகோளுக்கு மாறாக $x \approx y$
  என்பதைக் கொடுக்கும். எனவே
  $u \not\approx x$. இதே போன்ற வாதம்
  $u \not\approx y$ என்பதையும் தருகிறது.
\end{enumerate}
% SEGMENT OLP-0653-S13
எனவே நிலையற்ற தொகுதிகள் விகிதமுறு எண்களைப்
போல வரிசைப்படுத்தப்பட்டுள்ளன: அவை முனைகளற்ற
!!a{enumerable} செறிந்த நேரியல் வரிசையாக அமைகின்றன.
இதனால் மெய் எண்கணிதத்தின் எந்த இரு
!!{enumerable} நிலையற்ற மாதிரிகள்
$\Struct{M}_1$, $\Struct{M_2}$ ஆகியவற்றின்
$<$ மற்றும் $=$ ஐ மட்டுமே கொண்ட மொழிக்கான
குறைப்புகள் ஓரினமானவை என்று பின்வருகிறது.
உண்மையில் ஓரினச்சார்பு~$h$ ஐப்
பின்வருமாறு வரையறுக்கலாம்:
$\Struct{M_1}$, $\Struct{M_2}$
இன் நிலையான பகுதிகள் நிலையான மாதிரி
$\Struct{N}$ க்கு ஓரினமானவை;
எனவே ஒன்றுக்கொன்றும் ஓரினமானவை.
நிலையற்ற பகுதியை உருவாக்கும் தொகுதிகள்
விகிதமுறு எண்களைப் போலவே வரிசைப்படுத்தப்பட்டுள்ளன;
எனவே \olref[dlo]{thm:cantorQ} படி
அவையும் ஓரினமானவை. தொகுதிகளுக்கிடையிலான
ஓரினச்சார்பை, ஒவ்வொரு தொகுதியிலும்
தன்னிச்சையான உறுப்புகளைப் பொருத்தி,
பின்னர் $\Struct{M_1}$ இல்~$x$
இன் அடுத்தெண்ணின் படிமம்,
$\Struct{M_2}$ இல்~$x$ இன் படிமத்தின்
அடுத்தெண்ணாக இருக்குமாறு அமைத்து,
தொகுதிகளுக்கு \emph{உள்ளேயும்}
ஓரினச்சார்பாக நீட்டிக்கலாம்.
இதனால் $\mathfrak{M}_1$,
$\mathfrak{M}_2$ ஆகியவை எண்கணிதத்தின்
முழு மொழியில் ஓரினமானவை என்று
\emph{பின்வராது} என்பதைக் கவனிக்க:
ஓரினத்தன்மை எப்போதும் ஒரு
குறியீட்டமைப்பைச் சார்ந்தது;
$\Domain{M_1}$, $\Domain{M_2}$
மீதான கூட்டலையும் பெருக்கலையும்
ஓரினமற்ற வழிகளில் வரையறுக்கலாம்.
(ஒரு நிறைவான கோட்பாட்டின்
எண்ணத்தக்க மாதிரிகளின் எண்ணிக்கை~2
ஆக இருக்க முடியாது என்ற வாட்டின்
புகழ்பெற்ற தேற்றத்திலிருந்தும்
இது பின்வருகிறது.)

% SEGMENT OLP-0653-S14
\begin{prob}
எண்கணிதத்தின் நிலையற்ற மாதிரியில்
மீப்பெரு தொகுதி இருக்க முடியாது என்று காட்டுக.
\end{prob}

% SEGMENT OLP-0653-S15
\begin{prob}
$\Lang{L}$ என்பது $\eq$ தவிர
$<$ ஐ மட்டுமே !!{predicate}
ஆகக் கொண்ட முதல்தர !!{language}
ஆகட்டும்; $\Struct{N} = (\Nat,
<)$ என்க. $\Struct{N}$ இன்
முடிவுறும் அல்லது நிரப்பி முடிவுறும்
உட்கணங்கள் அனைத்தும் வரையறுக்கத்தக்கவை.
$\Struct{N}$ இன் வரையறுக்கத்தக்க
உட்கணங்கள் இவை \emph{மட்டுமே}
என்று காட்டுக.

(குறிப்பு: முதலில்
$\Obj{prc}(x,y)$ என்பது
“$x$ ஆனது $y$ இன் உடனடி முன்னெண்”
என்பதைச் சுருக்கி எழுதும்
$\Lang{L}$-வாய்பாடு என்க:
\[
x<y \land \lnot \lexists[z][(x<z \land z < y)].
\]
இப்போது $\Struct{N}$ இன் ஒவ்வொரு
வரையறுக்கத்தக்க உட்கணத்துக்கும்
$\Lang{L}$ இல் வாய்பாடு~$!A(x)$
ஒன்று உண்டு. அத்தகைய ஒவ்வொரு
$!A$ க்கும் வாக்கியம்~$!D$ ஐக் கருதுக:
\[
\lexists[x][\lforall[y][\lforall[z][((x<y \land x<z \land \Obj{prc}(y,z)
  \land !A(y)) \lif !A(z))]]].
\]
$!A$ வரையறுக்கும் $\Struct{N}$
இன் உட்கணம் முடிவுறுவதோ அல்லது
அதன் நிரப்பி முடிவுறுவதோ,
$\Sat{N}{!D}$ எனும்போது,
அப்போது மட்டுமே என்று காட்டுக.

இப்போது $\Struct{N}$ உடன்
அடிப்படைச் சமானமான நிலையற்ற
மாதிரி~$\Struct{M}$ ஐ எடுத்துக்கொள்க.
$a \in \Domain{M}$ நிலையற்றது எனில்,
$a$ ஐவிடப் பெரிய
$b, c \in \Domain{M}$ ஐ எடுத்துக்கொள்க;
$c$ இன் உடனடி முன்னெண்~$b$
ஆகட்டும். அப்போது $h(b)=c$
% SOURCE CORRECTION TA-MOD-006: h is an automorphism of the order reduct.
ஆக அமையும் $\Domain{M}$ இன்
வரிசைச் சுயஓரினச்சார்பு~$h$
ஒன்று உண்டு (ஏன்?).
எனவே $b$ ஆனது~$!A$ ஐ
நிறைவுறுத்தினால் $c$ உம்
நிறைவுறுத்தும் (ஏன்?).
இதனால் $!D$ ஆனது
$\Struct{M}$ இல் மெய்;
ஆகவே $\Struct{N}$ இலும் மெய்.
இதன்படி $!A$ வரையறுக்கும்
$\Struct{N}$ இன் உட்கணம்
முடிவுறுவதோ அல்லது
அதன் நிரப்பி முடிவுறுவதோ
ஆகும்.
\end{prob}


% END SOURCE UNIT OLP-0653

\clearpage
\noindent மூல அலகு: \texttt{OLP-0654}\par

% BEGIN SOURCE UNIT OLP-0654 content/normal-modal-logic/syntax-and-semantics/normal-modal-logics.tex
% Exact modular target bytes: 5967; SHA256: af5e915cfbaa62b774a52ba31b87a6ab2ad105db5729e6b806ef74f7331a3f91
% SEGMENT OLP-0654-S01


\olfileid{nml}{syn}{nor}

\olsection{இயல்பான வாய்ப்புநிலைத் தருக்கங்கள்}

தகுந்த ஒழுங்குப் பண்புகள் உள்ள !!{formula}s இன் எந்தக் கணத்தையும்
வாய்ப்புநிலைத் தருக்கமாகக் கருதலாம். சில வாய்ப்புநிலைத் தருக்கங்கள்
அவற்றின் சுவாரஸ்யமான பொருள்கோள்களாலும் பயன்பாடுகளாலும்,
முக்கியமாக அவற்றைப் பற்றிய கோட்பாட்டு அறிவு நன்கு வளர்ந்திருப்பதாலும்,
சிறப்பாகக் கவனிக்கப்படுகின்றன. இவையே இயல்பான வாய்ப்புநிலைத்
தருக்கங்கள் எனப்படுகின்றன.

% SEGMENT OLP-0654-S02
\begin{defn}\ollabel{defn:modal-logic}
\emph{வாய்ப்புநிலைத் தருக்கம்}~$\Log L$ என்பது அடிப்படை
வாய்ப்புநிலை மொழியின் !!{formula}s இன் கணம்; அது
\begin{enumerate}
\item கூற்றுத் தருக்க மெய்மங்கள் அனைத்தையும் கொண்டிருக்கும்;
\item சீரான பதிலீட்டின் கீழ் மூடப்பட்டிருக்கும்; மேலும்
\item மோடஸ் போனென்ஸின் கீழ் மூடப்பட்டிருக்கும்; அதாவது
  $!A$ மற்றும் $(!A \lif
  !B) \in \Log L$ எனில், $!B \in \Log L$ உம் பொருந்தும்.
\end{enumerate}
\end{defn}

% SEGMENT OLP-0654-S03
\begin{ex}
இவ்வரையறையை நிறைவுசெய்வதால் வாய்ப்புநிலைத் தருக்கங்களாகக்
கணிக்கப்படும், பெரிதாக ஆர்வமூட்டாத !!{formula}s இன்
கணங்களுக்குச் சில எடுத்துக்காட்டுகள்:
\begin{enumerate}
\item மெய்மங்கள் அனைத்தின் கணம்;
\item எல்லா !!{formula}s இன் கணம் (முரணுடைய வாய்ப்புநிலைத் தருக்கம்).
\end{enumerate}
\end{ex}

% SEGMENT OLP-0654-S04
\begin{defn}\label{defn:normal-modal-logic}
வாய்ப்புநிலைத் தருக்கம்~$\Log L$ \emph{இயல்பானது}
எனில், எனில் மட்டுமே,
\begin{enumerate}
\item அது பின்வரும் இரண்டையும் கொண்டிருக்கும்:
\begin{align}
\tag{K} \Box(p \lif q) & \lif (\Box p \lif \Box q) \\
\tag{Dual} \Diamond p & \liff \lnot \Box \lnot p
\end{align}
\item மேலும் அது இன்றியமையாக்கலின் கீழ் மூடப்பட்டிருக்கும்;
  அதாவது $!A \in
  \Log L$ எனும்போதெல்லாம் $ \Box !A \in \Log L$ உம் பொருந்தும்.
\end{enumerate}
\end{defn}

% SEGMENT OLP-0654-S05
வாய்ப்புநிலைத் தருக்கங்களை வரையறுக்கும் இரண்டு முதன்மை வழிகளை
நாம் காண்போம். ஒன்று நிறுவல்-கோட்பாட்டு வழி: குறிப்பிட்ட
!!{derivation} முறைமைகளில் வருவிக்கத்தக்க !!{formula}s இன்
கணங்களாக அவற்றைக் கொள்வது. மற்றொன்று மாதிரிக் கோட்பாட்டு வழி:
குறிப்பிட்ட மாதிரி வகுப்புகளில் செல்லுபடியாகும் !!{formula}s
இன் கணங்களாக அவற்றைக் கொள்வது. இது சாதாரண கூற்றுத்
தருக்கத்திலும் (அதேபோல் முதல்தரத் தருக்கத்திலும்) நிலவும்
நிலையைப் பிரதிபலிக்கிறது. அங்கும் ஓர் !!{formula}s இன்
கணத்தை இரு வேறு வழிகளில் குறிப்பிடலாம்: எடுத்துக்காட்டாக,
எல்லா மெய்மதிப்பு ஒதுக்கீடுகளிலும் மெய்யான மெய்மங்களின்
கணமாகவும், ஏதோவொரு !!{derivation} முறைமையில்
!!{derivable} ஆகும் எல்லா !!{formula}s இன்
கணமாகவும் குறிப்பிடலாம். இயல்பான வாய்ப்புநிலைத் தருக்கங்களுக்கு,
தொடர்புசார் மாதிரிகளையும் அடிகோள்சார் (ஹில்பர்ட் வகை)
நிறுவல் முறைமைகளையும் பயன்படுத்தி இந்த இருவழி
குறிப்பீட்டை அமைக்கலாம்.

% END SOURCE UNIT OLP-0654


\chapter{நிறுவல் கோட்பாட்டின் கூடுதல் பகுதிகள்}
\clearpage
\section{வெட்டு நீக்கலின் தனித்த மூலத் துணுக்கு}
\noindent மூல அலகு: \texttt{OLP-0660}.
இது முழு அத்தியாயமாக அமைக்கப்படாத தனித்த நிறுவல் துணுக்கு.
கோப்பின் பெயர் \texttt{intuitionistic.tex} என்றாலும், அதன்
வாய்பாடுகளில் பலமுடிவு \Log{G3c} தொடரணிகளே பயன்படுத்தப்படுகின்றன.
பெயரை வைத்து நிறுவல் முறைமை மாற்றப்படவில்லை.
\begingroup
\olfileid{pt}{cut}{int}

% BEGIN SOURCE UNIT OLP-0660 content/proof-theory/cut-elimination/intuitionistic.tex
% Exact modular target bytes: 3319; SHA256: e6822abbd7974e758d369eefd1f3fff21e487f58c13a0e7d6aa45fa2a46f25e6
% SEGMENT OLP-0660-S01
% Part: proof-theory
% Chapter: cut-elimination
% Auxiliary fragment: not imported into the reader

$!A \ident (!B \lif !C)$ எனில், $\pi$ பின்வருமாறு முடிகிறது:
\[
\AxiomC{}
\RightLabel{$\pi_1'$}
\Deduce$!B, \Gamma \fCenter \Delta, !C$
\RightLabel{\RightR{\lif}}
\UnaryInf$\Gamma \fCenter \Delta, !B \lif !C$
\LeftSubproofLabel{$\pi_1$}
\AxiomC{}
\RightLabel{$\pi_2'$}
\Deduce$!B \lif !C, \Pi \fCenter \Lambda, !B$
\AxiomC{}
\RightLabel{$\pi_2''$}
\Deduce$!C, \Pi \fCenter \Lambda$
\RightLabel{\LeftR{\lif}}
\BinaryInf$!B \lif !C, \Pi \fCenter \Lambda$
\RightSubproofLabel{$\pi_2$}
\RightLabel{\Cut}
\BinaryInf$\Gamma, \Pi \fCenter \Delta, \Lambda$
\DisplayProof
\]

% SEGMENT OLP-0660-S02
\Cut{} இல் முடியும் பின்வரும் !!{proof} ஐக் கருதுக:
\[
\AxiomC{}
\RightLabel{$\pi_1'$}
\Deduce$!B, \Gamma \fCenter \Delta, !C$
\RightLabel{\RightR{\lif}}
\UnaryInf$\Gamma \fCenter \Delta, !B \lif !C$
\LeftSubproofLabel{$\pi_1$}
\AxiomC{}
\RightLabel{$\pi_2'$}
\Deduce$!B \lif !C, \Pi \fCenter \Lambda, !B$
\RightLabel{\Cut}
\BinaryInf$\Gamma, \Pi \fCenter \Delta, \Lambda, !B$
\DisplayProof
\]
இதன் வெட்டு !!{height}, $\pi$ இன் வெட்டு உயரத்தைவிடக் குறைவு.
தொகுத்தறிதல் கருதுகோள் $\Gamma, \Pi \fCenter \Delta,
\Lambda, !B$ இன் !!a{proof}~$\pi_1''$ ஐத் தருகிறது.

% SEGMENT OLP-0660-S03
\Cut{} இல் முடியும் பின்வரும் !!{proof} ஐக் கருதுக:
\[
\AxiomC{}
\RightLabel{$\pi_1'$}
\Deduce$!B, \Gamma \fCenter \Delta, !C$
\AxiomC{}
\RightLabel{$\pi_2''$}
\Deduce$!C, \Pi \fCenter \Lambda$
\RightLabel{\Cut}
\BinaryInf$!B, \Gamma, \Pi \fCenter \Delta, \Lambda$
\DisplayProof
\]
இதன் வெட்டுத் தரமும் வெட்டு !!{height} உம் $\pi$ இன்
வெட்டுத் தரம், வெட்டு !!{height} ஆகிய இரண்டையும்விடக் குறைவு.
எனவே தொகுத்தறிதல் கருதுகோள் $!B, \Gamma, \Pi
\fCenter \Delta, \Lambda$ இன் \Log{G3c}-!!{proof}~$\pi_2'''$
ஐத் தருகிறது.

% SEGMENT OLP-0660-S04
இப்போது \Cut{} இல் முடியும் பின்வரும் !!{proof} க்குத்
தொகுத்தறிதல் கருதுகோளைப் பயன்படுத்துக:
\[
\AxiomC{}
\RightLabel{$\pi_1''$}
\Deduce$\Gamma, \Pi \fCenter \Delta, \Lambda, !B$
\AxiomC{}
\RightLabel{$\pi_2'''$}
\Deduce$!B, \Gamma, \Pi \fCenter \Delta, \Lambda$
\RightLabel{\Cut}
\BinaryInf$\Gamma, \Gamma, \Pi, \Pi \fCenter \Delta, \Delta, \Lambda, \Lambda$
\DisplayProof
\]
(இதன் வெட்டுத் தரம் $< \cutr{\pi}$.) ஆகவே
$\Gamma, \Gamma, \Pi, \Pi \fCenter \Delta, \Delta,
\Lambda, \Lambda$ இன் !!a{proof} கிடைக்கிறது.
\olref[seq][inv]{prop:G3c-cont-adm} ஐப் பயன்படுத்திச் சுருக்கினால்,
$\Gamma, \Pi \fCenter \Delta, \Lambda$ இன்
!!{proof}~$\pi'$ கிடைக்கிறது.

% END SOURCE UNIT OLP-0660

\endgroup
\clearpage
\noindent மூல அலகு: \texttt{OLP-0694}\par

% BEGIN SOURCE UNIT OLP-0694 content/proof-theory/propositions-as-types/sequent-natural-deduction.tex
% Exact modular target bytes: 6598; SHA256: 84ef6aa6702713458859311804d63e0738a850a4c4c0eebad0cd446664cc3641
% SEGMENT OLP-0694-S01


\olfileid{int}{pty}{snd}

\olsection{தொடரணிகளுடன் இயல்பான வருவித்தல்}

$\Gamma \Sequent !A$ என்று எழுதுவோம்: அதாவது,
$\Gamma$ என்பது !!{undischarged} எடுகோள்களின்
கணமாகவும் $!A$ முடிவாகவும் உள்ள இயல்பான வருவித்தல்
!!{derivation} ஒன்று உள்ளது. $\Gamma$ காலிக் கணம்
எனில், $\Sequent !A$ என்று எழுதுவோம்.

$\Gamma \cup \{!A_1, \dots, !A_n\}$ என்பதற்கு
$\Gamma, !A_1, \dots, !A_n$ என்றும்,
$\Gamma \cup \Delta$ என்பதற்கு $\Gamma, \Delta$
என்றும் எழுதுவோம்.

% SEGMENT OLP-0694-S02
$\Gamma \Sequent !A \land !B$ எனில்,
$\Gamma$ என்பதை !!{undischarged} எடுகோள்களாகவும்
$!A \land !B$ ஐ இறுதி-!!{formula} ஆகவும் கொண்ட
!!a{derivation} நமக்கு உள்ளது. அதன் அடியில்
$\Elim{\land}$ விதியைப் பயன்படுத்தி, அதே
!!{undischarged} எடுகோள்களிலிருந்து $!A$ ஐ முடிவாகக்
கொண்ட !!a{derivation} ஐப் பெறலாம். வேறு சொற்களில்,
$\Gamma \Sequent !A \land !B$ எனில்
$\Gamma \Sequent !A$:
\begin{prooftree}
  \Axiom$\Gamma \fCenter !A \land !B$
  \RightLabel{$\Elim{\land}$}
  \UnaryInf$\Gamma \fCenter !A$
  \DisplayProof\qquad\bottomAlignProof
  \Axiom$\Gamma \fCenter !A \land !B$
  \RightLabel{$\Elim{\land}$}
  \UnaryInf$\Gamma \fCenter !B$
\end{prooftree}
$\Elim{\land}$ என்ற அடையாளம், இயல்பான வருவித்தலில்
அதே பெயருடைய விதியுடனான தொடர்பைக் காட்டுகிறது.

% SEGMENT OLP-0694-S03
அதேபோல், $\Gamma, !A \Sequent !B$ எனக் கொள்வோம்.
அதாவது, $\Gamma, !A$ என்பவற்றை !!{undischarged}
எடுகோள்களாகவும் $!B$ ஐ இறுதி-!!{formula} ஆகவும்
கொண்ட !!a{derivation} உள்ளது. இதற்கு
$\Intro{\lif}$ விதியைப் பயன்படுத்தினால்,
$\Gamma$ என்பதை !!{undischarged} எடுகோள்களாகவும்
$!A \lif !B$ ஐ இறுதி-!!{formula} ஆகவும் கொண்ட
!!a{derivation} கிடைக்கும்; அதாவது,
$\Gamma \Sequent !A \lif !B$.
எடுகோள்களின் !!{discharge} இங்கு மேலும் தெளிவாகக்
காட்டப்படுவதைக் கவனிக்கவும்.
\begin{prooftree}
  \Axiom $\Gamma, !A \fCenter !B$
  \RightLabel{$\Intro{\lif}$}
  \UnaryInf $\Gamma \fCenter !A \lif !B$
\end{prooftree}

இதே முறையில் மற்ற விதிகளிலிருந்தும் முடிவுகளைக்
காணலாம்; அவை முழுவதும் பின்வருமாறு:
\begin{gather*}
  \Axiom $\Gamma \fCenter !A$
  \Axiom $\Delta \fCenter !B$
  \RightLabel{$\Intro{\land}$}
  \BinaryInf $\Gamma,\Delta \fCenter !A \land !B$
  \DisplayProof\\
  \Axiom $\Gamma \fCenter !A \land !B$
  \RightLabel{$\Elim{\land}_1$}
  \UnaryInf $\Gamma \fCenter !A$
  \DisplayProof
  \qquad
  \Axiom $\Gamma \fCenter !A \land !B$
  \RightLabel{$\Elim{\land}_2$}
  \UnaryInf $\Gamma \fCenter !B$
  \DisplayProof
  \\
  \Axiom $\Gamma \fCenter !A$
  \RightLabel{$\Intro{\lor}_1$}
  \UnaryInf $\Gamma \fCenter !A \lor !B$
  \DisplayProof\qquad
  \Axiom $\Gamma \fCenter !B$
  \RightLabel{$\Intro{\lor}_2$}
  \UnaryInf $\Gamma \fCenter !A \lor !B$
  \DisplayProof\\
  \Axiom $\Gamma \fCenter !A \lor !B$
  \Axiom $\Delta, !A \fCenter !C$
  \Axiom $\Delta', !B \fCenter !C$
  \RightLabel{$\Elim{\lor}$}
  \TrinaryInf $\Gamma, \Delta, \Delta' \fCenter !C$
  \DisplayProof
  \\
  \Axiom $\Gamma, !A \fCenter !B$
  \RightLabel{$\Intro{\lif}$}
  \UnaryInf $\Gamma \fCenter !A \lif !B$
  \DisplayProof
  \qquad
  \Axiom $\Delta \fCenter !A \lif !B$
  \Axiom $\Gamma \fCenter !A$
  \RightLabel{$\Elim{\lif}$}
  \BinaryInf $\Gamma, \Delta \fCenter !B$
  \DisplayProof
  \\
  \Axiom $\Gamma \fCenter \lfalse$
  \RightLabel{$\FalseInt$}
  \UnaryInf $\Gamma \fCenter !A$
  \DisplayProof
\end{gather*}

% SEGMENT OLP-0694-S04
ஓர் எடுகோள் தனியே $!A$ இலிருந்து $!A$ இன்
!!a{derivation} ஆகும்; எனவே $!A \Sequent !A$
எப்போதும் கிடைக்கும்.

\begin{prooftree}
  \AxiomC{$ $}
  \UnaryInfC{$!A \fCenter !A$}
\end{prooftree}

இவ்விதிகளை ஒருங்கே எடுத்தால், எந்த இயல்பான வருவித்தல்
!!{derivation}s உள்ளன என்பதை விவரிக்கும் கலனமாகக்
கொள்ளலாம். இயல்பான வருவித்தலின் குறியீட்டு மாற்றுருவாகவும்
கொள்ளலாம்: ஒவ்வொரு படியும் !!{derive} செய்யப்பட்ட
!!{formula} ஐ மட்டுமன்றி, அது எந்த !!{undischarged}
எடுகோள்களிலிருந்து !!{derive} செய்யப்பட்டது என்பதையும்
பதிவுசெய்கிறது.

எடுத்துக்காட்டாக, பின்வரும் நிறுவலைக் காண்க:
\begin{prooftree}
  \Axiom$!A \fCenter !A$
  \UnaryInf $!A \fCenter !A \lor (!A \lif \lfalse)$
  \Axiom $!B \fCenter !B$
  \BinaryInf$!A, !B \fCenter \lfalse$
  \UnaryInf$!B \fCenter !A \lif \lfalse$
  \UnaryInf$!B \fCenter !A \lor (!A \lif \lfalse)$
  \Axiom$!B \fCenter !B$
  \BinaryInf$!B \fCenter \lfalse$
  \UnaryInf$\fCenter !B \lif \lfalse$
\end{prooftree}
இங்கு $!B$ என்பது
$(!A \lor (!A \lif \lfalse))\lif \lfalse$
என்பதன் சுருக்கம்.


% END SOURCE UNIT OLP-0694


\clearpage
\section{மாற்று தொடர்கணித விதி அட்டவணைகள்}
\noindent மூல அலகுகள்:
\texttt{OLP-0705}, \texttt{OLP-0708}, \texttt{OLP-0709}, \texttt{OLP-0710}.
மூலத்தில் வழங்கப்பட்ட விதிகளும் நிபந்தனைகளும் அப்படியே பாதுகாக்கப்படுகின்றன.
\begingroup
\olfileid{pt}{seq}{rai}

% BEGIN SOURCE UNIT OLP-0705 content/proof-theory/sequent-calculus/rules-G1i.tex
% Exact modular target bytes: 3787; SHA256: 6a52040593872550089ff34972ffe1e78a694fb7c06cb3742d254cdd694fac4f
% SEGMENT OLP-0705-S01


\begin{table}
\begin{tabular}{@{}c@{}c@{}}
\multicolumn{2}{@{}c@{}}{அடிக்கோள்கள்}\\[2ex]
$!A \Sequent !A$ & $\lfalse \Sequent \quad$
\\[2ex]
% SEGMENT OLP-0705-S02
\multicolumn{2}{@{}c@{}}{கட்டமைப்பு விதிகள்}\\[2ex]
\Axiom$ \Gamma \fCenter \Delta $
\RightLabel{\LeftR{\Weakening}}
\UnaryInf$ !A, \Gamma \fCenter \Delta$
\DisplayProof
&
\Axiom$ \Gamma \fCenter $
\RightLabel{\RightR{\Weakening}}
\UnaryInf$ \Gamma \fCenter !A$
\DisplayProof
\\[4ex]
\Axiom$ !A, !A, \Gamma \fCenter \Delta $
\RightLabel{\LeftR{\Contraction}}
\UnaryInf$ !A, \Gamma \fCenter \Delta$
\DisplayProof
&
\\[4ex]
% SEGMENT OLP-0705-S03
\multicolumn{2}{@{}c@{}}{தருக்க விதிகள்}\\[2ex]
\Axiom$ \Gamma \fCenter !A $
\RightLabel{\LeftR{\lnot}}
\UnaryInf$ \lnot !A, \Gamma \fCenter $
\DisplayProof
&
\Axiom$!A, \Gamma \fCenter $
\RightLabel{\RightR{\lnot}}
\UnaryInf$ \Gamma \fCenter \lnot !A $
\DisplayProof
\\[4ex]\begin{tabular}[t]{@{}l@{}}
\Axiom$ !A, \Gamma \fCenter \Delta$
\RightLabel{\LeftR{\land}}
\UnaryInf$ !A \land !B, \Gamma \fCenter \Delta$
\DisplayProof
\\[3ex]
\Axiom$!B, \Gamma \fCenter \Delta$
\RightLabel{\LeftR{\land}}
\UnaryInf$!A \land !B, \Gamma \fCenter \Delta$
\DisplayProof\\[3ex]
\end{tabular}
&
\Axiom$\Gamma \fCenter !A$
\Axiom$ \Gamma \fCenter !B$
\RightLabel{\RightR{\land}}
\BinaryInf$ \Gamma \fCenter !A \land !B $
\DisplayProof
\\[4ex]\Axiom$!A, \Gamma \fCenter \Delta$
\Axiom$!B, \Gamma \fCenter \Delta$
\RightLabel{\LeftR{\lor}}
\BinaryInf$!A \lor !B, \Gamma \fCenter \Delta$
\DisplayProof
&
\begin{tabular}[t]{@{}r@{}}
\Axiom$\Gamma \fCenter !A$
\RightLabel{\RightR{\lor}}
\UnaryInf$ \Gamma \fCenter !A \lor !B$
\DisplayProof
\\[3ex]
\Axiom$ \Gamma \fCenter !B$
\RightLabel{\RightR{\lor}}
\UnaryInf$ \Gamma \fCenter !A \lor !B$
\DisplayProof\\[3ex]
\end{tabular}
\\[4ex]
\Axiom$ \Gamma \fCenter !A$
\Axiom$ !B, \Gamma \fCenter \Delta$
\RightLabel{\LeftR{\lif}}
\BinaryInf$ !A \lif !B, \Gamma \fCenter \Delta$
\DisplayProof
&
\Axiom$ !A, \Gamma \fCenter !B$
\RightLabel{\RightR{\lif}}
\UnaryInf$ \Gamma \fCenter !A \lif !B $
\DisplayProof
\\[4ex]
\Axiom$ !A(t), \Gamma \fCenter \Delta$
\RightLabel{\LeftR{\lforall}}
\UnaryInf$ \lforall[x][!A(x)],\Gamma \fCenter \Delta$
\DisplayProof
&
\Axiom$ \Gamma \fCenter !A(a) $
\RightLabel{\RightR{\lforall}}
\UnaryInf$ \Gamma \fCenter \lforall[x][!A(x)]$
\DisplayProof
\\[4ex]
\Axiom$ !A(a), \Gamma \fCenter \Delta $
\RightLabel{\LeftR{\lexists}}
\UnaryInf$ \lexists[x][!A(x)], \Gamma \fCenter \Delta$
\DisplayProof
&
\Axiom$ \Gamma \fCenter !A(t) $
\RightLabel{\RightR{\lexists}}
\UnaryInf$ \Gamma \fCenter \lexists[x][!A(x)]$
\DisplayProof
\end{tabular}
\caption{உள்ளுணர்வுவாதத் தருக்கத்திற்கான \Log{G1i} விதிகள்.
$\RightR{\forall}$ மற்றும் $\LeftR{\exists}$ ஆகிய விதிகளில்,
!!{constant}~$a$ முடிவுத் தொடரணியில் இடம்பெறக் கூடாது.
$\Gamma$ என்பது !!{formula}s இன் பன்மைக் கணம்;
$\Delta$ வெறுமையாகவோ ஒரே ஓர் !!{formula} மட்டும் கொண்டதாகவோ
இருக்கலாம். \RightR{\Contraction} விதி இல்லை.
\Log{G1m} என்பது \RightR{\Weakening} விதியும்
$\lfalse, \Gamma \Sequent \Delta$ என்ற அடிக்கோள்களும்
நீக்கப்பட்ட \Log{G1i} ஆகும்.}
\ollabel{tab:G1i}
\end{table}


% END SOURCE UNIT OLP-0705

\clearpage

% BEGIN SOURCE UNIT OLP-0708 content/proof-theory/sequent-calculus/rules-G3i.tex
% Exact modular target bytes: 3293; SHA256: e7d819fc86f53ea95a692d5c5eca5f403d8739c448316a4f74ac8dca3d81602f
% SEGMENT OLP-0708-S01


\begin{table}
\begin{tabular}{@{}c@{}c@{}}
\multicolumn{2}{@{}c@{}}{அடிக்கோள்கள்}\\[2ex]
$!C, \Gamma \Sequent !C$ & $\lfalse, \Gamma \Sequent \Delta$
\\[2ex]
% SEGMENT OLP-0708-S02
\multicolumn{2}{@{}c@{}}{தருக்க விதிகள்}\\[2ex]
\Axiom$\lnot !A, \Gamma \fCenter !A $
\RightLabel{\LeftR{\lnot}}
\UnaryInf$\lnot !A, \Gamma \fCenter \Delta$
\DisplayProof
&
\Axiom$!A, \Gamma \fCenter$
\RightLabel{\RightR{\lnot}}
\UnaryInf$ \Gamma \fCenter \lnot !A $
\DisplayProof
\\[4ex]
\Axiom$ !A, !B, \Gamma \fCenter \Delta$
\RightLabel{\LeftR{\land}}
\UnaryInf$ !A \land !B, \Gamma \fCenter \Delta$
\DisplayProof
&
\Axiom$\Gamma \fCenter !A$
\Axiom$ \Gamma \fCenter !B$
\RightLabel{\RightR{\land}}
\BinaryInf$ \Gamma \fCenter !A \land !B $
\DisplayProof
\\[4ex]\Axiom$!A, \Gamma \fCenter \Delta$
\Axiom$!B, \Gamma \fCenter \Delta$
\RightLabel{\LeftR{\lor}}
\BinaryInf$!A \lor !B, \Gamma \fCenter \Delta$
\DisplayProof
&
\begin{tabular}[t]{@{}r@{}}
\Axiom$\Gamma \fCenter !A$
\RightLabel{\RightR{\lor}}
\UnaryInf$ \Gamma \fCenter !A \lor !B$
\DisplayProof
\\[3ex]
\Axiom$\Gamma \fCenter !B$
\RightLabel{\RightR{\lor}}
\UnaryInf$ \Gamma \fCenter !A \lor !B$
\DisplayProof\\[3ex]
\end{tabular}
\\[4ex]
\Axiom$!A \lif !B, \Gamma \fCenter !A$
\Axiom$!B, \Gamma \fCenter \Delta$
\RightLabel{\LeftR{\lif}}
\BinaryInf$!A \lif !B, \Gamma \fCenter \Delta$
\DisplayProof
&
\Axiom$!A, \Gamma \fCenter !B$
\RightLabel{\RightR{\lif}}
\UnaryInf$ \Gamma \fCenter !A \lif !B $
\DisplayProof
\\[4ex]
\Axiom$ !A(t), \lforall[x][!A(x)],\Gamma \fCenter \Delta$
\RightLabel{\LeftR{\lforall}}
\UnaryInf$ \lforall[x][!A(x)],\Gamma \fCenter \Delta$
\DisplayProof
&
\Axiom$ \Gamma \fCenter !A(a) $
\RightLabel{\RightR{\lforall}}
\UnaryInf$ \Gamma \fCenter \lforall[x][!A(x)]$
\DisplayProof
\\[4ex]
\Axiom$ !A(a), \Gamma \fCenter \Delta $
\RightLabel{\LeftR{\lexists}}
\UnaryInf$ \lexists[x][!A(x)], \Gamma \fCenter \Delta$
\DisplayProof
&
\Axiom$ \Gamma \fCenter !A(t) $
\RightLabel{\RightR{\lexists}}
\UnaryInf$ \Gamma \fCenter \lexists[x][!A(x)]$
\DisplayProof
\end{tabular}
\caption{உள்ளுணர்வுவாதத் தருக்கத்திற்கான \Log{G3i} விதிகள்.
$\RightR{\forall}$ மற்றும் $\LeftR{\exists}$ ஆகிய விதிகளில்,
!!{constant}~$a$ முடிவுத் தொடரணியில் இடம்பெறக் கூடாது.
$\Gamma$ என்பது !!{formula}s இன் பன்மைக் கணம்;
$\Delta$ வெறுமையாகவோ ஒரே ஓர் !!{formula} மட்டும் கொண்டதாகவோ
இருக்கலாம். அடிக்கோள்களில் !!{formula}~$!C$
அணு வாய்பாடாக இருக்க வேண்டும். \Log{G1m} என்பது
\Log{G1i} இல் \RightR{\Weakening} விதியும்
$\lfalse, \Gamma \Sequent \Delta$ என்ற அடிக்கோள்களும்
நீக்கப்பட்ட முறைமை.}
\ollabel{tab:G3i}
\end{table}


% END SOURCE UNIT OLP-0708

\clearpage

% BEGIN SOURCE UNIT OLP-0709 content/proof-theory/sequent-calculus/rules-LK.tex
% Exact modular target bytes: 3860; SHA256: 135f0f89418bb735382e65ecff8ad426a268b1d97c891f360ec37946f050c226
% SEGMENT OLP-0709-S01


\begin{table}
\begin{tabular}{@{}c@{}c@{}}
\multicolumn{2}{@{}c@{}}{அடிக்கோள்கள்}\\[2ex]
$!A \Sequent !A$ & $\lfalse \Sequent \quad$
\\[2ex]
% SEGMENT OLP-0709-S02
\multicolumn{2}{@{}c@{}}{கட்டமைப்பு விதிகள்}\\[2ex]
\Axiom$ \Gamma \fCenter \Delta $
\RightLabel{\LeftR{\Weakening}}
\UnaryInf$ !A, \Gamma \fCenter \Delta$
\DisplayProof
&
\Axiom$ \Gamma \fCenter \Delta$
\RightLabel{\RightR{\Weakening}}
\UnaryInf$ \Gamma \fCenter \Delta, !A$
\DisplayProof
\\[4ex]
\Axiom$ !A, !A, \Gamma \fCenter \Delta $
\RightLabel{\LeftR{\Contraction}}
\UnaryInf$ !A, \Gamma \fCenter \Delta$
\DisplayProof
&
\Axiom$ \Gamma \fCenter \Delta, !A, !A$
\RightLabel{\RightR{\Contraction}}
\UnaryInf$ \Gamma \fCenter \Delta, !A$
\DisplayProof
\\[4ex]
\Axiom$ \Gamma, !A, !B, \Pi \fCenter \Delta $
\RightLabel{\LeftR{\Exchange}}
\UnaryInf$\Gamma, !B, !A, \Pi \fCenter \Delta$
\DisplayProof
&
\Axiom$ \Gamma \fCenter \Delta, !A, !B, \Lambda$
\RightLabel{\RightR{\Exchange}}
\UnaryInf$ \Gamma \fCenter \Delta, !B, !A, \Lambda$
\DisplayProof
\\[4ex]
% SEGMENT OLP-0709-S03
\multicolumn{2}{@{}c@{}}{தருக்க விதிகள்}\\[2ex]
\Axiom$ \Gamma \fCenter \Delta, !A $
\RightLabel{\LeftR{\lnot}}
\UnaryInf$ \lnot !A, \Gamma \fCenter \Delta$
\DisplayProof
&
\Axiom$!A, \Gamma \fCenter \Delta$
\RightLabel{\RightR{\lnot}}
\UnaryInf$ \Gamma \fCenter \Delta, \lnot !A $
\DisplayProof
\\[4ex]\begin{tabular}[t]{@{}l@{}}
\Axiom$ !A, \Gamma \fCenter \Delta$
\RightLabel{\LeftR{\land}}
\UnaryInf$ !A \land !B, \Gamma \fCenter \Delta$
\DisplayProof
\\[3ex]
\Axiom$!B, \Gamma \fCenter \Delta$
\RightLabel{\LeftR{\land}}
\UnaryInf$!A \land !B, \Gamma \fCenter \Delta$
\DisplayProof\\[3ex]
\end{tabular}
&
\Axiom$\Gamma \fCenter \Delta, !A$
\Axiom$ \Gamma \fCenter \Delta, !B$
\RightLabel{\RightR{\land}}
\BinaryInf$ \Gamma \fCenter \Delta, !A \land !B $
\DisplayProof
\\[4ex]\Axiom$!A, \Gamma \fCenter \Delta$
\Axiom$!B, \Gamma \fCenter \Delta$
\RightLabel{\LeftR{\lor}}
\BinaryInf$!A \lor !B, \Gamma \fCenter \Delta$
\DisplayProof
&
\begin{tabular}[t]{@{}r@{}}
\Axiom$\Gamma \fCenter \Delta, !A$
\RightLabel{\RightR{\lor}}
\UnaryInf$ \Gamma \fCenter \Delta, !A \lor !B$
\DisplayProof
\\[3ex]
\Axiom$ \Gamma \fCenter \Delta, !B$
\RightLabel{\RightR{\lor}}
\UnaryInf$ \Gamma \fCenter \Delta, !A \lor !B$
\DisplayProof\\[3ex]
\end{tabular}
\\[4ex]
\Axiom$ \Gamma \fCenter \Delta, !A$
\Axiom$ !B, \Gamma \fCenter \Delta$
\RightLabel{\LeftR{\lif}}
\BinaryInf$ !A \lif !B, \Gamma \fCenter \Delta$
\DisplayProof
&
\Axiom$ !A, \Gamma \fCenter \Delta, !B$
\RightLabel{\RightR{\lif}}
\UnaryInf$ \Gamma \fCenter \Delta, !A \lif !B $
\DisplayProof
\\[4ex]
\Axiom$ !A(t), \Gamma \fCenter \Delta$
\RightLabel{\LeftR{\lforall}}
\UnaryInf$ \lforall[x][!A(x)],\Gamma \fCenter \Delta$
\DisplayProof
&
\Axiom$ \Gamma \fCenter \Delta, !A(a) $
\RightLabel{\RightR{\lforall}}
\UnaryInf$ \Gamma \fCenter \Delta, \lforall[x][!A(x)]$
\DisplayProof
\\[4ex]
\Axiom$ !A(a), \Gamma \fCenter \Delta $
\RightLabel{\LeftR{\lexists}}
\UnaryInf$ \lexists[x][!A(x)], \Gamma \fCenter \Delta$
\DisplayProof
&
\Axiom$ \Gamma \fCenter \Delta, !A(t) $
\RightLabel{\RightR{\lexists}}
\UnaryInf$ \Gamma \fCenter \Delta, \lexists[x][!A(x)]$
\DisplayProof
\end{tabular}
\caption{\Log{LK} இன் விதிகள். $\RightR{\forall}$ மற்றும்
$\LeftR{\exists}$ ஆகிய விதிகளில், !!{constant}~$a$
முடிவுத் தொடரணியில் இடம்பெறக் கூடாது. $\Gamma$, $\Delta$,
$\Pi$, $\Lambda$ ஆகியவை !!{formula}s இன் வரிசைகள்.}
\ollabel{tab:LK}
\end{table}


% END SOURCE UNIT OLP-0709

\clearpage

% BEGIN SOURCE UNIT OLP-0710 content/proof-theory/sequent-calculus/rules-mG3i.tex
% Exact modular target bytes: 3098; SHA256: 9800548758f58b99f905996abe4f623cf564beb2b7f6601972137b98997dbc5e
% SEGMENT OLP-0710-S01


\begin{table}
\begin{tabular}{@{}c@{}c@{}}
\multicolumn{2}{@{}c@{}}{அடிக்கோள்கள்}\\[2ex]
$!C, \Gamma \Sequent \Delta, !C$ & $\lfalse, \Gamma \Sequent \Delta$
\\[2ex]
% SEGMENT OLP-0710-S02
\multicolumn{2}{@{}c@{}}{தருக்க விதிகள்}\\[2ex]
\Axiom$\lnot !A, \Gamma \fCenter \Delta, !A $
\RightLabel{\LeftR{\lnot}}
\UnaryInf$\lnot !A, \Gamma \fCenter \Delta$
\DisplayProof
&
\Axiom$!A, \Gamma \fCenter$
\RightLabel{\RightR{\lnot}}
\UnaryInf$ \Gamma \fCenter \Delta, \lnot !A $
\DisplayProof
\\[4ex]
\Axiom$ !A, !B, \Gamma \fCenter \Delta$
\RightLabel{\LeftR{\land}}
\UnaryInf$ !A \land !B, \Gamma \fCenter \Delta$
\DisplayProof
&
\Axiom$\Gamma \fCenter \Delta, !A$
\Axiom$ \Gamma \fCenter \Delta, !B$
\RightLabel{\RightR{\land}}
\BinaryInf$ \Gamma \fCenter \Delta, !A \land !B $
\DisplayProof
\\[4ex]\Axiom$!A, \Gamma \fCenter \Delta$
\Axiom$!B, \Gamma \fCenter \Delta$
\RightLabel{\LeftR{\lor}}
\BinaryInf$!A \lor !B, \Gamma \fCenter \Delta$
\DisplayProof
&
\Axiom$\Gamma \fCenter \Delta, !A, !B$
\RightLabel{\RightR{\lor}}
\UnaryInf$ \Gamma \fCenter \Delta, !A \lor !B$
\DisplayProof
\\[4ex]
\Axiom$!A \lif !B, \Gamma \fCenter \Delta, !A$
\Axiom$ !B, \Gamma \fCenter \Delta$
\RightLabel{\LeftR{\lif}}
\BinaryInf$!A \lif !B, \Gamma \fCenter \Delta$
\DisplayProof
&
\Axiom$ !A, \Gamma \fCenter !B$
\RightLabel{\RightR{\lif}}
\UnaryInf$ \Gamma \fCenter \Delta, !A \lif !B $
\DisplayProof
\\[4ex]
\Axiom$ !A(t), \lforall[x][!A(x)],\Gamma \fCenter \Delta$
\RightLabel{\LeftR{\lforall}}
\UnaryInf$ \lforall[x][!A(x)],\Gamma \fCenter \Delta$
\DisplayProof
&
\Axiom$ \Gamma \fCenter !A(a) $
\RightLabel{\RightR{\lforall}}
\UnaryInf$ \Gamma \fCenter \lforall[x][!A(x)]$
\DisplayProof
\\[4ex]
\Axiom$ !A(a), \Gamma \fCenter \Delta $
\RightLabel{\LeftR{\lexists}}
\UnaryInf$ \lexists[x][!A(x)], \Gamma \fCenter \Delta$
\DisplayProof
&
\Axiom$ \Gamma \fCenter \Delta, \lexists[x][!A(x)], !A(t) $
\RightLabel{\RightR{\lexists}}
\UnaryInf$ \Gamma \fCenter \Delta, \lexists[x][!A(x)]$
\DisplayProof
\end{tabular}
\caption{உள்ளுணர்வுவாதத் தருக்கத்திற்கான பல-முடிவு
\Log{mG3i} விதிகள். $\RightR{\forall}$ மற்றும்
$\LeftR{\exists}$ ஆகிய விதிகளில், !!{constant}~$a$
முடிவுத் தொடரணியில் இடம்பெறக் கூடாது. $\Gamma$ மற்றும்
$\Delta$ ஆகியவை !!{formula}s இன் பன்மைக் கணங்கள்.
அடிக்கோள்களில் !!{formula}~$!C$ அணு வாய்பாடாக இருக்க வேண்டும்.
\Log{mG1m} என்பது $\lfalse, \Gamma \Sequent \Delta$
என்ற அடிக்கோள்கள் நீக்கப்பட்ட \Log{mG1i} ஆகும்.}
\ollabel{tab:mG3i}
\end{table}


% END SOURCE UNIT OLP-0710

\clearpage
\endgroup

\chapter{கூற்றுத் தருக்கமும் இரண்டாம்தரத் தருக்கமும்}
\clearpage
\noindent மூல அலகு: \texttt{OLP-0715}\par

% BEGIN SOURCE UNIT OLP-0715 content/propositional-logic/syntax-and-semantics/soundness.tex
% Exact modular target bytes: 1810; SHA256: d0bd58391df12477c251b31194d5b081cd7fd39c49d8288d8f930af717b559b1
% SEGMENT OLP-0715-S01


\olfileid{pl}{prp}{snd}

\olsection{கூற்றுத் தருக்கத்தின் சரித்தன்மை}

\begin{thm}[சரித்தன்மை]
\ollabel{thm:soundness}
$\Gamma\Proves!A$ எனில் $\Gamma \Entails !A$.
\end{thm}

% SEGMENT OLP-0715-S02
\begin{proof}
தேற்றங்களின் மீது தொகுத்தறிதல் செய்கிறோம். $!A$ ஓர் அடிகோள் எனில்,
$\Entails !A$; எனவே குறிப்பாக $\Gamma \Entails!A$.
அதேபோல் $!A \in \Gamma$ எனில், $\Gamma \Entails!A$.
தொகுத்தறிதல் படியில், $!C$ மற்றும் $!C\lif!A$ ஆகியவற்றிலிருந்து
\emph{மோடஸ் போனென்ஸ்} விதியால் $!A$ பெறப்படுகிறது என்க;
மேலும் $!C\lif !A$ க்கும் $!C$ க்கும் தேற்றம் பொருந்துகிறது என்பது
தொகுத்தறிதல் எடுகோள் என்க. அப்போது
\olref[pl][syn][sem]{prop:semanticalfacts} இன் இரண்டாவது உருப்படி
தேவையான முடிவைக் கொடுக்கிறது.
\end{proof}

% SEGMENT OLP-0715-S03
\begin{cor}
$\Gamma$ நிறைவுறத்தக்கது எனில், $\Gamma$ முரணற்றது.
எனவே கூற்றுத் தருக்கம் முரணற்றது.
\end{cor}


% END SOURCE UNIT OLP-0715

\clearpage
\noindent மூல அலகு: \texttt{OLP-0714}\par

% BEGIN SOURCE UNIT OLP-0714 content/propositional-logic/syntax-and-semantics/completeness.tex
% Exact modular target bytes: 8822; SHA256: fa8d05f37caf7da9861d769a04caa3e3e7ed6f35c711cbf16edfc5d043c164a0
% SEGMENT OLP-0714-S01


\olfileid{pl}{prp}{com}

\olsection{கூற்றுத் தருக்கத்தின் நிறைவுத்தன்மை}

\begin{defn}
\ollabel{def:MaxCon}
!!{formula}s கொண்ட கணம்~$\Gamma$, முரணற்றதாகவும்,
$\Gamma \subseteq \Delta$ ஆக இருக்கும் முரணற்ற
கணம்~$\Delta$ எதற்கும் $\Gamma = \Delta$ ஆகவும்
இருந்தால், அது \emph{அதிகபட்ச முரணற்றது} எனப்படும்.
\end{defn}

% SEGMENT OLP-0714-S02
\begin{lem}[மெய்மைத் துணைத்தேற்றம்]
  \ollabel{lem:truth}
$\Gamma$ அதிகபட்ச முரணற்றது என்க. அப்போது:
\begin{enumerate}
\item \ollabel{lem:truth-1} $!A \in \Gamma$ அப்போதும்
  அப்போது மட்டுமே $\lnot !A \notin \Gamma$;
\item \ollabel{lem:truth-2} $\Gamma \Proves !A$ அப்போதும்
  அப்போது மட்டுமே $!A \in \Gamma$;
\item \ollabel{lem:truth-3} $!A \lif !B \in \Gamma$
  அப்போதும் அப்போது மட்டுமே $!A \notin \Gamma$
  அல்லது $!B \in \Gamma$.
\end{enumerate}
\end{lem}

% SEGMENT OLP-0714-S03
\begin{proof}
\begin{enumerate}
\item $!A \in \Gamma$ என்க. மேலும் $\lnot !A \in \Gamma$
  எனில் $\Gamma$ முரணுடையது. $!A$ உம்
  $\lnot!A$ உம் இரண்டுமே $\Gamma$ இல் இல்லை எனில்,
  $\Gamma\cup\{!A\}$, $\Gamma \cup \{\lnot !A\}$
  என்ற இரு விரிவாக்கங்களும் முரணுடையவையாக இருக்க
  முடியாது: அப்படியாயின் வழக்கமான இருநிலை
  முரண்-பிரிப்பு வாதத்தால் ஆரம்பக் கணமே முரணுடையதாகும்.
  ஆகவே அவற்றில் ஒன்று முரணற்றது; இது $\Gamma$
  அதிகபட்சமானது என்பதற்கு முரண். மறுதிசையும்
  இதே முறையில் நிறுவப்படுகிறது.

% SEGMENT OLP-0714-S04
\item $!A \in \Gamma$ எனில் உடனே
  $\Gamma \Proves !A$. மறுபுறம், $\Gamma \Proves !A$
  என்க. $!A \notin \Gamma$ எனில்
  \olref{lem:truth-1} இன் படி $\lnot !A \in \Gamma$.
  ஆகவே $\Gamma \Proves !A$ மற்றும்
  $\Gamma \Proves \lnot !A$ இரண்டும் கிடைக்கும்;
  இது முரண்பாடு.
\item பயிற்சி.
\end{enumerate}
\end{proof}

% SEGMENT OLP-0714-S05
\olref{lem:truth}\olref{lem:truth-2} இன் இடமிருந்து
வலத்திற்கான திசை, $\Gamma$ என்பது
\emph{வருவித்தலின் கீழ் மூடப்பட்டது} எனக் காட்டுகிறது:
எந்த~$!A$ க்கும் $\Gamma \Proves !A$ எனில்
$!A \in \Gamma$.

\begin{prob}
  \olref{lem:truth} இன் நிறுவலை நிறைவு செய்க.
\end{prob}

% SEGMENT OLP-0714-S06
\begin{thm}[நிறைவுத்தன்மை]
\ollabel{thm:completeness}
$\Gamma$ முரணற்றது எனில் $\Gamma$ நிறைவுறத்தக்கது.
\end{thm}

\begin{proof}
மொழியின் எல்லா !!{formula}s ஐயும் முழுமையாகப்
பட்டியலிடும் $!A_0$, $!A_1$,~\dots என்ற வரிசையை
எடுத்துக் கொள்வோம். !!{formula}s கொண்ட
அதிகரிக்கும் கணங்களின் வரிசை $\Gamma_0$,
$\Gamma_1$,~\dots ஐப் பின்வருமாறு தொகுத்தறிதலாக
வரையறுக்கிறோம்:
\begin{align*}
\Gamma_0 & = \Gamma\\
\Gamma_{n+1} &  =
\begin{cases}
  \Gamma_n \cup \{!A_n\} & \text{$\Gamma_n \cup \{!A_n\}$ முரணற்றது எனில்;}\\
  \Gamma_n \cup \{\lnot !A_n\} & \text{இல்லாவிட்டால்}.
\end{cases}
\end{align*}
பின்னர்
\[
\Gamma^* = \bigcup_{0\le n}\Gamma_n.
\]
என வரையறுக்கிறோம்.

% SEGMENT OLP-0714-S07
இப்போது பின்வரும் உண்மைகளை வரிசையாக நிறுவ வேண்டும்:
\begin{enumerate}
\item ஒவ்வொரு $n$ க்கும் $\Gamma_n$ முரணற்றது.
  இதை $n$ மீது தொகுத்தறிதலாலும், ஒரு முரணற்ற
  கணத்துடன் ஒரு வாய்பாட்டையோ அதன் மறுப்பையோ
  சேர்க்கும் இரு விரிவாக்கங்களில் குறைந்தது ஒன்று
  முரணற்றது என்ற மேலுள்ள தேர்வு வாதத்தாலும் பெறுகிறோம்;
\item $\Gamma^*$ முரணற்றது;
\item $\Gamma^*$ அதிகபட்சமானது.
\end{enumerate}
பின்னர் $p_i \in \Gamma^*$ அப்போதும் அப்போது
மட்டுமே $v(p_i) = \True$ என்றவாறு மெய்மதிப்பு
ஒதுக்கீடு~$v$ ஐ வரையறுக்கிறோம். $!A$ மீது
தொகுத்தறிதலால், சிக்கலான !!{sentence}s க்கும்
$\Gamma^*$ இல் இருப்பதும் $v$ இன் படி மெய்யாக
இருப்பதும் ஒன்றே, அதாவது
$\overline{v}(!A) = \True$ அப்போதும் அப்போது
மட்டுமே $!A \in \Gamma^*$ என்று காட்டுகிறோம்.
குறிப்பாக, $v$ ஆனது $\Gamma^*$ ஐ நிறைவுறுத்துகிறது;
மேலும் $\Gamma \subseteq \Gamma^*$ என்பதால்
வேண்டியபடி $\Gamma$ உம் நிறைவுறத்தக்கது.
\end{proof}

% SEGMENT OLP-0714-S08
\begin{cor}
$\Gamma \Entails !A$ எனில் $\Gamma \Proves !A$.
\end{cor}

\begin{proof}
$\Gamma\Proves/ !A$ எனில்
$\Gamma \cup \{\lnot !A\}$ முரணற்றது:
இது முரணுடையதாக இருந்தால், நிறுவத்தக்கதன்மைக்கும்
முரணுக்குமான வழக்கமான தொடர்பால் $\Gamma$ விலிருந்து
அந்த வாய்பாடு நிறுவப்படும்; அது கருதுகோளுக்கு முரண்.
ஆகவே நிறைவுத்தன்மைத் தேற்றத்தின்படி
$\Gamma \cup \{\lnot !A\}$ நிறைவுறத்தக்கது.
எனவே $\Gamma \Entails/ !A$.
\end{proof}

% SEGMENT OLP-0714-S09
\begin{prop}[இறுக்கத்தன்மைத் தேற்றம்]
\ollabel{prop:compactness}
$\Gamma$ நிறைவுறத்தக்கது அப்போதும் அப்போது
மட்டுமே, $\Gamma$ வின் ஒவ்வொரு
\emph{முடிவுறு} உட்கணமும்~$\Gamma_0$
நிறைவுறத்தக்கது.
\end{prop}

\begin{proof}
$\Gamma$ நிறைவுற இயலாதது அப்போதும் அப்போது
மட்டுமே அது முரணுடையது; அப்போதும் அப்போது
மட்டுமே $\Gamma$ வின் ஏதோ ஒரு முடிவுறு
உட்கணம்~$\Gamma_0$ முரணுடையது; அப்போதும்
அப்போது மட்டுமே அந்த முடிவுறு
உட்கணம்~$\Gamma_0$ நிறைவுற இயலாதது.
\end{proof}

% SEGMENT OLP-0714-S10
\begin{cor}
$\Gamma \Entails !A$ அப்போதும் அப்போது
மட்டுமே, $\Gamma$ வின் ஏதோ ஒரு முடிவுறு
உட்கணம்~$\Gamma_0$ க்கு
$\Gamma_0 \Entails !A$.
\end{cor}


% END SOURCE UNIT OLP-0714

\clearpage
\noindent மூல அலகு: \texttt{OLP-0716}\par

% BEGIN SOURCE UNIT OLP-0716 content/second-order-logic/syntax-and-semantics/language-of-sol.tex
% Exact modular target bytes: 3870; SHA256: 8ff252e13fc325ef1db25b40aa6c5eafc8686ef2ed27a5d8d3d6c6ab81f421ff
% SEGMENT OLP-0716-S01


\olfileid{sol}{syn}{lan}

\olsection{இரண்டாம்தரத் தருக்கத்தின் மொழி}

முதல்தரத் தருக்கத்தைப் போலவே, இரண்டாம்தரத் தருக்கத்தின் கோவைகள்
\emph{!!{variable}s}, \emph{!!{constant}s}, \emph{!!{predicate}s},
சில நேரங்களில் \emph{!!{function}s} ஆகியவற்றைக் கொண்ட அடிப்படைச்
சொற்கோவையிலிருந்து கட்டியெழுப்பப்படுகின்றன. இவற்றுடன் தருக்க
இணைப்பிகள், அளவையடைகள், அடைப்புக்குறிகள் மற்றும் காற்புள்ளிகள் போன்ற
நிறுத்தற்குறிகளைச் சேர்த்து, \emph{உறுப்புகளும்}
\emph{!!{formula}s உம்} அமைக்கப்படுகின்றன. பொருள்களுக்கான மாறிகளோடு,
உறவுகளுக்கும் சார்புகளுக்குமான மாறிகளையும் இரண்டாம்தரத் தருக்கம்
கொண்டிருப்பதும், அவற்றின்மீது அளவையிடலை அனுமதிப்பதுமே வேறுபாடு.
எனவே இரண்டாம்தரத் தருக்கத்தின் தருக்கக் குறிகள் முதல்தரத்
தருக்கத்தின் குறிகளுடன் பின்வருவனவற்றையும் கொண்டவை:

% SEGMENT OLP-0716-S02
\begin{enumerate}
\item ஒவ்வோர் இட எண்ணிக்கை~$n$ க்கும் இரண்டாம்தர உறவு
  !!{variable}s இன் ஓர் !!{denumerable}s கணம்:
  $\Obj V_0^n$, $\Obj V_1^n$, $\Obj V_2^n$, \dots
\item இரண்டாம்தரச் சார்பு !!{variable}s இன் ஓர் !!{denumerable}s கணம்:
  $\Obj u_0^n$, $\Obj u_1^n$, $\Obj u_2^n$, \dots
\end{enumerate}

% SEGMENT OLP-0716-S03
முதல்தரத் தருக்கத்தில் !!{identity}~$\eq$ பொதுவாகச் சேர்க்கப்படுகிறது.
முதல்தரத் தருக்கத்தில், ஒரு மொழி~$\Lang{L}$ இன் தருக்கமல்லாத குறிகள்
சுவையான எதையும் வெளிப்படுத்த இன்றியமையாதவை. தருக்கமல்லாத குறிகள்
எதையும் பயன்படுத்தாத !!{sentence}s நிச்சயமாக உள்ளன; ஆனால்~$\eq$
மட்டும் கொண்டு சுவையான எதையும் கூறுவது கடினம். இரண்டாம்தரத்
தருக்கத்தில், உறவு மற்றும் சார்பு மாறிகள் வரம்பின்றிக் கிடைப்பதால்,
தருக்கமல்லாத குறிகளின் தனிப்பட்ட தொகுதி இல்லாமலேயே ஒரு முதல்தர
மொழியில் சொல்லக்கூடிய எதையும் நம்மால் சொல்ல முடியும்.


% END SOURCE UNIT OLP-0716


\chapter{கணங்கள், தொடர்புகள், சார்புகளின் மாற்று வடிவங்கள்}
\clearpage
\noindent மூல அலகு: \texttt{OLP-0717}\par

% BEGIN SOURCE UNIT OLP-0717 content/sets-functions-relations/functions/isomorphic-functions.tex
% Exact modular target bytes: 3441; SHA256: 24ae25a6f9e0bf71d07cf6d5b03169e69b5072b47b453b8979574ed1816cf339
% SEGMENT OLP-0717-S01


\olfileid{sfr}{fun}{iso}
\olsection{ஓரினச்சார்பு}

\begin{explain}
\emph{ஓரினச்சார்பு} என்பது, அது தொடர்புபடுத்தும் கணங்களின்
அமைப்பைப் பேணும் ஓர் இருபுறச் சார்பு; இங்கு அமைப்பு என்பது
அக்கணங்களின் !!{element}s இடையே நிலவும் உறவுகளைச் சார்ந்தது.
பின்வரும் இரண்டு கணங்களைக் கருதுக: $X=\{1,2,3\}$ மற்றும்
$Y=\{4,5,6\}$. இரு கணங்களுக்கும் தொடரி, சிறியது,
பெரியது ஆகிய உறவுகளால் அமைப்பு வழங்கப்படுகிறது.
இரு கணங்களுக்கிடையிலான ஓரினச்சார்பு அவ்வமைப்புகளைப் பேணும்
!!a{bijection} ஆகும். எனவே !!a{bijective} சார்பு $f \colon X
\to Y$ ஆனது $i<j$ அப்போதும் அப்போது மட்டுமே $f(i)<f(j)$,
$i>j$ அப்போதும் அப்போது மட்டுமே $f(i)>f(j)$, மேலும் $j$ ஆனது
$i$ இன் தொடரி அப்போதும் அப்போது மட்டுமே $f(j)$ ஆனது
$f(i)$ இன் தொடரி எனில், அது ஓரினச்சார்பு ஆகும்.
\end{explain}

% SEGMENT OLP-0717-S02
\begin{defn}[ஓரினச்சார்பு]
$U$ என்பது $\langle X, R\rangle$ என்ற ஜோடியாகவும், $V$ என்பது
$\langle Y, S\rangle$ என்ற ஜோடியாகவும் இருக்கட்டும்; இங்கு
$X$, $Y$ கணங்கள், $R$, $S$ முறையே $X$, $Y$ மீதான உறவுகள்.
$X$ இலிருந்து $Y$ க்குச் செல்லும் !!a{bijection}~$f$ உறவு
அமைப்பைப் பேணினால், அப்போதும் அப்போது மட்டுமே, அது $U$ இலிருந்து
$V$ க்குச் செல்லும் \emph{ஓரினச்சார்பு} ஆகும். அதாவது,
$X$ இல் உள்ள எந்த $x_{1}$, $x_{2}$ க்கும்
$\tuple{x_1,x_2} \in R$ அப்போதும் அப்போது மட்டுமே
$\tuple{f(x_1),f(x_2)} \in S$.
\end{defn}

% SEGMENT OLP-0717-S03
\begin{ex}
பின்வரும் இரண்டு கணங்கள் $X=\{1,2,3\}$,
$Y=\{4,5,6\}$ மற்றும் சிறியது, பெரியது என்ற உறவுகளைக்
கருதுக. $f(x) = 7-x$ என வரையறுக்கப்படும்
$f\colon X \to Y$ என்ற சார்பு, $\tuple{X,<}$ க்கும்
$\tuple{Y,>}$ க்கும் இடையிலான ஓரினச்சார்பு ஆகும்.
\end{ex}


% END SOURCE UNIT OLP-0717

\clearpage
\noindent மூல அலகு: \texttt{OLP-0718}\par

% BEGIN SOURCE UNIT OLP-0718 content/sets-functions-relations/inductive-defs-proofs/introduction.tex
% Exact modular target bytes: 11259; SHA256: 4a13308fd2b322051cebdb0d26346b99a2f6d4e0fd4f211a0b4eaea940179ccd
% SEGMENT OLP-0718-S01


\olfileid{sfr}{ind}{int}

\olsection{அறிமுகம்}

\begin{explain}
தருக்கத்தில் பரவலாகப் பயன்படுத்தப்படும் கணித நிறுவல் உத்தி
தொகுத்தறிதல். கணிதத்தில் இயல் எண்கள் மீது தொகுத்தறிதல் செய்வதை
நீங்கள் அறிந்திருக்கலாம். ஆனால் கணிதத்திலும் தருக்கத்திலும்
பொருள்களைக் கொண்ட பல்வேறு வகைக் கணங்கள் தொகுத்தறிதலுக்கு ஏற்றவை.

பொதுவாகத் தொகுத்தறிதல், அனைத்தையும் பற்றிய ஒரு கூற்றை நிறுவ
உதவுகிறது: ஒரு குறிப்பிட்ட வகையைச் சேர்ந்த ஒவ்வொரு பொருளுக்கும்
ஒரு பண்பு உள்ளது என்று காட்டுகிறது. குறிப்பாக,
``அனைத்தளவை அறிமுகம்'' வகை நிறுவல் முறை மிகவும் சிக்கலாக இருக்கும்
இடங்களில் தொகுத்தறிதல் பயனுள்ளதாக இருக்கும். சிறப்பு முறைகளில்
அமைக்கப்பட்ட கணிதப் பொருள்களுக்கே தொகுத்தறிதல் பொருந்தும்:
ஒரு கணத்தின் ஒவ்வோர் உறுப்பும் அடிப்படை உறுப்பாகவோ,
அடிப்படை உறுப்புகளிலிருந்து உருவாக்கப்பட்டதாகவோ இருந்தால்,
அதன்மீது தொகுத்தறிதல் செய்யலாம். இன்னும் துல்லியமாகச் சொல்வதானால்:
\end{explain}

% SEGMENT OLP-0718-S02
\begin{defn}
ஒரு கணம் $S$, ஒரு சார்பு $f$ இன் கீழ் \emph{மூடப்பட்டது}
என்பது, அப்போதும் அப்போது மட்டுமே, $x \in S$ எனும்போதெல்லாம்
$f(x) \in S$ ஆகும்.
\end{defn}

\begin{explain}
எடுத்துக்காட்டாக, இயல் எண்களின் கணம் ($\mathbb{N}$),
தொடரிச் சார்பு $f(x) = x+1$ இன் கீழ் மூடப்பட்டது.
$x$ ஓர் இயல் எண் எனில், $x+1$ உம் இயல் எண்.
ஆனால் முன்னிச் சார்பு $f(x) = x-1$ இன் கீழ் இயல் எண்களின்
கணம் மூடப்படவில்லை; ஏனெனில் $0$ ஓர் இயல் எண்,
ஆனால் $-1$ இயல் எண் அல்ல.
\end{explain}

% SEGMENT OLP-0718-S03
\begin{defn}
ஒரு பண்பு $P$, ஒரு சார்பு $f(x)$ இன் கீழ் \emph{பேணப்படுகிறது}
எனப்படுவது, $f$ இன் சார்பகத்தில் உள்ள எந்தப் பொருள் $a$ க்கும்,
$P(a)$ எனில் $P(f(a))$ ஆகும்போது ஆகும்
($a$ க்குப் பண்பு $P$ இருந்தால், $f(a)$ க்கும் அப்பண்பு உண்டு).
\end{defn}

\begin{explain}
``இரட்டை எண்ணாக இருத்தல்'' என்ற பண்பு, $f(x) = x+2$
என்ற சார்பின் கீழ் பேணப்படுகிறது; ஏனெனில் $x$ இரட்டை எண்
எனில், $x+2$ உம் இரட்டை எண். ஆனால் இப்பண்பு தொடரிச்
சார்பின் கீழ் பேணப்படுவதில்லை; $x$ இரட்டை எண் எனில்,
$x+1$ ஒற்றை எண்.

ஒவ்வோர் இயல் எண்ணையும் 0 இலிருந்து தொடரிச் சார்பை முடிவுறு
எண்ணிக்கையிலான முறை பயன்படுத்தி அடையலாம் என்பதால்தான்
இயல் எண்கள் மீது தொகுத்தறிதல் செய்ய முடிகிறது.
தொகுத்தறிதல் செய்யக்கூடிய எந்தக் கணத்திற்கும் இத்தகைய
உருவாக்கப் பண்பு உண்டு.
\end{explain}

% SEGMENT OLP-0718-S04
\begin{defn}[$\Nat$ மீது தொகுத்தறிதல்]
0 க்கு ஒரு பண்பு $P$ இருந்து, $P$ தொடரிச் சார்பின் கீழ்
பேணப்பட்டால், ஒவ்வோர் இயல் எண்ணுக்கும் பண்பு $P$ உண்டு.
\end{defn}

% SEGMENT OLP-0718-S05
\begin{ex}
பூச்சியத்திலிருந்து $n$ வரையிலான இயல் எண்களின் கூட்டுத்தொகை
$\frac{n(n+1)}{2}$. \\

\textbf{அடிப்படை நிலை.} பூச்சியத்திலிருந்து பூச்சியம் வரையிலான கூட்டுத்தொகை
$0=\frac{0\cdot 1}{2}$.\\

\textbf{தொகுத்தறிதலின் அடுத்த படி.} $k$ க்குப் பண்பு
பொருந்துகிறது என்க. அது $k+1$ க்கும் பொருந்துகிறது என்று
காட்டுவோம். அதாவது, $P(k)$ பொருந்துகிறது என்க:
\[ 0 + 1 + \ldots + k = \frac{k(k+1)}{2} \]
$P(k+1)$ பொருந்துகிறது என்று, அதாவது
\[ 0 + 1 + \ldots + k + (k+1) = \frac{(k+1)(k+2)}{2} \]
என்று, $P(k)$ என்ற எடுகோளைப் பயன்படுத்திக் காட்டுவோம்.

% SEGMENT OLP-0718-S06
\begin{align*}
0 + 1 + \ldots + k &= \frac{k(k+1)}{2} \tag{எடுகோள்} \\
0 + 1 + \ldots + k + (k+1) &= \frac{k(k+1)}{2} + (k+1)
\tag{இரு பக்கங்களிலும் $k+1$ ஐச் சேர்க்க} \\
&= \frac{k(k+1) + 2(k+1)}{2} \\
&= \frac{k^2 + k + 2k +2}{2} \\
&=\frac{(k+1)(k+2)}{2}
\end{align*}
எனவே தொகுத்தறிதலின் அடுத்த படி நிறுவப்பட்டது;
இதனால் கூற்றை நிறுவியுள்ளோம்.
\end{ex}

% SEGMENT OLP-0718-S07
\begin{explain}
!!{formula}s க்கான அடிப்படை உறுப்புகள் அணு !!{formula}s.
எளிய !!{formula}s ஐச் செயலிகளால் இணைத்து மேலும் சிக்கலான
!!{formula}s பெறப்படுகின்றன. ஒவ்வொரு செயலிக்கும் அதற்குரிய
ஒன்று அல்லது இரண்டு !!{formula}s ஐ உள்ளீடாகப் பெற்று,
அச்செயலியைப் பயன்படுத்திப் புதிய !!{formula} ஒன்றை உருவாக்கும்
சார்பு இருப்பதாகக் கருதலாம். எடுத்துக்காட்டாக,
$!A$, $!B$ என்ற இரண்டு !!{formula}s ஐ இணையல் வாய்பாடாக
இணைக்கும் சார்பை $\mathcal E _\land (!A,
!B) = !A \land !B$ என்று எழுதலாம். இவற்றை
\emph{வாய்பாடுகளை உருவாக்கும் சார்புகள்} எனலாம்.
!!{formula}s இன் கணம் இச்சார்புகளின் கீழ் மூடப்பட்டது:
$!A$, $!B$ ஆகியவை !!{formula}s எனும்போதெல்லாம்,
$\lnot !A$, $!A \land !B$ முதலியவையும் இக்கணத்தில் அடங்கும்.
\end{explain}

% SEGMENT OLP-0718-S08
\begin{defn}[!!{formula}s மீது தொகுத்தறிதல்]
ஒவ்வோர் அணு !!{formula} க்கும் பண்பு $P$ இருந்து,
!!{formula} உருவாக்கும் சார்புகளின் கீழ் $P$ பேணப்பட்டால்,
ஒவ்வொரு !!{formula} க்கும் பண்பு $P$ உண்டு.
\end{defn}

% SEGMENT OLP-0718-S09
\begin{explain}
இந்த வரையறை, !!{formula}s மீதான தொகுத்தறிதல் நிறுவல்களுக்கு
ஒரு செய்முறையைத் தருகிறது. ஒவ்வொரு !!{formula} க்கும்
பண்பு $P$ உண்டு என்று காட்ட வேண்டுமெனில், பின்வரும்
படிகளைப் பின்பற்றுக:\\

\textbf{அடிப்படை நிலை.} $!A$ ஓர் அணு !!{formula} என்க.
[\ldots] எனவே $!A$ க்குப் பண்பு $P$ உண்டு.

% SEGMENT OLP-0718-S10
\textbf{தொகுத்தறிதலின் அடுத்த படி.} $!A$, $!B$ இரண்டும்
பண்பு $P$ உடைய !!{formula}s என்க.

$\lnot$ நிகழ்வு: [\ldots] எனவே $\lnot !A$ க்குப் பண்பு $P$ உண்டு.

$\land$ நிகழ்வு: [\ldots] எனவே $!A \land !B$ க்குப் பண்பு $P$ உண்டு.

$\lor$ நிகழ்வு: [\ldots] எனவே $!A \lor !B$ க்குப் பண்பு $P$ உண்டு.

$\lif$ நிகழ்வு: [\ldots] எனவே $!A \lif !B$ க்குப் பண்பு $P$ உண்டு.

$\liff$ நிகழ்வு: [\ldots] எனவே $!A \liff !B$ க்குப் பண்பு $P$ உண்டு.

$\lforall[]$ நிகழ்வு: [\ldots] எனவே $\lforall[x] !A$ க்குப் பண்பு $P$ உண்டு.

$\lexists[]$ நிகழ்வு: [\ldots] எனவே $\lexists[x] !A$ க்குப் பண்பு $P$ உண்டு. \\

எனவே ஒவ்வொரு !!{formula} க்கும் பண்பு $P$ உண்டு.
\end{explain}


% END SOURCE UNIT OLP-0718

\clearpage
\noindent மூல அலகு: \texttt{OLP-0721}\par

% BEGIN SOURCE UNIT OLP-0721 content/sets-functions-relations/sets/proofs-about-sets.tex
% Exact modular target bytes: 9988; SHA256: 944a2320276e4fc7cadfe8660e196f999a4a763588fa4a6b0af551b17dd6df9c
% SEGMENT OLP-0721-S01


\olfileid{sfr}{set}{prf}
\olsection{கணங்கள் பற்றிய நிறுவல்கள்}

\begin{editorial}
இப்பிரிவுக்குப் பதிலாக \verb|content/methods/proofs| பயன்படுத்தப்படுகிறது;
எனவே இயல்பான அமைப்பில் இப்பிரிவு இந்த அத்தியாயத்திலிருந்து
நீக்கப்பட்டுள்ளது.
\end{editorial}

% SEGMENT OLP-0721-S02
\begin{explain}
கணங்களும் இதுவரை அறிமுகப்படுத்திய குறியீடுகளும்,
கணங்களைக் குறிப்பிடவும் அவற்றுக்கிடையிலான தொடர்புகளை
வெளிப்படுத்தவும் வசதியான சுருக்கக் குறியீடுகளைத் தருகின்றன.
அத்தகைய தொடர்புகள் பற்றிய கூற்றுகளை நிறுவுவதும் அடிக்கடி
தேவைப்படும். கணித நிறுவல்களை நீங்கள் அறிந்திருக்காவிட்டால்,
இது புதிதாக இருக்கலாம். ஆகவே ஓர் எளிய எடுத்துக்காட்டைப்
படிப்படியாகப் பார்ப்போம். எந்தக் கணங்கள் $X$, $Y$ க்கும்
$X \cap (X \cup Y) = X$ எப்போதும் பொருந்தும் என்று
நிறுவுவோம். இத்தகைய கணச் சமத்துவத்தை எவ்வாறு நிறுவுவது?
கணங்கள் அவற்றின் !!{element}s ஆல் மட்டுமே தீர்மானிக்கப்படுகின்றன
என்பதை நினைவுகூர்க. அதாவது, கணங்கள் ஒரே !!{element}s ஐக்
கொண்டிருந்தால், அப்போதும் அப்போது மட்டுமே, அவை சமம்.
எனவே இங்கு (a) $X \cap (X \cup Y)$ இன் ஒவ்வோர் !!{element}
உம் $X$ இன் !!{element} என்றும், மறுதிசையில்
(b) $X$ இன் ஒவ்வோர் !!{element} உம் $X \cap (X \cup Y)$ இன்
!!{element} என்றும் நிறுவ வேண்டும். வேறு சொற்களில்,
(a) $X \cap (X \cup Y) \subseteq X$ என்பதையும்
(b) $X \subseteq X \cap (X \cup Y)$ என்பதையும் காட்டுகிறோம்.

% SEGMENT OLP-0721-S03
``$X \cap (X \cup Y)$ இன் ஒவ்வோர் !!{element}~$z$ உம்
$X$ இன் !!{element}'' போன்ற பொதுக் கூற்றை நிறுவும்போது,
முதலில் தன்னிச்சையான $z \in X \cap (X \cup Y)$ கொடுக்கப்பட்டுள்ளது
என்று எடுத்துக்கொண்டு, அந்த எடுகோளிலிருந்து $z \in X$ என்பதை
நிறுவுகிறோம். இவ்வடிவத்தை ``பொது நிபந்தனை நிறுவல்'' என்று
நீங்கள் அறிந்திருக்கலாம். இந்நிறுவலில், எடுகோளிலும் முடிவிலும்
பயன்படும் வரையறைகளையும் பயன்படுத்த வேண்டும்; இங்கு
``$\cap$'', ``$\cup$'' ஆகியவற்றின் வரையறைகள் அவை.
எனவே (a) நிகழ்வை பின்வருமாறு நிறுவலாம்:

% SEGMENT OLP-0721-S04
\begin{quote}
(a) முதலில் $X \cap (X \cup Y) \subseteq X$ என்பதைக்
காட்ட விரும்புகிறோம். அதாவது, $\subseteq$ இன் வரையறையின்படி,
எந்த~$z$ க்கும் $z \in X \cap (X \cup Y)$ எனில்
$z \in X$ என்று காட்ட வேண்டும். ஆகவே
$z \in X \cap (X \cup Y)$ என்க. $z$ என்ற !!{element},
இரண்டு கணங்களின் வெட்டுக் கணத்தில் இருப்பது, இரு கணங்களிலும்
!!{element} ஆக இருந்தால், அப்போதும் அப்போது மட்டுமே ஆகும்.
எனவே $z \in X$ என்றும் $z \in X \cup Y$ என்றும் முடிவு
செய்யலாம். குறிப்பாக, $z \in X$; இதையே காட்ட விரும்பினோம்.
\end{quote}

% SEGMENT OLP-0721-S05
இதனால் நிறுவலின் முதல் பாதி முடிந்தது. கடைசிப் படியில்,
ஓர் எடுகோளிலிருந்து ஓர் இணையல் ($z \in X$ மற்றும்
$z \in X \cup Y$) பின்தொடர்ந்தால், அதே எடுகோளிலிருந்து
அதன் ஒவ்வோர் இணைக் கூற்றும் பின்தொடரும் என்ற உண்மையைப்
பயன்படுத்தினோம். இவ்விதியை ``இணையல் நீக்கல்'' அல்லது
$\land$ நீக்கல் என்று நீங்கள் அறிந்திருக்கலாம்.
இப்போது (b) ஐ நிறுவுவோம்:

% SEGMENT OLP-0721-S06
\begin{quote}
(b) இப்போது $X \subseteq X \cap (X \cup Y)$ என்பதை
நிறுவுகிறோம். அதாவது, $\subseteq$ இன் வரையறையின்படி,
எந்த~$z$ க்கும் $z \in X$ எனில்
$z \in X \cap (X \cup Y)$ உம் ஆகும்.
$z \in X$ என்க. $z \in X \cap (X \cup Y)$ என்பதை
நிறுவ, ``$\cap$'' இன் வரையறையின்படி, (i) $z \in X$
என்பதையும் (ii) $z \in X \cup Y$ என்பதையும் காட்ட வேண்டும்.
(i) நமது எடுகோளே; ஆகவே மேலும் நிறுவ வேண்டியது ஏதுமில்லை.
(ii) க்கு, $z$, கணங்களின் சேர்ப்புக் கணத்தின் !!{element}
ஆக இருப்பது, அக்கணங்களில் குறைந்தது ஒன்றின் !!{element}
ஆக இருந்தால், அப்போதும் அப்போது மட்டுமே என்பதை நினைவுகூர்க.
$z \in X$ என்பதாலும், $X \cup Y$ என்பது $X$, $Y$ இன்
சேர்ப்புக் கணம் என்பதாலும், இங்கு அவ்வாறே உள்ளது.
ஆகவே $z \in X \cup Y$. (i) $z \in X$ என்பதையும்
(ii) $z \in X \cup Y$ என்பதையும் காட்டிவிட்டோம்;
எனவே ``$\cap$'' இன் வரையறையின்படி
$z \in X \cap (X \cup Y)$.
\end{quote}

% SEGMENT OLP-0721-S07
இது சற்றே விரிவான விளக்கம்; ஆனால் கணங்களையும்
அவற்றுக்கிடையிலான தொடர்புகளையும் பற்றி எவ்வாறு காரணம்
கூறுகிறோம் என்பதைக் காட்டுகிறது. வழக்கமாக இவ்வளவு
வெளிப்படையாக ஒவ்வொரு படியையும் கூறுவதில்லை; குறிப்பாக,
எல்லா வரையறைகளையும் மீண்டும் கூறாமல் இருக்கலாம்.
மேம்பட்ட பாடநூலில் இம்முடிவுக்கான நிறுவல் மிகவும்
சுருக்கமாக இருக்கும். அது பின்வருமாறு அமையலாம்.
\end{explain}

% SEGMENT OLP-0721-S08
\begin{prop}[உட்கவர் பண்பு]
எல்லாக் கணங்கள் $X$, $Y$ க்கும்,
\[
X \cap (X \cup Y) = X
\]
\end{prop}

\begin{proof}
(a) $z \in X \cap (X \cup Y)$ என்க. அப்போது $z \in X$;
ஆகவே $X \cap (X \cup Y) \subseteq X$.

(b) இப்போது $z \in X$ என்க. அப்போது $z \in X \cup Y$
உம், எனவே $z \in X \cap (X \cup Y)$ உம் ஆகும்.
ஆகவே $X \subseteq X \cap (X \cup Y)$.
\end{proof}

% SEGMENT OLP-0721-S09
\begin{prob}
$X \cup (X \cap Y) = X$ என்பதை விரிவாக நிறுவுக.
பின்னர் அதன் சுருக்கமான நிறுவலைத் தருக.
(குறிப்பு: $X \cup (X \cap Y) \subseteq X$ என்ற திசைக்கு,
நிகழ்வுகளைப் பிரித்து நிறுவுதல், அதாவது $\lor$ நீக்கல்,
தேவைப்படும்.)
\end{prob}


% END SOURCE UNIT OLP-0721

\clearpage
\noindent மூல அலகு: \texttt{OLP-0719}\par
\begingroup
\TADriverHeadingMode
\olchapterid{sfr}{rel}

% BEGIN SOURCE UNIT OLP-0719 content/sets-functions-relations/relations/relations.tex
% Exact modular target bytes: 421; SHA256: cf0096bb6fd694a4bb0f7106ded9a8f9faa2971f92dddedced7636e12f9deeab
% SEGMENT OLP-0719-S01


\olchapter{sfr}{rel}{தொடர்புகள்}

\TADriverImport{relations-as-sets}

\TADriverImport{special-properties}

\TADriverImport{equivalence-relations}

\TADriverImport{orders}

\TADriverImport{graphs}

\TADriverImport{trees}

\TADriverImport{operations}

\OLEndChapterHook


% END SOURCE UNIT OLP-0719

\endgroup

\clearpage
\noindent மூல அலகு: \texttt{OLP-0720}\par
\begingroup
\TADriverHeadingMode
\olchapterid{sfr}{}

% BEGIN SOURCE UNIT OLP-0720 content/sets-functions-relations/sets-functions-relations.tex
% Exact modular target bytes: 2177; SHA256: 0b6ccd73855b62bf8aecade1078d78b84df3b598c463b93863947d6778cf9f4c
% SEGMENT OLP-0720-S01


\olpart{sfr}{கணங்கள், தொடர்புகள், சார்புகள்}

\begin{editorial}
இந்தக் கோப்பு, முறைசாராத கணக்கோட்பாடு பற்றிய பகுதியின்
முதல் நான்கு அத்தியாயங்களை, படிப்படியாக விரித்துரைக்கும்
மூல வடிவத்தில் சேர்க்கிறது. முழுப் பகுதியும்
\verb|sets-functions-relations-complete.tex| மூலம் சேர்க்கப்படுகிறது.
இப்பகுதியின் பாடப்பொருள், அடிப்படை முறைசாராத கணக்கோட்பாட்டிற்கு
ஓரளவு முழுமையான அறிமுகம். மாணவர்களுக்கு இப்பின்னணி
இருக்கும் என்று கருத முடியாவிட்டால், இப்பாடப்பொருளை,
குறைந்தது கணங்கள், சார்புகள், தொடர்புகள் பற்றியவற்றை,
மீள்பார்ப்பதுடன் பாடத்திட்டத்தைத் தொடங்குவது நல்லது.
இதனால் மற்ற தருக்கப் பகுதிகளில் எதற்கும் தேவையான
கணிதக் குறியீடுகளைப் பயன்படுத்தும் அடிப்படைத் திறன்
எல்லா மாணவர்களுக்கும் இருப்பதை உறுதிப்படுத்தலாம்.
\end{editorial}

% SEGMENT OLP-0720-S02
\TADriverImport[sets]{sets}

\TADriverImport[relations]{relations}

\TADriverImport[functions]{functions}

\TADriverImport[size-of-sets]{size-of-sets}

\OLEndPartHook

% END SOURCE UNIT OLP-0720

\endgroup

\clearpage
\noindent மூல அலகு: \texttt{OLP-0722}\par
\begingroup
\TADriverHeadingMode
\olchapterid{sfr}{siz}

% BEGIN SOURCE UNIT OLP-0722 content/sets-functions-relations/size-of-sets/size-of-sets.tex
% Exact modular target bytes: 1574; SHA256: 26c74a88d9c075fc7c1cd00a436331be60056857bb4af8557c66dd991b32921e
% SEGMENT OLP-0722-S01


\olchapter{sfr}{siz}{கணங்களின் அளவு}

\begin{editorial}
இந்த அத்தியாயம் பட்டியலாக்கங்கள், எண்ணிடத்தக்க தன்மை,
எண்ணிடத்தகாத தன்மை ஆகியவற்றை விவாதிக்கிறது.
சில பிரிவுகளுக்கு இரண்டு வடிவங்கள் உள்ளன:
ஒரு தொடக்கநிலை வடிவமும், கணக்கோட்பாட்டை இன்னும்
விரிவாக விவாதிக்கும்போது பயன்படுத்துவதற்கான ஒரு மேம்பட்ட
வடிவமும். இக்கோப்பு தொடக்கநிலை வடிவங்களை மட்டுமே சேர்க்கிறது.
மேம்பட்ட வடிவங்கள் \verb|size-of-sets-complete| மூலம் சேர்க்கப்படுகின்றன.
\end{editorial}

% SEGMENT OLP-0722-S02
\TADriverImport{introduction}

\TADriverImport{enumerability}

\TADriverImport{zig-zag}

\TADriverImport{pairing}

\TADriverImport{pairing-alt}

\TADriverImport{non-enumerability}

\TADriverImport{reduction}

\TADriverImport{equinumerous-sets}

\TADriverImport{comparing-size}

\TADriverImport{schroder-bernstein}

\OLEndChapterHook


% END SOURCE UNIT OLP-0722

\endgroup


\nocite{Frege1953}
\bibliographystyle{../upstream/bib/natbib-oup}
\bibliography{../upstream/bib/open-logic}
\end{document}

